\documentclass{article}%
\usepackage[T1]{fontenc}%
\usepackage[utf8]{inputenc}%
\usepackage{lmodern}%
\usepackage{textcomp}%
\usepackage{lastpage}%
\usepackage{geometry}%
\geometry{margin=1in,headheight=14pt,headsep=25pt}%
\usepackage{hyperref}%
\usepackage{enumitem}%
\usepackage{fancyhdr}%
\usepackage{xcolor}%
%

            \hypersetup{
                colorlinks=true,
                linkcolor=blue,
                filecolor=magenta,
                urlcolor=blue,
            }
            
            \pagestyle{fancy}
            \fancyhf{}
            \rhead{Generated on \today}
            \lhead{Course Summary}
            \cfoot{\thepage}
            
            % Configure itemize settings globally
            \setlist[itemize]{nosep}  % Reduce vertical spacing between items
            \setlist[itemize]{leftmargin=*}  % Align with left margin
        %
\title{Linear Programming Course Summary}%
\author{Generated by Scribe.AI}%
\date{\today}%
%
\begin{document}%
\normalsize%
\maketitle%
\section{Key Terms}%
\label{sec:KeyTerms}%
\begin{itemize}[leftmargin=*]%
\item \href{https://www.math.purdue.edu/~yipn/421/2024-08-27-ExSimplex.pdf#page=1}{\textbf{Simplex Method}}: An iterative algorithm used to solve linear programming problems by moving from one vertex of the feasible region to another until an optimal solution is found.%
\item \href{https://www.math.purdue.edu/~yipn/421/2024-08-27-ExSimplex.pdf#page=1}{\textbf{Slack Variable}}: A non-negative variable added to a less-than-or-equal-to constraint to transform it into an equality constraint.%
\item \href{https://www.math.purdue.edu/~yipn/421/2024-08-27-ExSimplex.pdf#page=2}{\textbf{Feasible Region}}: The set of all points that satisfy all the constraints of a linear programming problem.%
\item \href{https://www.math.purdue.edu/~yipn/421/2024-08-27-ExSimplex.pdf#page=9}{\textbf{Basic Variable}}: In the Simplex method, a variable that is expressed in terms of non-basic variables in a dictionary.%
\item \href{https://www.math.purdue.edu/~yipn/421/2024-08-27-ExSimplex.pdf#page=6}{\textbf{Non{-}basic Variable}}: In the Simplex method, a variable set to zero in a given iteration of the algorithm.%
\item \href{https://www.math.purdue.edu/~yipn/421/2024-08-27-ExSimplex.pdf#page=9}{\textbf{Pivot Operation}}: The process of exchanging a basic and a non-basic variable in the Simplex method to move to a new vertex in the feasible region.%
\item \href{https://www.math.purdue.edu/~yipn/421/2024-08-27-ExSimplex.pdf#page=26}{\textbf{Dictionary}}: A representation of a linear programming problem where basic variables are expressed as functions of non-basic variables.%
\item \href{https://www.math.purdue.edu/~yipn/421/2024-08-27-ExSimplex.pdf#page=32}{\textbf{Auxiliary Problem}}: A modified linear programming problem introduced to find an initial feasible solution when the origin is not feasible in the original problem.%
\item \href{https://www.math.purdue.edu/~yipn/421/2024-09-05-SimplexEfficiency.pdf#page=3}{\textbf{Largest{-}Coefficient Rule}}: A pivot rule in the simplex method that selects the entering variable with the largest coefficient in the objective function row. %
\item \href{https://www.math.purdue.edu/~yipn/421/2024-09-05-SimplexEfficiency.pdf#page=6}{\textbf{Largest{-}Increase Rule}}: A pivot rule in the simplex method that selects the entering variable that produces the largest increase in the objective function value. %
\item \href{https://www.math.purdue.edu/~yipn/421/2024-09-05-SimplexEfficiency.pdf#page=2}{\textbf{Klee{-}Minty Example}}: A worst-case example for the simplex method, demonstrating exponential time complexity for certain pivot rules. %
\item \href{https://www.math.purdue.edu/~yipn/421/2024-09-05-SimplexEfficiency.pdf#page=10}{\textbf{Bland's Rule}}: A pivot rule in the simplex method designed to avoid cycling, ensuring termination. %
\item \href{https://www.math.purdue.edu/~yipn/421/2024-09-05-SimplexEfficiency.pdf#page=10}{\textbf{Lexicographic Method}}: A method for choosing the leaving variable in the simplex method to prevent cycling. %
\item \href{https://www.math.purdue.edu/~yipn/421/2024-09-05-SimplexEfficiency.pdf#page=11}{\textbf{Smoothed Complexity}}: A measure of the complexity of an algorithm that considers the effect of small random perturbations to the input. %
\item \href{https://www.math.purdue.edu/~yipn/421/2024-09-05-SimplexEfficiency.pdf#page=10}{\textbf{Randomized Pivot Rule}}: A pivot rule in the simplex method that randomly permutes the order of variables before applying another rule. %
\item \href{https://www.math.purdue.edu/~yipn/421/2024-09-05-SimplexEfficiency.pdf#page=12}{\textbf{Interior Method}}: A class of algorithms for linear programming that traverse the interior of the feasible region, unlike the simplex method which follows the boundary. %
\item \href{https://www.math.purdue.edu/~yipn/421/2024-09-12-Duality.pdf#page=1}{\textbf{Primal Problem}}: A linear programming problem in its standard maximization form, aiming to maximize a linear objective function subject to linear inequality constraints and non-negativity constraints.%
\item \href{https://www.math.purdue.edu/~yipn/421/2024-09-12-Duality.pdf#page=1}{\textbf{Dual Problem}}: A linear programming problem derived from the primal problem, typically in minimization form, with constraints and objective function related to the primal problem through transposition and sign changes.%
\item \href{https://www.math.purdue.edu/~yipn/421/2024-09-12-Duality.pdf#page=2}{\textbf{Weak Duality}}: A theorem stating that the objective function value of any feasible solution to the primal problem is always less than or equal to the objective function value of any feasible solution to the dual problem.%
\item \href{https://www.math.purdue.edu/~yipn/421/2024-09-12-Duality.pdf#page=10}{\textbf{Strong Duality}}: A theorem stating that if either the primal or dual problem has an optimal solution, then so does the other, and their optimal objective function values are equal.%
\item \href{https://www.math.purdue.edu/~yipn/421/2024-09-12-Duality.pdf#page=5}{\textbf{Complementary Slackness}}: A condition stating that for optimal primal and dual solutions, the product of each primal variable and its corresponding dual slack variable is zero, and the product of each dual variable and its corresponding primal slack variable is zero.%
\item \href{https://www.math.purdue.edu/~yipn/421/2024-09-12-Duality.pdf#page=6}{\textbf{Negative Transpose Property}}: The property that the coefficient matrix of the dual problem is the negative transpose of the coefficient matrix of the primal problem, preserved throughout simplex iterations.%
\item \href{https://www.math.purdue.edu/~yipn/421/2024-09-17-DualityGeneral.pdf#page=1}{\textbf{Primal Problem}}: A linear programming problem in its standard maximization form, aiming to maximize an objective function subject to a set of constraints.%
\item \href{https://www.math.purdue.edu/~yipn/421/2024-09-17-DualityGeneral.pdf#page=1}{\textbf{Dual Problem}}: A linear programming problem derived from the primal problem, aiming to minimize an objective function subject to a set of constraints that are related to the primal constraints.%
\item \href{https://www.math.purdue.edu/~yipn/421/2024-09-17-DualityGeneral.pdf#page=2}{\textbf{Lagrange Multiplier}}: A variable introduced in the Lagrangian function to incorporate constraints into the optimization problem.%
\item \href{https://www.math.purdue.edu/~yipn/421/2024-09-17-DualityGeneral.pdf#page=2}{\textbf{Lagrangian Function}}: A function that combines the objective function and constraints of a linear programming problem using Lagrange multipliers, allowing for unconstrained optimization.%
\item \href{https://www.math.purdue.edu/~yipn/421/2024-09-17-DualityGeneral.pdf#page=5}{\textbf{Weak Duality}}: A property stating that the optimal value of the primal problem is always less than or equal to the optimal value of the dual problem.%
\item \href{https://www.math.purdue.edu/~yipn/421/2024-09-17-DualityGeneral.pdf#page=17}{\textbf{Strong Duality}}: A property stating that if the primal problem has a finite optimal solution, then the dual problem also has a finite optimal solution, and their optimal values are equal.%
\item \href{https://www.math.purdue.edu/~yipn/421/2024-09-17-DualityGeneral.pdf#page=22}{\textbf{Complementary Slackness}}: A condition that holds at optimality for primal and dual problems, stating that the product of each primal variable and its corresponding dual slack variable is zero.%
\item \href{https://www.math.purdue.edu/~yipn/421/2024-09-17-DualityGeneral.pdf#page=1}{\textbf{Objective Function}}: The function to be maximized or minimized in a linear programming problem.%
\item \href{https://www.math.purdue.edu/~yipn/421/2024-09-17-DualityGeneral.pdf#page=2}{\textbf{Equality Constraints}}: Constraints in a linear programming problem that must be satisfied exactly.%
\item \href{https://www.math.purdue.edu/~yipn/421/2024-09-17-DualityGeneral.pdf#page=7}{\textbf{Inequality Constraints}}: Constraints in a linear programming problem that allow for less-than-or-equal-to relationships.%
\item \href{https://www.math.purdue.edu/~yipn/421/2024-09-17-DualityGeneral.pdf#page=11}{\textbf{Feasible Region}}: The set of all points that satisfy the constraints of a linear programming problem.%
\item \href{https://www.math.purdue.edu/~yipn/421/2024-09-19-DualityExamples.pdf#page=1}{\textbf{Primal Problem}}: A linear programming problem formulated to maximize or minimize an objective function subject to a set of constraints.%
\item \href{https://www.math.purdue.edu/~yipn/421/2024-09-19-DualityExamples.pdf#page=1}{\textbf{Dual Problem}}: A linear programming problem derived from a primal problem, offering an alternative perspective on the same optimization problem.%
\item \href{https://www.math.purdue.edu/~yipn/421/2024-09-19-DualityExamples.pdf#page=2}{\textbf{Diet Problem}}: A classic linear programming example where the goal is to minimize the cost of a diet while meeting minimum nutritional requirements.%
\item \href{https://www.math.purdue.edu/~yipn/421/2024-09-19-DualityExamples.pdf#page=1}{\textbf{Negative Transpose Property}}: The relationship between the constraint matrices and vectors of the primal and dual problems, where the dual's constraint matrix is the transpose of the primal's, with signs reversed for inequality constraints.%
\item \href{https://www.math.purdue.edu/~yipn/421/2024-09-24-SimplexMatrix.pdf#page=1}{\textbf{Objective Function}}: A function to be maximized or minimized in a linear programming problem. %
\item \href{https://www.math.purdue.edu/~yipn/421/2024-09-24-SimplexMatrix.pdf#page=1}{\textbf{Constraints}}: Restrictions or limitations on the values of the variables in a linear programming problem, expressed as inequalities or equalities. %
\item \href{https://www.math.purdue.edu/~yipn/421/2024-09-24-SimplexMatrix.pdf#page=1}{\textbf{Slack Variable}}: A non-negative variable added to an inequality constraint to transform it into an equality constraint. %
\item \href{https://www.math.purdue.edu/~yipn/421/2024-09-24-SimplexMatrix.pdf#page=4}{\textbf{Basic Variable}}: A variable in a linear programming problem that is assigned a value determined by the current basis in the simplex method. %
\item \href{https://www.math.purdue.edu/~yipn/421/2024-09-24-SimplexMatrix.pdf#page=4}{\textbf{Non{-}basic Variable}}: A variable in a linear programming problem that is set to zero in the current basis in the simplex method. %
\item \href{https://www.math.purdue.edu/~yipn/421/2024-09-24-SimplexMatrix.pdf#page=1}{\textbf{Simplex Method}}: An iterative algorithm used to solve linear programming problems by moving from one feasible solution to another, improving the objective function value at each step. %
\item \href{https://www.math.purdue.edu/~yipn/421/2024-09-24-SimplexMatrix.pdf#page=17}{\textbf{Pivot Operation}}: A step in the simplex method that involves selecting a pivot element in the simplex tableau and performing row operations to update the tableau. %
\item \href{https://www.math.purdue.edu/~yipn/421/2024-09-24-SimplexMatrix.pdf#page=16}{\textbf{Feasible Region}}: The set of all points that satisfy the constraints of a linear programming problem. %
\item \href{https://www.math.purdue.edu/~yipn/421/2024-09-24-SimplexMatrix.pdf#page=16}{\textbf{Graphical Representation}}: A visual method for solving linear programming problems by plotting the constraints and the objective function on a graph. %
\item \href{https://www.math.purdue.edu/~yipn/421/2024-09-24-SimplexMatrix.pdf#page=1}{\textbf{Matrix Representation}}: A way to represent a linear programming problem using matrices and vectors, which is particularly useful for implementing the simplex method. %
\item \href{https://www.math.purdue.edu/~yipn/421/2024-09-24-SimplexMatrix.pdf#page=14}{\textbf{Complementary Slackness}}: A condition that relates the primal and dual solutions of a linear programming problem. %
\item \href{https://www.math.purdue.edu/~yipn/421/2024-09-24-SimplexMatrix.pdf#page=14}{\textbf{Primal Problem}}: The original linear programming problem. %
\item \href{https://www.math.purdue.edu/~yipn/421/2024-09-24-SimplexMatrix.pdf#page=14}{\textbf{Dual Problem}}: A related linear programming problem derived from the primal problem. %
\item \href{https://www.math.purdue.edu/~yipn/421/2024-09-24-SimplexMatrix.pdf#page=1}{\textbf{Equality Constraints}}: Constraints in a linear programming problem expressed as equalities. %
\item \href{https://www.math.purdue.edu/~yipn/421/2024-09-24-SimplexMatrix.pdf#page=1}{\textbf{Inequality Constraints}}: Constraints in a linear programming problem expressed as inequalities. %
\item \href{https://www.math.purdue.edu/~yipn/421/2024-09-24-SimplexMatrix.pdf#page=15}{\textbf{Dictionary}}: A representation of the simplex tableau used to track the values of the variables and the objective function during the simplex method iterations. %
\item \href{https://www.math.purdue.edu/~yipn/421/2024-09-24-SimplexMatrix.pdf#page=15}{\textbf{Auxiliary Problem}}: A modified linear programming problem used to find an initial feasible solution for the simplex method when the origin is not feasible. %
\item \href{https://www.math.purdue.edu/~yipn/421/2024-09-26-Sensitivity.pdf#page=1}{\textbf{Sensitivity Analysis<L{[}2024{-}09{-}26{-}Sensitivity{]} 1><L{[}2024{-}09{-}26{-}Sensitivity{]} 4><L{[}2024{-}09{-}26{-}Sensitivity{]} 7><L{[}2024{-}09{-}26{-}Sensitivity{]} 8><L{[}2024{-}09{-}26{-}Sensitivity{]} 11>}}: The study of how changes in the objective function coefficients or constraint values affect the optimal solution of a linear programming problem.%
\item \href{https://www.math.purdue.edu/~yipn/421/2024-09-26-Sensitivity.pdf#page=1}{\textbf{Parametric Analysis<L{[}2024{-}09{-}26{-}Sensitivity{]} 1><L{[}2024{-}09{-}26{-}Sensitivity{]} 4><L{[}2024{-}09{-}26{-}Sensitivity{]} 8>}}: A technique used in linear programming to analyze the effect of systematically changing one or more parameters (e.g., objective function coefficients, constraint values) on the optimal solution.%
\item \href{https://www.math.purdue.edu/~yipn/421/2024-09-26-Sensitivity.pdf#page=2}{\textbf{Optimal Solution<L{[}2024{-}09{-}26{-}Sensitivity{]} 2><L{[}2024{-}09{-}26{-}Sensitivity{]} 3><L{[}2024{-}09{-}26{-}Sensitivity{]} 5><L{[}2024{-}09{-}26{-}Sensitivity{]} 6><L{[}2024{-}09{-}26{-}Sensitivity{]} 9><L{[}2024{-}09{-}26{-}Sensitivity{]} 10>}}: The set of values for the decision variables that maximizes (or minimizes) the objective function while satisfying all constraints of a linear programming problem.%
\item \href{https://www.math.purdue.edu/~yipn/421/2024-09-26-Sensitivity.pdf#page=2}{\textbf{Basic Variables<L{[}2024{-}09{-}26{-}Sensitivity{]} 2><L{[}2024{-}09{-}26{-}Sensitivity{]} 3><L{[}2024{-}09{-}26{-}Sensitivity{]} 7><L{[}2024{-}09{-}26{-}Sensitivity{]} 11>}}: In the simplex method, the variables associated with the columns forming the identity matrix in the tableau. Their values are directly determined by the current solution. %
\item \href{https://www.math.purdue.edu/~yipn/421/2024-09-26-Sensitivity.pdf#page=2}{\textbf{Non{-}basic Variables<L{[}2024{-}09{-}26{-}Sensitivity{]} 2><L{[}2024{-}09{-}26{-}Sensitivity{]} 3><L{[}2024{-}09{-}26{-}Sensitivity{]} 7><L{[}2024{-}09{-}26{-}Sensitivity{]} 11>}}: In the simplex method, the variables whose values are set to zero in a given iteration. %
\item \href{https://www.math.purdue.edu/~yipn/421/2024-09-26-Sensitivity.pdf#page=7}{\textbf{Reduced Costs<L{[}2024{-}09{-}26{-}Sensitivity{]} 7>}}: In the simplex method, the values that indicate how much the objective function value would change if a non-basic variable were to become basic.%
\item \href{https://www.math.purdue.edu/~yipn/421/2024-09-26-Sensitivity.pdf#page=8}{\textbf{Right{-}hand Side (RHS) Values<L{[}2024{-}09{-}26{-}Sensitivity{]} 8><L{[}2024{-}09{-}26{-}Sensitivity{]} 11>}}: The constants on the right-hand side of the constraints in a linear programming problem.%
\item \href{https://www.math.purdue.edu/~yipn/421/2024-09-26-Sensitivity.pdf#page=1}{\textbf{Objective Function<L{[}2024{-}09{-}26{-}Sensitivity{]} 1><L{[}2024{-}09{-}26{-}Sensitivity{]} 4><L{[}2024{-}09{-}26{-}Sensitivity{]} 8>}}: The function to be maximized or minimized in a linear programming problem. It represents the goal of the optimization. %
\item \href{https://www.math.purdue.edu/~yipn/421/2024-09-26-Sensitivity.pdf#page=1}{\textbf{Constraints<L{[}2024{-}09{-}26{-}Sensitivity{]} 1><L{[}2024{-}09{-}26{-}Sensitivity{]} 4><L{[}2024{-}09{-}26{-}Sensitivity{]} 8>}}: The limitations or restrictions on the values of the decision variables in a linear programming problem. %
\item \href{https://www.math.purdue.edu/~yipn/421/2024-09-26-Sensitivity.pdf#page=1}{\textbf{Decision Variables<L{[}2024{-}09{-}26{-}Sensitivity{]} 1>}}: The unknown quantities that are to be determined in a linear programming problem. Their values are optimized to achieve the objective. %
\item \href{https://www.math.purdue.edu/~yipn/421/2024-09-26-Sensitivity.pdf#page=2}{\textbf{Slack Variables<L{[}2024{-}09{-}26{-}Sensitivity{]} 2><L{[}2024{-}09{-}26{-}Sensitivity{]} 3><L{[}2024{-}09{-}26{-}Sensitivity{]} 5><L{[}2024{-}09{-}26{-}Sensitivity{]} 6><L{[}2024{-}09{-}26{-}Sensitivity{]} 9><L{[}2024{-}09{-}26{-}Sensitivity{]} 10>}}: Variables added to inequality constraints to transform them into equalities. They represent the unused resources. %
\item \href{https://www.math.purdue.edu/~yipn/421/2024-09-26-Sensitivity.pdf#page=4}{\textbf{Parameter λ<L{[}2024{-}09{-}26{-}Sensitivity{]} 4>}}: A parameter introduced to represent a change in the objective function coefficient.%
\item \href{https://www.math.purdue.edu/~yipn/421/2024-09-26-Sensitivity.pdf#page=8}{\textbf{Parameters Δbᵢ<L{[}2024{-}09{-}26{-}Sensitivity{]} 8>}}: Parameters introduced to represent changes in the right-hand side values of the constraints.%
\item \href{https://www.math.purdue.edu/~yipn/421/2024-09-26-Sensitivity.pdf#page=7}{\textbf{Basis Matrix (B)<L{[}2024{-}09{-}26{-}Sensitivity{]} 7><L{[}2024{-}09{-}26{-}Sensitivity{]} 11>}}: A matrix formed by the columns of the constraint matrix corresponding to the basic variables.%
\item \href{https://www.math.purdue.edu/~yipn/421/2024-09-26-Sensitivity.pdf#page=7}{\textbf{Inverse of the Basis Matrix (B⁻¹)<L{[}2024{-}09{-}26{-}Sensitivity{]} 7><L{[}2024{-}09{-}26{-}Sensitivity{]} 11>}}: The inverse of the basis matrix, used in calculations within the simplex method.%
\item \href{https://www.math.purdue.edu/~yipn/421/2024-09-26-Sensitivity.pdf#page=3}{\textbf{Matrix Operations<L{[}2024{-}09{-}26{-}Sensitivity{]} 3><L{[}2024{-}09{-}26{-}Sensitivity{]} 7><L{[}2024{-}09{-}26{-}Sensitivity{]} 11>}}: Mathematical operations involving matrices, used extensively in the simplex method.%
\item \href{https://www.math.purdue.edu/~yipn/421/2024-09-26-Sensitivity.pdf#page=3}{\textbf{Simplex Tableau<L{[}2024{-}09{-}26{-}Sensitivity{]} 3><L{[}2024{-}09{-}26{-}Sensitivity{]} 7>}}: A tabular representation of the linear programming problem used in the simplex method.%
\item \href{https://www.math.purdue.edu/~yipn/421/2024-09-26-Sensitivity.pdf#page=2}{\textbf{Dictionary<L{[}2024{-}09{-}26{-}Sensitivity{]} 2><L{[}2024{-}09{-}26{-}Sensitivity{]} 3><L{[}2024{-}09{-}26{-}Sensitivity{]} 5><L{[}2024{-}09{-}26{-}Sensitivity{]} 6><L{[}2024{-}09{-}26{-}Sensitivity{]} 9><L{[}2024{-}09{-}26{-}Sensitivity{]} 10>}}: A representation of the optimal solution in terms of basic and non-basic variables.%
\item \href{https://www.math.purdue.edu/~yipn/421/2024-10-01-DualGeneralLP.pdf#page=1}{\textbf{Dual of General LP Problem}}: The optimization problem derived from a primal linear program, where the objective function is minimized and the constraints are related to the primal problem's constraints through transposition and inequality reversal. %
\item \href{https://www.math.purdue.edu/~yipn/421/2024-10-01-DualGeneralLP.pdf#page=2}{\textbf{General LP}}: A linear programming problem where the decision variables are not necessarily constrained to be non-negative. %
\item \href{https://www.math.purdue.edu/~yipn/421/2024-10-01-DualGeneralLP.pdf#page=4}{\textbf{Lagrangian Function}}: A function that combines the objective function and constraints of an optimization problem, incorporating Lagrange multipliers to handle the constraints. %
\item \href{https://www.math.purdue.edu/~yipn/421/2024-10-01-DualGeneralLP.pdf#page=5}{\textbf{Weak Duality}}: A theorem stating that the objective function value of any feasible solution to the primal problem is less than or equal to the objective function value of any feasible solution to the dual problem. %
\item \href{https://www.math.purdue.edu/~yipn/421/2024-10-01-DualGeneralLP.pdf#page=1}{\textbf{Primal Problem}}: The original linear programming problem, typically a maximization problem. %
\item \href{https://www.math.purdue.edu/~yipn/421/2024-10-01-DualGeneralLP.pdf#page=1}{\textbf{Dual Problem}}: The optimization problem derived from a primal linear program, typically a minimization problem. %
\item \href{https://www.math.purdue.edu/~yipn/421/2024-10-01-DualGeneralLP.pdf#page=1}{\textbf{Objective Function}}: The function to be maximized or minimized in a linear programming problem. %
\item \href{https://www.math.purdue.edu/~yipn/421/2024-10-01-DualGeneralLP.pdf#page=1}{\textbf{Constraints}}: Restrictions or limitations on the decision variables in a linear programming problem. %
\item \href{https://www.math.purdue.edu/~yipn/421/2024-10-01-DualGeneralLP.pdf#page=1}{\textbf{Decision Variables}}: The variables whose values are to be determined in a linear programming problem. %
\item \href{https://www.math.purdue.edu/~yipn/421/2024-10-01-DualGeneralLP.pdf#page=2}{\textbf{Basic Variables}}: A subset of variables in a linear program that are used to express the other variables (non-basic variables) in terms of them. %
\item \href{https://www.math.purdue.edu/~yipn/421/2024-10-01-DualGeneralLP.pdf#page=2}{\textbf{Non{-}basic Variables}}: Variables in a linear program that are set to zero in a basic feasible solution. %
\item \href{https://www.math.purdue.edu/~yipn/421/2024-10-01-DualGeneralLP.pdf#page=2}{\textbf{Slack Variables}}: Variables added to inequality constraints to transform them into equality constraints. %
\item \textbf{Feasible Region}: The set of all points that satisfy the constraints of a linear programming problem.%
\item \href{https://www.math.purdue.edu/~yipn/421/2024-10-01-DualGeneralLP.pdf#page=2}{\textbf{Optimal Solution}}: A feasible solution that optimizes (maximizes or minimizes) the objective function. %
\item \href{https://www.math.purdue.edu/~yipn/421/2024-10-01-SensitivityDuality.pdf#page=1}{\textbf{Primal Problem}}: A linear programming problem in its standard form, typically aiming to maximize an objective function subject to constraints. %
\item \href{https://www.math.purdue.edu/~yipn/421/2024-10-01-SensitivityDuality.pdf#page=2}{\textbf{Dual Problem}}: A linear programming problem derived from the primal problem, typically aiming to minimize an objective function subject to constraints. The optimal solution of the dual problem provides valuable information about the primal problem. %
\item \href{https://www.math.purdue.edu/~yipn/421/2024-10-01-SensitivityDuality.pdf#page=3}{\textbf{Complementary Slackness}}: A theorem stating that at the optimal solution of a linear programming problem, either a constraint is binding (holds with equality) or the corresponding dual variable is zero. %
\item \href{https://www.math.purdue.edu/~yipn/421/2024-10-01-SensitivityDuality.pdf#page=1}{\textbf{Sensitivity Analysis}}: The study of how changes in the parameters of a linear programming problem (e.g., objective function coefficients, constraint values) affect the optimal solution. %
\item \href{https://www.math.purdue.edu/~yipn/421/2024-10-01-SensitivityDuality.pdf#page=10}{\textbf{Shadow Price}}: The rate of change in the optimal objective function value for a unit change in the right-hand side of a constraint. It's represented by the optimal dual variable. %
\item \href{https://www.math.purdue.edu/~yipn/421/2024-10-01-SensitivityDuality.pdf#page=7}{\textbf{Basis Matrix (B)}}: A square matrix formed by the columns of the constraint matrix corresponding to the basic variables in a linear programming problem. %
\item \href{https://www.math.purdue.edu/~yipn/421/2024-10-01-SensitivityDuality.pdf#page=7}{\textbf{Inverse of the Basis Matrix (B⁻¹)}}: The inverse of the basis matrix, used in the simplex method to update the solution and calculate the optimal values. %
\item \href{https://www.math.purdue.edu/~yipn/421/2024-10-01-SensitivityDuality.pdf#page=2}{\textbf{Basic Variables}}: Variables that form a basis in the simplex method. Their values are determined by solving a system of linear equations. %
\item \href{https://www.math.purdue.edu/~yipn/421/2024-10-01-SensitivityDuality.pdf#page=2}{\textbf{Non{-}basic Variables}}: Variables that are not in the basis in the simplex method. Their values are set to zero. %
\item \href{https://www.math.purdue.edu/~yipn/421/2024-10-01-SensitivityDuality.pdf#page=7}{\textbf{Reduced Costs}}: In the simplex method, the values that indicate how much the objective function would improve if a non-basic variable were to enter the basis. %
\item \href{https://www.math.purdue.edu/~yipn/421/2024-10-01-SensitivityDuality.pdf#page=8}{\textbf{Right{-}Hand Side (RHS) Values}}: The values on the right-hand side of the inequality or equality constraints in a linear programming problem. %
\item \href{https://www.math.purdue.edu/~yipn/421/2024-10-01-SensitivityDuality.pdf#page=1}{\textbf{Objective Function}}: The function to be maximized or minimized in a linear programming problem. %
\item \href{https://www.math.purdue.edu/~yipn/421/2024-10-01-SensitivityDuality.pdf#page=1}{\textbf{Constraints}}: Restrictions or limitations on the values of the decision variables in a linear programming problem. %
\item \href{https://www.math.purdue.edu/~yipn/421/2024-10-01-SensitivityDuality.pdf#page=1}{\textbf{Decision Variables}}: The variables whose values are to be determined in a linear programming problem. %
\item \href{https://www.math.purdue.edu/~yipn/421/2024-10-01-SensitivityDuality.pdf#page=5}{\textbf{Slack Variables}}: Non-negative variables added to inequality constraints to transform them into equalities. %
\item \href{https://www.math.purdue.edu/~yipn/421/2024-10-03-DualSimplexMatrix.pdf#page=2}{\textbf{Primal Problem}}: The original linear programming problem, typically a maximization problem. %
\item \href{https://www.math.purdue.edu/~yipn/421/2024-10-03-DualSimplexMatrix.pdf#page=2}{\textbf{Dual Problem}}: A linear programming problem derived from the primal problem, typically a minimization problem. The optimal solution of the dual problem provides valuable information about the primal problem. %
\item \href{https://www.math.purdue.edu/~yipn/421/2024-10-03-DualSimplexMatrix.pdf#page=1}{\textbf{Strong Duality}}: The theorem stating that if both the primal and dual linear programs have feasible solutions, then their optimal objective function values are equal. %
\item \href{https://www.math.purdue.edu/~yipn/421/2024-10-03-DualSimplexMatrix.pdf#page=1}{\textbf{Weak Duality}}: The theorem stating that the objective function value of any feasible solution to the dual problem is always less than or equal to the objective function value of any feasible solution to the primal problem. %
\item \href{https://www.math.purdue.edu/~yipn/421/2024-10-03-DualSimplexMatrix.pdf#page=2}{\textbf{Complementary Slackness}}: A condition that relates the optimal solutions of the primal and dual problems. It states that if a constraint in the primal problem is not binding at the optimum, then the corresponding dual variable is zero, and vice versa. %
\item \href{https://www.math.purdue.edu/~yipn/421/2024-10-03-DualSimplexMatrix.pdf#page=2}{\textbf{Dual Simplex Method}}: An algorithm for solving linear programming problems that starts with a dual feasible solution and iteratively improves it until an optimal solution is found. %
\item \href{https://www.math.purdue.edu/~yipn/421/2024-10-03-DualSimplexMatrix.pdf#page=7}{\textbf{Simplex Method}}: An iterative algorithm for solving linear programming problems. It moves from one feasible solution to another, improving the objective function value at each step, until an optimal solution is reached. %
\item \href{https://www.math.purdue.edu/~yipn/421/2024-10-03-DualSimplexMatrix.pdf#page=7}{\textbf{Basis Matrix (B)}}: A matrix formed by the columns of the constraint matrix corresponding to the basic variables in a linear programming problem. %
\item \href{https://www.math.purdue.edu/~yipn/421/2024-10-03-DualSimplexMatrix.pdf#page=7}{\textbf{Inverse of the Basis Matrix (B⁻¹)}}: The inverse of the basis matrix, used in the simplex method to efficiently update the solution at each iteration. %
\item \href{https://www.math.purdue.edu/~yipn/421/2024-10-03-DualSimplexMatrix.pdf#page=2}{\textbf{Basic Variables}}: Variables that form a basis in the simplex method. Their values are determined by the constraints. %
\item \href{https://www.math.purdue.edu/~yipn/421/2024-10-03-DualSimplexMatrix.pdf#page=2}{\textbf{Non{-}basic Variables}}: Variables that are not in the basis in the simplex method. Their values are set to zero. %
\item \href{https://www.math.purdue.edu/~yipn/421/2024-10-03-DualSimplexMatrix.pdf#page=7}{\textbf{Reduced Costs}}: Values that indicate how much the objective function would change if a non-basic variable were to enter the basis. %
\item \href{https://www.math.purdue.edu/~yipn/421/2024-10-03-DualSimplexMatrix.pdf#page=1}{\textbf{Objective Function}}: The function that is being maximized or minimized in a linear programming problem. %
\item \href{https://www.math.purdue.edu/~yipn/421/2024-10-03-DualSimplexMatrix.pdf#page=1}{\textbf{Constraints}}: Restrictions or limitations on the values of the decision variables in a linear programming problem. %
\item \href{https://www.math.purdue.edu/~yipn/421/2024-10-03-DualSimplexMatrix.pdf#page=2}{\textbf{Optimal Solution}}: A feasible solution to a linear programming problem that yields the best possible value of the objective function. %
\item \href{https://www.math.purdue.edu/~yipn/421/2024-10-03-DualSimplexMatrix.pdf#page=2}{\textbf{Slack Variables}}: Variables added to inequality constraints to transform them into equality constraints. %
\item \href{https://www.math.purdue.edu/~yipn/421/2024-10-03-DualSimplexMatrix.pdf#page=7}{\textbf{Matrix Representation of Linear Programs}}: A way to represent linear programming problems using matrices and vectors, which is useful for implementing algorithms like the simplex method. %
\item \href{https://www.math.purdue.edu/~yipn/421/2024-10-03-DualSimplexMatrix.pdf#page=8}{\textbf{Right{-}Hand Side (RHS) Values}}: The values on the right-hand side of the constraints in a linear programming problem. %
\item \href{https://www.math.purdue.edu/~yipn/421/2024-10-03-DualSimplexMatrix.pdf#page=1}{\textbf{Sensitivity Analysis}}: The study of how changes in the parameters of a linear programming problem affect the optimal solution. %
\item \href{https://www.math.purdue.edu/~yipn/421/2024-10-03-DualSimplexMatrix.pdf#page=1}{\textbf{Parametric Analysis}}: A type of sensitivity analysis that examines the effect of systematically changing one or more parameters of a linear programming problem. %
\item \href{https://www.math.purdue.edu/~yipn/421/2024-10-03-DualSimplexMatrix.pdf#page=2}{\textbf{Negative Transpose Property}}: The relationship between the coefficient matrices of the primal and dual problems. The coefficient matrix of the dual problem is the negative transpose of the coefficient matrix of the primal problem. %
\item \href{https://www.math.purdue.edu/~yipn/421/2024-10-03-DualSimplexMatrix.pdf#page=4}{\textbf{Lagrangian Function}}: A function used in optimization to incorporate constraints into the objective function. %
\item \href{https://www.math.purdue.edu/~yipn/421/2024-10-03-DualSimplexMatrix.pdf#page=7}{\textbf{Partitioning Variables into Basic and Non{-}basic Sets}}: The process of dividing the variables in a linear programming problem into basic and non-basic sets, which is a key step in the simplex method. %
\item \href{https://www.math.purdue.edu/~yipn/421/2024-10-03-DualSimplexMatrix.pdf#page=7}{\textbf{Expressing the Objective Function in Terms of Basic and Non{-}basic Variables}}: A technique used in the simplex method to simplify the objective function and express it in terms of only the non-basic variables. %
\item \href{https://www.math.purdue.edu/~yipn/421/2024-10-03-DualSimplexMatrix.pdf#page=11}{\textbf{Updating the Simplex Tableau Using Matrix Operations}}: The process of updating the simplex tableau using matrix operations, which is a key step in the simplex method. %
\item \href{https://www.math.purdue.edu/~yipn/421/2024-10-22-Convexity.pdf#page=1}{\textbf{Convex Set}}: A set $S \subseteq \mathbb{R}^m$ such that for any $z_1, z_2 \in S$, $\lambda z_1 + (1-\lambda) z_2 \in S$ for $0 \le \lambda \le 1$.%
\item \href{https://www.math.purdue.edu/~yipn/421/2024-10-22-Convexity.pdf#page=1}{\textbf{Convex Combination}}: A weighted sum of points $z_1, z_2, \dots, z_n \in \mathbb{R}^m$ of the form $\sum_{i=1}^n \lambda_i z_i$, where $\lambda_i \ge 0$ and $\sum_{i=1}^n \lambda_i = 1$.%
\item \href{https://www.math.purdue.edu/~yipn/421/2024-10-22-Convexity.pdf#page=5}{\textbf{Convex Hull}}: The smallest convex set containing a given set $S \subseteq \mathbb{R}^m$, denoted as $\text{Conv}(S) = \bigcap \{C \text{ convex } | S \subseteq C\}$.%
\item \href{https://www.math.purdue.edu/~yipn/421/2024-10-22-Convexity.pdf#page=8}{\textbf{Caratheodory Theorem}}: If $z \in \text{Conv}(S)$, then $z$ can be written as a convex combination of at most $m+1$ points from $S$.%
\item \href{https://www.math.purdue.edu/~yipn/421/2024-10-22-Convexity.pdf#page=10}{\textbf{Separation Theorem}}: If $P$ and $\bar{P}$ are two disjoint nonempty polyhedra in $\mathbb{R}^n$, then there exist disjoint half-spaces $H$ and $\bar{H}$ such that $P \subset H$ and $\bar{P} \subset \bar{H}$.%
\item \href{https://www.math.purdue.edu/~yipn/421/2024-10-22-Convexity.pdf#page=12}{\textbf{Farkas' Lemma}}: The system $Ax \le b$ has no solution if and only if there exists a vector $y \ge 0$ such that $y^T A = 0$ and $y^T b < 0$.%
\item \href{https://www.math.purdue.edu/~yipn/421/2024-10-22-Regression.pdf#page=1}{\textbf{Mean}}: The average of a set of numbers. %
\item \href{https://www.math.purdue.edu/~yipn/421/2024-10-22-Regression.pdf#page=1}{\textbf{Median}}: The middle value in a sorted set of numbers. %
\item \href{https://www.math.purdue.edu/~yipn/421/2024-10-22-Regression.pdf#page=1}{\textbf{Mid{-}range}}: The average of the minimum and maximum values in a set of numbers. %
\item \href{https://www.math.purdue.edu/~yipn/421/2024-10-22-Regression.pdf#page=2}{\textbf{Linear Regression}}: A statistical method to model the relationship between a dependent variable and one or more independent variables by fitting a linear equation to observed data. %
\item \href{https://www.math.purdue.edu/~yipn/421/2024-10-22-Regression.pdf#page=4}{\textbf{L² Regression (Least Square)}}: A type of linear regression that minimizes the sum of the squares of the differences between the observed and predicted values. %
\item \href{https://www.math.purdue.edu/~yipn/421/2024-10-22-Regression.pdf#page=5}{\textbf{L¹ Regression}}: A type of linear regression that minimizes the sum of the absolute differences between the observed and predicted values. %
\item \href{https://www.math.purdue.edu/~yipn/421/2024-10-22-Regression.pdf#page=6}{\textbf{L∞ Regression}}: A type of linear regression that minimizes the maximum absolute difference between the observed and predicted values. %
\item \href{https://www.math.purdue.edu/~yipn/421/2024-10-22-Regression.pdf#page=7}{\textbf{Binary Classification}}: A type of supervised machine learning where the goal is to classify data points into two categories. %
\item \href{https://www.math.purdue.edu/~yipn/421/2024-10-22-Regression.pdf#page=8}{\textbf{Linear Separability}}: The property of a dataset where its points can be perfectly separated into two categories by a hyperplane. %
\item \href{https://www.math.purdue.edu/~yipn/421/2024-10-22-Regression.pdf#page=9}{\textbf{Maximum Margin Classifier}}: A classifier that finds the hyperplane that maximizes the distance to the nearest data points of each class. %
\item \href{https://www.math.purdue.edu/~yipn/421/2024-10-22-Regression.pdf#page=10}{\textbf{Support Vector Machine (SVM)}}: A supervised machine learning algorithm that uses a hyperplane to maximize the margin between classes. %
\item \href{https://www.math.purdue.edu/~yipn/421/2024-10-22-Regression.pdf#page=10}{\textbf{Support Vectors}}: The data points closest to the hyperplane in an SVM. %
\item \href{https://www.math.purdue.edu/~yipn/421/2024-10-22-Regression.pdf#page=15}{\textbf{Maximum Inscribed Circle in a Convex Polyhedron}}: The problem of finding the largest circle that can fit inside a convex polyhedron. %
\item \href{https://www.math.purdue.edu/~yipn/421/2024-10-22-Regression.pdf#page=22}{\textbf{Smallest Ball Problem}}: The problem of finding the smallest-radius ball that encloses a set of points in a d-dimensional space. %
\item \href{https://www.math.purdue.edu/~yipn/421/2024-10-22-Regression.pdf#page=7}{\textbf{Hyperplane}}: A subspace of one dimension less than its ambient space. %
\item \href{https://www.math.purdue.edu/~yipn/421/2024-10-28-ApplsLinProg.pdf#page=1}{\textbf{Monthly Demands}}: The predicted sales of ice cream for each month of the year.%
\item \href{https://www.math.purdue.edu/~yipn/421/2024-10-28-ApplsLinProg.pdf#page=2}{\textbf{Average Demand}}: The average monthly sales of ice cream over the entire year.%
\item \href{https://www.math.purdue.edu/~yipn/421/2024-10-28-ApplsLinProg.pdf#page=3}{\textbf{Production at month i (Xᵢ)}}: The amount of ice cream produced in month i.%
\item \href{https://www.math.purdue.edu/~yipn/421/2024-10-28-ApplsLinProg.pdf#page=3}{\textbf{Surplus at end of Month i (Sᵢ)}}: The amount of ice cream left over at the end of month i.%
\item \href{https://www.math.purdue.edu/~yipn/421/2024-10-28-ApplsLinProg.pdf#page=3}{\textbf{Cost of Changing Production}}: The cost associated with altering the production level from one month to the next.%
\item \href{https://www.math.purdue.edu/~yipn/421/2024-10-28-ApplsLinProg.pdf#page=3}{\textbf{Cost of Storage}}: The cost of storing surplus ice cream.%
\item \href{https://www.math.purdue.edu/~yipn/421/2024-10-28-ApplsLinProg.pdf#page=9}{\textbf{Annual Return of ith Investment (Rᵢ)}}: The return on investment i over a year.%
\item \href{https://www.math.purdue.edu/~yipn/421/2024-10-28-ApplsLinProg.pdf#page=9}{\textbf{Proportion of Investment in i (Xᵢ)}}: The fraction of the total investment allocated to investment i.%
\item \href{https://www.math.purdue.edu/~yipn/421/2024-10-28-ApplsLinProg.pdf#page=10}{\textbf{Fluctuation (R̃ⱼ)}}: The deviation of the return of investment j from its expected return.%
\item \href{https://www.math.purdue.edu/~yipn/421/2024-10-28-ApplsLinProg.pdf#page=17}{\textbf{Smallest Enclosing Ball Problem}}: Finding the smallest ball (hypersphere) that contains a given set of points.%
\item \href{https://www.math.purdue.edu/~yipn/421/2024-10-28-ApplsLinProg.pdf#page=19}{\textbf{Maximum Weight Matching}}: Finding a matching in a weighted graph that maximizes the total weight of the edges in the matching.%
\item \href{https://www.math.purdue.edu/~yipn/421/2024-10-28-ApplsLinProg.pdf#page=25}{\textbf{Cutting Stock Problem}}: Determining how to cut large rolls of material into smaller rolls of specified widths to minimize waste.%
\item \href{https://www.math.purdue.edu/~yipn/421/2024-10-28-ApplsLinProg.pdf#page=30}{\textbf{Sparse Solution of Underdetermined System of Linear Equations}}: Finding a solution to a system of linear equations with fewer equations than unknowns, such that the solution has a limited number of non-zero elements.%
\item \href{https://www.math.purdue.edu/~yipn/421/2024-10-28-Farkas.pdf#page=1}{\textbf{Farkas' Lemma}}: A fundamental theorem in linear programming stating conditions for the existence or non-existence of solutions to systems of linear inequalities. %
\item \href{https://www.math.purdue.edu/~yipn/421/2024-10-28-Farkas.pdf#page=1}{\textbf{Primal Problem}}: In linear programming, the original optimization problem to be solved. %
\item \href{https://www.math.purdue.edu/~yipn/421/2024-10-28-Farkas.pdf#page=1}{\textbf{Dual Problem}}: In linear programming, a related optimization problem derived from the primal problem, often providing valuable insights and alternative solution methods. %
\item \href{https://www.math.purdue.edu/~yipn/421/2024-10-28-Farkas.pdf#page=1}{\textbf{Objective Function}}: The function to be maximized or minimized in a linear programming problem. %
\item \href{https://www.math.purdue.edu/~yipn/421/2024-10-28-Farkas.pdf#page=1}{\textbf{Constraints}}: Restrictions or limitations on the decision variables in a linear programming problem, expressed as linear inequalities or equations. %
\item \href{https://www.math.purdue.edu/~yipn/421/2024-10-28-Farkas.pdf#page=1}{\textbf{Decision Variables}}: The unknown quantities in a linear programming problem that are to be determined to optimize the objective function. %
\item \href{https://www.math.purdue.edu/~yipn/421/2024-10-28-Farkas.pdf#page=2}{\textbf{Feasible Region}}: The set of all points that satisfy all the constraints of a linear programming problem. %
\item \href{https://www.math.purdue.edu/~yipn/421/2024-10-28-Farkas.pdf#page=2}{\textbf{Optimal Solution}}: A solution to a linear programming problem that achieves the best possible value of the objective function within the feasible region. %
\item \href{https://www.math.purdue.edu/~yipn/421/2024-10-28-Farkas.pdf#page=2}{\textbf{Complementary Slackness}}: A condition that relates the optimal solutions of the primal and dual problems in linear programming. %
\item \href{https://www.math.purdue.edu/~yipn/421/2024-10-28-Farkas.pdf#page=2}{\textbf{Negative Transpose Property}}: The relationship between the coefficient matrices of the primal and dual problems in linear programming. %
\item \href{https://www.math.purdue.edu/~yipn/421/2024-10-28-Farkas.pdf#page=4}{\textbf{Convex Set}}: A set where any line segment connecting two points in the set is entirely contained within the set. %
\item \href{https://www.math.purdue.edu/~yipn/421/2024-10-28-Farkas.pdf#page=4}{\textbf{Convex Hull}}: The smallest convex set containing a given set of points. %
\item \href{https://www.math.purdue.edu/~yipn/421/2024-10-28-Farkas.pdf#page=9}{\textbf{Hyperplane}}: A flat, n-1 dimensional subspace in n-dimensional space. %
\item \href{https://www.math.purdue.edu/~yipn/421/2024-10-28-Farkas.pdf#page=2}{\textbf{Basic Variables}}: In the simplex method, a subset of variables that form a basis for the solution space. %
\item \href{https://www.math.purdue.edu/~yipn/421/2024-10-28-Farkas.pdf#page=2}{\textbf{Non{-}basic Variables}}: In the simplex method, variables that are not in the basis and are set to zero. %
\item \href{https://www.math.purdue.edu/~yipn/421/2024-10-28-Farkas.pdf#page=7}{\textbf{Basis Matrix (B)}}: A matrix formed by the columns of the constraint matrix corresponding to the basic variables in the simplex method. %
\item \href{https://www.math.purdue.edu/~yipn/421/2024-10-28-Farkas.pdf#page=7}{\textbf{Inverse of the Basis Matrix (B⁻¹)}}: The inverse of the basis matrix, used in the simplex method to update the solution. %
\item \href{https://www.math.purdue.edu/~yipn/421/2024-10-28-Farkas.pdf#page=7}{\textbf{Reduced Costs}}: In the simplex method, values indicating the change in the objective function value if a non-basic variable is increased from zero. %
\item \href{https://www.math.purdue.edu/~yipn/421/2024-10-28-Farkas.pdf#page=2}{\textbf{Slack Variables}}: Variables added to inequality constraints to transform them into equalities. %
\item \href{https://www.math.purdue.edu/~yipn/421/2024-10-28-Farkas.pdf#page=1}{\textbf{Sensitivity Analysis}}: The study of how changes in the parameters of a linear programming problem affect the optimal solution. %
\item \href{https://www.math.purdue.edu/~yipn/421/2024-10-28-Farkas.pdf#page=4}{\textbf{Parametric Analysis}}: A type of sensitivity analysis that examines the effect of systematically changing one or more parameters of a linear programming problem. %
\item \href{https://www.math.purdue.edu/~yipn/421/2024-10-28-Farkas.pdf#page=4}{\textbf{Lagrangian Function}}: A function used in optimization to incorporate constraints into the objective function. %
\item \href{https://www.math.purdue.edu/~yipn/421/2024-10-28-Farkas.pdf#page=4}{\textbf{Lagrange Multiplier}}: A parameter in the Lagrangian function that represents the penalty for violating a constraint. %
\item \href{https://www.math.purdue.edu/~yipn/421/2024-10-31-IntLinProg.pdf#page=1}{\textbf{Integer Programming}}: A linear programming problem where some or all of the decision variables are restricted to integer values.%
\item \href{https://www.math.purdue.edu/~yipn/421/2024-10-31-IntLinProg.pdf#page=2}{\textbf{Relaxed Problem}}: The linear programming problem obtained by removing the integer constraints from an integer programming problem.%
\item \href{https://www.math.purdue.edu/~yipn/421/2024-10-31-IntLinProg.pdf#page=5}{\textbf{Branch and Bound}}: A method for solving integer programming problems by recursively partitioning the feasible region and bounding the objective function value.%
\item \href{https://www.math.purdue.edu/~yipn/421/2024-10-31-IntLinProg.pdf#page=15}{\textbf{Gomory Cuts}}: Additional constraints added to a linear programming relaxation of an integer program to cut off non-integer solutions.%
\item \href{https://www.math.purdue.edu/~yipn/421/2024-11-05-NetFlows.pdf#page=1}{\textbf{Network Flow}}: The study of flows through a network, often modeled as a directed graph with capacities and costs on the arcs. %
\item \href{https://www.math.purdue.edu/~yipn/421/2024-11-05-NetFlows.pdf#page=1}{\textbf{Nodes (N)}}: The set of points in a network flow problem representing locations or entities. %
\item \href{https://www.math.purdue.edu/~yipn/421/2024-11-05-NetFlows.pdf#page=1}{\textbf{Arcs (A)}}: The set of directed edges connecting nodes in a network flow problem, representing pathways or connections. %
\item \href{https://www.math.purdue.edu/~yipn/421/2024-11-05-NetFlows.pdf#page=2}{\textbf{Supply/Demand (\$b\_i\$)}}: The net flow at each node in a network, positive values representing supply and negative values representing demand. %
\item \href{https://www.math.purdue.edu/~yipn/421/2024-11-05-NetFlows.pdf#page=2}{\textbf{Flow (\$X\_\{ij\}\$)}}: The amount of flow transported from node $i$ to node $j$ along an arc. %
\item \href{https://www.math.purdue.edu/~yipn/421/2024-11-05-NetFlows.pdf#page=2}{\textbf{Cost (\$C\_\{ij\}\$)}}: The cost of transporting one unit of flow from node $i$ to node $j$ along an arc. %
\item \href{https://www.math.purdue.edu/~yipn/421/2024-11-05-NetFlows.pdf#page=2}{\textbf{Objective Function}}: The function to be minimized or maximized in a linear programming problem, often representing total cost or profit. %
\item \href{https://www.math.purdue.edu/~yipn/421/2024-11-05-NetFlows.pdf#page=3}{\textbf{Balanced Equation}}: An equation representing the conservation of flow at each node in a network, stating that the inflow plus supply equals the outflow. %
\item \href{https://www.math.purdue.edu/~yipn/421/2024-11-05-NetFlows.pdf#page=4}{\textbf{Incidence Matrix (A)}}: A matrix representing the structure of a network, where rows correspond to nodes and columns correspond to arcs. %
\item \href{https://www.math.purdue.edu/~yipn/421/2024-11-05-NetFlows.pdf#page=6}{\textbf{Rank(A)}}: The rank of the incidence matrix, indicating the number of linearly independent equations in the network flow problem. %
\item \href{https://www.math.purdue.edu/~yipn/421/2024-11-05-NetFlows.pdf#page=7}{\textbf{Primal Problem (P)}}: The original minimization problem in linear programming. %
\item \href{https://www.math.purdue.edu/~yipn/421/2024-11-05-NetFlows.pdf#page=7}{\textbf{Dual Problem (D)}}: The maximization problem associated with the primal problem in linear programming. %
\item \href{https://www.math.purdue.edu/~yipn/421/2024-11-05-NetFlows.pdf#page=10}{\textbf{Path}}: A sequence of connected nodes in a network. %
\item \href{https://www.math.purdue.edu/~yipn/421/2024-11-05-NetFlows.pdf#page=10}{\textbf{Cycle}}: A path that starts and ends at the same node. %
\item \href{https://www.math.purdue.edu/~yipn/421/2024-11-05-NetFlows.pdf#page=10}{\textbf{Tree}}: A connected graph with no cycles. %
\item \href{https://www.math.purdue.edu/~yipn/421/2024-11-05-NetFlows.pdf#page=10}{\textbf{Spanning Tree}}: A tree that includes all nodes in a network. %
\item \href{https://www.math.purdue.edu/~yipn/421/2024-11-05-NetFlows.pdf#page=11}{\textbf{Tree Solution}}: A solution to a network flow problem where the flow along arcs in a spanning tree is zero. %
\item \href{https://www.math.purdue.edu/~yipn/421/2024-11-05-NetFlows.pdf#page=15}{\textbf{Basis Matrix (B)}}: A square submatrix of the incidence matrix corresponding to a spanning tree. %
\item \href{https://www.math.purdue.edu/~yipn/421/2024-11-05-NetFlows.pdf#page=18}{\textbf{Complementary Slackness}}: A condition relating primal and dual variables in linear programming, indicating optimality. %
\item \href{https://www.math.purdue.edu/~yipn/421/2024-11-05-NetFlows.pdf#page=22}{\textbf{Primal Simplex Method}}: An iterative algorithm for solving linear programming problems. %
\item \href{https://www.math.purdue.edu/~yipn/421/2024-11-07-NetSimplex.pdf#page=1}{\textbf{Network Simplex Method}}: A specialized simplex method for solving linear programs representing network flow problems. %
\item \href{https://www.math.purdue.edu/~yipn/421/2024-11-07-NetSimplex.pdf#page=1}{\textbf{Network (N, A)}}: A representation of a network flow problem as a graph, where N is the set of nodes and A is the set of arcs (edges). %
\item \href{https://www.math.purdue.edu/~yipn/421/2024-11-07-NetSimplex.pdf#page=1}{\textbf{Spanning Tree T = (N, Ā)}}: A connected, acyclic subgraph of a network that includes all nodes. %
\item \href{https://www.math.purdue.edu/~yipn/421/2024-11-07-NetSimplex.pdf#page=2}{\textbf{Primal Flow}}: The flow values assigned to the arcs in a network flow problem that satisfy the supply/demand constraints. %
\item \href{https://www.math.purdue.edu/~yipn/421/2024-11-07-NetSimplex.pdf#page=3}{\textbf{Dual Flow}}: The dual variables associated with the nodes in a network flow problem that satisfy the dual constraints. %
\item \href{https://www.math.purdue.edu/~yipn/421/2024-11-07-NetSimplex.pdf#page=3}{\textbf{Cost (cij)}}: The cost associated with the flow along arc (i,j) in a network flow problem. %
\item \href{https://www.math.purdue.edu/~yipn/421/2024-11-07-NetSimplex.pdf#page=6}{\textbf{Entering Arc}}: The arc selected to enter the spanning tree during an iteration of the network simplex method. %
\item \href{https://www.math.purdue.edu/~yipn/421/2024-11-07-NetSimplex.pdf#page=6}{\textbf{Leaving Arc}}: The arc selected to leave the spanning tree during an iteration of the network simplex method. %
\item \href{https://www.math.purdue.edu/~yipn/421/2024-11-07-NetSimplex.pdf#page=4}{\textbf{Cycle}}: A closed path in a network. %
\item \href{https://www.math.purdue.edu/~yipn/421/2024-11-07-NetSimplex.pdf#page=4}{\textbf{Path}}: A sequence of connected arcs in a network. %
\item \href{https://www.math.purdue.edu/~yipn/421/2024-11-07-NetSimplex.pdf#page=1}{\textbf{Nodes (N)}}: The set of points in a network. %
\item \href{https://www.math.purdue.edu/~yipn/421/2024-11-07-NetSimplex.pdf#page=1}{\textbf{Arcs (A)}}: The set of directed edges in a network. %
\item \href{https://www.math.purdue.edu/~yipn/421/2024-11-07-NetSimplex.pdf#page=5}{\textbf{Primal Network Simplex}}: The primal simplex method applied to network flow problems. %
\item \href{https://www.math.purdue.edu/~yipn/421/2024-11-07-NetSimplex.pdf#page=32}{\textbf{Dual Network Simplex}}: The dual simplex method applied to network flow problems. %
\item \href{https://www.math.purdue.edu/~yipn/421/2024-11-07-NetSimplex.pdf#page=1}{\textbf{Tree}}: A connected, acyclic graph. %
\item \href{https://www.math.purdue.edu/~yipn/421/2024-11-07-NetSimplex.pdf#page=2}{\textbf{Flow (xij)}}: The amount of flow along arc (i,j) in a network flow problem. %
\item \href{https://www.math.purdue.edu/~yipn/421/2024-11-07-NetSimplex.pdf#page=2}{\textbf{Basic Variables}}: Variables associated with arcs in the spanning tree. %
\item \href{https://www.math.purdue.edu/~yipn/421/2024-11-07-NetSimplex.pdf#page=2}{\textbf{Non{-}basic Variables}}: Variables associated with arcs not in the spanning tree. %
\item \href{https://www.math.purdue.edu/~yipn/421/2024-11-07-NetSimplex.pdf#page=3}{\textbf{Dual Problem (D)}}: The dual linear program associated with a given primal linear program. %
\item \href{https://www.math.purdue.edu/~yipn/421/2024-11-07-NetSimplex.pdf#page=3}{\textbf{Primal Problem (P)}}: The original linear program. %
\item \href{https://www.math.purdue.edu/~yipn/421/2024-11-07-NetSimplex.pdf#page=1}{\textbf{Objective Function}}: The function to be minimized or maximized in a linear program. %
\item \href{https://www.math.purdue.edu/~yipn/421/2024-11-07-NetSimplex.pdf#page=1}{\textbf{Constraints}}: The limitations or restrictions on the variables in a linear program. %
\item \href{https://www.math.purdue.edu/~yipn/421/2024-11-07-NetSimplex.pdf#page=2}{\textbf{Optimal Solution}}: A feasible solution that optimizes the objective function. %
\item \href{https://www.math.purdue.edu/~yipn/421/2024-11-07-NetSimplex.pdf#page=2}{\textbf{Supply/Demand (bi)}}: The supply or demand at node i in a network flow problem. %
\item \href{https://www.math.purdue.edu/~yipn/421/2024-11-07-NetSimplex.pdf#page=4}{\textbf{Parametric Analysis}}: A technique for analyzing how the optimal solution changes as problem parameters vary. %
\item \href{https://www.math.purdue.edu/~yipn/421/2024-11-07-NetSimplex.pdf#page=7}{\textbf{Basis Matrix (B)}}: A matrix formed by the columns of the constraint matrix corresponding to the basic variables. %
\item \href{https://www.math.purdue.edu/~yipn/421/2024-11-07-NetSimplex.pdf#page=7}{\textbf{Inverse of the Basis Matrix (B⁻¹)}}: The inverse of the basis matrix, used in simplex calculations. %
\item \href{https://www.math.purdue.edu/~yipn/421/2024-11-07-NetSimplex.pdf#page=7}{\textbf{Reduced Costs}}: The costs of non-basic variables after accounting for the current basis. %
\item \href{https://www.math.purdue.edu/~yipn/421/2024-11-07-NetSimplex.pdf#page=35}{\textbf{Dual Slack}}: The difference between the left-hand side and right-hand side of a dual constraint. %
\item \href{https://www.math.purdue.edu/~yipn/421/2024-11-14-NetAppls.pdf#page=1}{\textbf{Transportation Problem}}: A linear programming problem where the goal is to minimize the cost of transporting goods from multiple sources to multiple destinations, subject to supply and demand constraints.%
\item \href{https://www.math.purdue.edu/~yipn/421/2024-11-14-NetAppls.pdf#page=3}{\textbf{Dummy Node}}: A node added to a network flow model to handle excess supply or unmet demand, converting inequality constraints into equalities.%
\item \href{https://www.math.purdue.edu/~yipn/421/2024-11-14-NetAppls.pdf#page=7}{\textbf{Incidence Matrix}}: A matrix representing the connections in a network, where rows correspond to nodes and columns correspond to arcs.%
\item \href{https://www.math.purdue.edu/~yipn/421/2024-11-14-NetAppls.pdf#page=9}{\textbf{Upper Bounded Transshipment Problem}}: A network flow problem where the flow on each arc is constrained by an upper bound.%
\item \href{https://www.math.purdue.edu/~yipn/421/2024-11-14-NetAppls.pdf#page=13}{\textbf{Shortest Path Problem}}: A graph problem where the goal is to find the path with the minimum total weight or cost between two nodes in a network.%
\item \href{https://www.math.purdue.edu/~yipn/421/2024-11-14-NetAppls.pdf#page=16}{\textbf{Bellman's Equation}}: A recursive equation used to solve the shortest path problem by iteratively finding the shortest distance from each node to the destination node.%
\item \href{https://www.math.purdue.edu/~yipn/421/2024-11-14-NetAppls.pdf#page=20}{\textbf{Dijkstra's Algorithm}}: An algorithm for finding the shortest paths from a single source node to all other nodes in a graph.%
\item \href{https://www.math.purdue.edu/~yipn/421/2024-11-14-NetAppls.pdf#page=2}{\textbf{Slack Variable}}: A variable added to an inequality constraint to convert it into an equality constraint.%
\item \href{https://www.math.purdue.edu/~yipn/421/2024-11-14-NetAppls.pdf#page=1}{\textbf{Objective Function}}: The function to be minimized or maximized in a linear programming problem.%
\item \href{https://www.math.purdue.edu/~yipn/421/2024-11-14-NetAppls.pdf#page=1}{\textbf{Constraints}}: Restrictions or limitations on the values of the decision variables in a linear programming problem.%
\item \href{https://www.math.purdue.edu/~yipn/421/2024-11-14-NetAppls.pdf#page=1}{\textbf{Decision Variables}}: Variables representing the unknowns in a linear programming problem.%
\item \href{https://www.math.purdue.edu/~yipn/421/2024-11-14-NetAppls.pdf#page=7}{\textbf{Primal Problem}}: The original linear programming problem.%
\item \href{https://www.math.purdue.edu/~yipn/421/2024-11-14-NetAppls.pdf#page=7}{\textbf{Dual Problem}}: A linear programming problem derived from the primal problem, providing a lower bound on the optimal value of the primal problem.%
\item \href{https://www.math.purdue.edu/~yipn/421/2024-11-14-NetAppls.pdf#page=7}{\textbf{Complementary Slackness}}: A condition that must hold at the optimal solution of a linear programming problem and its dual.%
\item \href{https://www.math.purdue.edu/~yipn/421/2024-11-14-NetAppls.pdf#page=1}{\textbf{Network Flow}}: A type of linear programming problem where the goal is to find the optimal flow of goods or resources through a network.%
\item \href{https://www.math.purdue.edu/~yipn/421/2024-11-14-NetAppls.pdf#page=1}{\textbf{Nodes}}: Points or locations in a network.%
\item \href{https://www.math.purdue.edu/~yipn/421/2024-11-14-NetAppls.pdf#page=1}{\textbf{Arcs}}: Connections or edges between nodes in a network.%
\item \href{https://www.math.purdue.edu/~yipn/421/2024-11-14-NetAppls.pdf#page=2}{\textbf{Flow}}: The amount of goods or resources passing through an arc in a network.%
\item \href{https://www.math.purdue.edu/~yipn/421/2024-11-14-NetAppls.pdf#page=2}{\textbf{Supply/Demand (\$b\_i\$)}}: The amount of supply at a source node or demand at a sink node in a network flow problem.%
\item \href{https://www.math.purdue.edu/~yipn/421/2024-11-14-NetAppls.pdf#page=6}{\textbf{Cost (\$c\_\{ij\}\$)}}: The cost associated with transporting goods or resources along an arc in a network flow problem.%
\item \href{https://www.math.purdue.edu/~yipn/421/2024-11-14-NetAppls.pdf#page=2}{\textbf{Basic Variables}}: Variables that form a basis in the simplex method.%
\item \href{https://www.math.purdue.edu/~yipn/421/2024-11-14-NetAppls.pdf#page=2}{\textbf{Non{-}basic Variables}}: Variables that are not in the basis in the simplex method.%
\item \href{https://www.math.purdue.edu/~yipn/421/2024-11-14-NetAppls.pdf#page=7}{\textbf{Basis Matrix (B)}}: A matrix formed by the columns of the constraint matrix corresponding to the basic variables.%
\item \href{https://www.math.purdue.edu/~yipn/421/2024-11-14-NetAppls.pdf#page=7}{\textbf{Inverse of the Basis Matrix (B⁻¹)}}: The inverse of the basis matrix, used in the simplex method.%
\item \href{https://www.math.purdue.edu/~yipn/421/2024-11-14-NetAppls.pdf#page=7}{\textbf{Reduced Costs}}: The costs associated with bringing non-basic variables into the basis in the simplex method.%
\item \href{https://www.math.purdue.edu/~yipn/421/2024-11-14-NetAppls.pdf#page=7}{\textbf{Simplex Method}}: An iterative algorithm for solving linear programming problems.%
\item \href{https://www.math.purdue.edu/~yipn/421/2024-11-14-NetAppls.pdf#page=7}{\textbf{Dual Simplex Method}}: A variant of the simplex method that starts with a dual feasible solution.%
\item \href{https://www.math.purdue.edu/~yipn/421/2024-11-14-NetAppls.pdf#page=7}{\textbf{Primal Simplex Method}}: A variant of the simplex method that starts with a primal feasible solution.%
\item \href{https://www.math.purdue.edu/~yipn/421/2024-11-14-NetAppls.pdf#page=7}{\textbf{Network Simplex Method}}: A specialized simplex method for solving network flow problems.%
\item \href{https://www.math.purdue.edu/~yipn/421/2024-11-14-NetAppls.pdf#page=7}{\textbf{Spanning Tree}}: A tree that connects all nodes in a network.%
\item \href{https://www.math.purdue.edu/~yipn/421/2024-11-14-NetAppls.pdf#page=7}{\textbf{Convex Set}}: A set where any line segment connecting two points in the set is also contained in the set.%
\item \href{https://www.math.purdue.edu/~yipn/421/2024-11-14-NetAppls.pdf#page=7}{\textbf{Convex Hull}}: The smallest convex set containing a given set of points.%
\item \href{https://www.math.purdue.edu/~yipn/421/2024-11-14-NetAppls.pdf#page=7}{\textbf{Hyperplane}}: A flat, n-dimensional subspace in an (n+1)-dimensional space.%
\item \href{https://www.math.purdue.edu/~yipn/421/2024-11-14-NetAppls.pdf#page=7}{\textbf{Linear Separability}}: The ability to separate two sets of points with a hyperplane.%
\item \href{https://www.math.purdue.edu/~yipn/421/2024-11-14-NetAppls.pdf#page=7}{\textbf{Support Vectors}}: The data points closest to the hyperplane in a support vector machine (SVM).%
\item \href{https://www.math.purdue.edu/~yipn/421/2024-11-14-NetAppls.pdf#page=7}{\textbf{Support Vector Machine (SVM)}}: A machine learning algorithm that uses hyperplanes to perform binary classification.%
\item \href{https://www.math.purdue.edu/~yipn/421/2024-11-14-NetAppls.pdf#page=7}{\textbf{Lagrangian Function}}: A function used in optimization problems to incorporate constraints into the objective function.%
\item \href{https://www.math.purdue.edu/~yipn/421/2024-11-14-NetAppls.pdf#page=7}{\textbf{Lagrange Multiplier}}: A parameter in the Lagrangian function that represents the weight of a constraint.%
\item \href{https://www.math.purdue.edu/~yipn/421/2024-11-14-NetAppls.pdf#page=7}{\textbf{Parameter λ}}: A parameter used in optimization problems, often representing a trade-off between different objectives.%
\item \href{https://www.math.purdue.edu/~yipn/421/2024-11-14-NetAppls.pdf#page=13}{\textbf{Path}}: A sequence of arcs connecting two nodes in a network.%
\item \href{https://www.math.purdue.edu/~yipn/421/2024-11-14-NetAppls.pdf#page=7}{\textbf{Balanced Equation}}: An equation representing the conservation of flow in a network.%
\item \href{https://www.math.purdue.edu/~yipn/421/2024-11-14-NetAppls.pdf#page=7}{\textbf{Farkas' Lemma}}: A theorem in linear programming that provides conditions for the existence of a solution to a system of linear inequalities.%
\item \href{https://www.math.purdue.edu/~yipn/421/2024-11-14-NetAppls.pdf#page=7}{\textbf{Caratheodory's Theorem}}: A theorem stating that any point in the convex hull of a set of points can be expressed as a convex combination of at most n+1 points, where n is the dimension of the space.%
\item \href{https://www.math.purdue.edu/~yipn/421/2024-11-14-NetAppls.pdf#page=7}{\textbf{Separation Theorem}}: A theorem stating that two disjoint convex sets can be separated by a hyperplane.%
\item \href{https://www.math.purdue.edu/~yipn/421/2024-11-14-NetAppls.pdf#page=7}{\textbf{Weak Duality Theorem}}: A theorem stating that the objective function value of the dual problem is always less than or equal to the objective function value of the primal problem.%
\item \href{https://www.math.purdue.edu/~yipn/421/2024-11-14-NetAppls.pdf#page=7}{\textbf{Strong Duality Theorem}}: A theorem stating that if both the primal and dual problems have feasible solutions, then their optimal objective function values are equal.%
\item \href{https://www.math.purdue.edu/~yipn/421/2024-11-14-NetAppls.pdf#page=7}{\textbf{Sensitivity Analysis}}: The study of how changes in the parameters of a linear programming problem affect the optimal solution.%
\item \href{https://www.math.purdue.edu/~yipn/421/2024-11-14-NetAppls.pdf#page=7}{\textbf{Parametric Analysis}}: A type of sensitivity analysis that examines the effect of changes in a parameter over a range of values.%
\item \href{https://www.math.purdue.edu/~yipn/421/2024-11-14-NetAppls.pdf#page=7}{\textbf{Gomory Cuts}}: Cutting planes used in integer programming to improve the linear programming relaxation.%
\item \href{https://www.math.purdue.edu/~yipn/421/2024-11-14-NetAppls.pdf#page=7}{\textbf{Branch and Bound}}: An algorithm for solving integer programming problems by recursively partitioning the feasible region.%
\item \href{https://www.math.purdue.edu/~yipn/421/2024-11-14-NetAppls.pdf#page=7}{\textbf{Cutting Stock Problem}}: A problem of cutting large rolls of material into smaller pieces to meet specific demands.%
\item \href{https://www.math.purdue.edu/~yipn/421/2024-11-14-NetAppls.pdf#page=7}{\textbf{Knapsack Problem}}: A problem of selecting a subset of items with weights and values to maximize the total value while not exceeding a weight limit.%
\item \href{https://www.math.purdue.edu/~yipn/421/2024-11-14-NetAppls.pdf#page=7}{\textbf{Diet Problem}}: A problem of selecting a combination of foods to meet nutritional requirements at minimum cost.%
\item \href{https://www.math.purdue.edu/~yipn/421/2024-11-14-NetAppls.pdf#page=7}{\textbf{Portfolio Selection (Markowitz Model)}}: A problem of selecting a portfolio of investments to maximize expected return while minimizing risk.%
\item \href{https://www.math.purdue.edu/~yipn/421/2024-11-14-NetAppls.pdf#page=7}{\textbf{Machine Scheduling}}: A problem of assigning jobs to machines to minimize the total completion time.%
\item \href{https://www.math.purdue.edu/~yipn/421/2024-11-14-NetAppls.pdf#page=7}{\textbf{Workforce Planning}}: A problem of determining the optimal number of workers to hire to meet demand.%
\item \href{https://www.math.purdue.edu/~yipn/421/2024-11-14-NetAppls.pdf#page=7}{\textbf{Ice Cream Production Planning}}: A problem of determining the optimal production plan for different ice cream flavors to meet demand and minimize costs.%
\item \href{https://www.math.purdue.edu/~yipn/421/2024-11-14-NetAppls.pdf#page=7}{\textbf{L¹ Regression}}: A type of regression that minimizes the sum of absolute errors.%
\item \href{https://www.math.purdue.edu/~yipn/421/2024-11-14-NetAppls.pdf#page=7}{\textbf{L² Regression (Least Square)}}: A type of regression that minimizes the sum of squared errors.%
\item \href{https://www.math.purdue.edu/~yipn/421/2024-11-14-NetAppls.pdf#page=7}{\textbf{L∞ Regression}}: A type of regression that minimizes the maximum absolute error.%
\item \href{https://www.math.purdue.edu/~yipn/421/2024-11-14-NetAppls.pdf#page=7}{\textbf{Matrix Representation}}: Representing a linear program using matrices and vectors.%
\item \href{https://www.math.purdue.edu/~yipn/421/2024-11-14-NetAppls.pdf#page=7}{\textbf{Right{-}Hand Side (RHS) Values}}: The values on the right-hand side of the constraints in a linear programming problem.%
\item \href{https://www.math.purdue.edu/~yipn/421/2024-11-14-NetAppls.pdf#page=7}{\textbf{Auxiliary Problem}}: A problem used to find an initial feasible solution for a linear program.%
\item \href{https://www.math.purdue.edu/~yipn/421/2024-11-14-NetAppls.pdf#page=7}{\textbf{Relaxed Problem}}: A linear programming problem obtained by relaxing the integrality constraints of an integer programming problem.%
\item \href{https://www.math.purdue.edu/~yipn/421/2024-11-14-NetAppls.pdf#page=2}{\textbf{Optimal Solution}}: The solution that optimizes the objective function while satisfying all constraints.%
\item \href{https://www.math.purdue.edu/~yipn/421/2024-11-14-NetAppls.pdf#page=7}{\textbf{Graphical Solution Method}}: A method for solving linear programming problems graphically.%
\item \href{https://www.math.purdue.edu/~yipn/421/2024-11-14-NetAppls.pdf#page=7}{\textbf{Pivot Operation}}: An operation in the simplex method that moves from one basic feasible solution to another.%
\item \href{https://www.math.purdue.edu/~yipn/421/2024-11-14-NetAppls.pdf#page=7}{\textbf{Entering Arc}}: An arc that enters the basis in the network simplex method.%
\item \href{https://www.math.purdue.edu/~yipn/421/2024-11-14-NetAppls.pdf#page=7}{\textbf{Leaving Arc}}: An arc that leaves the basis in the network simplex method.%
\item \href{https://www.math.purdue.edu/~yipn/421/2024-11-14-NetAppls.pdf#page=7}{\textbf{Cycle}}: A closed path in a network.%
\item \href{https://www.math.purdue.edu/~yipn/421/2024-11-14-NetAppls.pdf#page=7}{\textbf{Tree}}: A connected graph without cycles.%
\item \href{https://www.math.purdue.edu/~yipn/421/2024-11-14-NetAppls.pdf#page=7}{\textbf{Tree Solution}}: A solution to a network flow problem represented as a tree.%
\item \href{https://www.math.purdue.edu/~yipn/421/2024-11-14-NetAppls.pdf#page=7}{\textbf{Surplus at End of Month i (sᵢ)}}: The amount of excess production at the end of month i.%
\item \href{https://www.math.purdue.edu/~yipn/421/2024-11-14-NetAppls.pdf#page=7}{\textbf{Production at Month i (xᵢ)}}: The amount of production during month i.%
\item \href{https://www.math.purdue.edu/~yipn/421/2024-11-14-NetAppls.pdf#page=7}{\textbf{Monthly Demands}}: The demand for a product in each month.%
\item \href{https://www.math.purdue.edu/~yipn/421/2024-11-14-NetAppls.pdf#page=7}{\textbf{Cost of Storage}}: The cost of storing excess production.%
\item \href{https://www.math.purdue.edu/~yipn/421/2024-11-14-NetAppls.pdf#page=7}{\textbf{Cost of Changing Production}}: The cost associated with changing the production level from one month to the next.%
\item \href{https://www.math.purdue.edu/~yipn/421/2024-11-14-NetAppls.pdf#page=7}{\textbf{Annual Return of ith Investment (rᵢ)}}: The annual return of the ith investment in a portfolio.%
\item \href{https://www.math.purdue.edu/~yipn/421/2024-11-14-NetAppls.pdf#page=7}{\textbf{Proportion of Investment in i (xᵢ)}}: The proportion of the total investment allocated to the ith investment.%
\item \href{https://www.math.purdue.edu/~yipn/421/2024-11-14-NetAppls.pdf#page=7}{\textbf{Fluctuation (r̃ⱼ)}}: The fluctuation of the jth variable.%
\item \href{https://www.math.purdue.edu/~yipn/421/2024-11-14-NetAppls.pdf#page=7}{\textbf{Shadow Price}}: The change in the optimal objective function value per unit change in a constraint's right-hand side value.%
\item \href{https://www.math.purdue.edu/~yipn/421/2024-11-14-NetAppls.pdf#page=7}{\textbf{Average Demand}}: The average demand for a product over a period of time.%
\item \href{https://www.math.purdue.edu/~yipn/421/2024-11-14-NetAppls.pdf#page=7}{\textbf{Median}}: The middle value in a sorted dataset.%
\item \href{https://www.math.purdue.edu/~yipn/421/2024-11-14-NetAppls.pdf#page=7}{\textbf{Mean}}: The average value in a dataset.%
\item \href{https://www.math.purdue.edu/~yipn/421/2024-11-14-NetAppls.pdf#page=7}{\textbf{Smallest Enclosing Ball}}: The smallest sphere that encloses a given set of points.%
\item \href{https://www.math.purdue.edu/~yipn/421/2024-11-14-NetAppls.pdf#page=7}{\textbf{Maximum Inscribed Circle in a Convex Polyhedron}}: The largest circle that can be inscribed within a convex polyhedron.%
\item \href{https://www.math.purdue.edu/~yipn/421/2024-11-14-NetAppls.pdf#page=7}{\textbf{Maximum Weight Matching}}: A problem of finding a matching in a graph that maximizes the sum of weights of the edges in the matching.%
\item \href{https://www.math.purdue.edu/~yipn/421/2024-11-14-NetAppls.pdf#page=7}{\textbf{Rank(A)}}: The rank of matrix A.%
\item \href{https://www.math.purdue.edu/~yipn/421/2024-11-14-NetAppls.pdf#page=7}{\textbf{Lexicographic Method}}: A method for resolving degeneracy in the simplex method.%
\item \href{https://www.math.purdue.edu/~yipn/421/2024-11-14-NetAppls.pdf#page=7}{\textbf{Bland's Rule}}: A rule for selecting entering and leaving variables in the simplex method to prevent cycling.%
\item \href{https://www.math.purdue.edu/~yipn/421/2024-11-14-NetAppls.pdf#page=7}{\textbf{Largest{-}Coefficient Rule}}: A rule for selecting the entering variable in the simplex method.%
\item \href{https://www.math.purdue.edu/~yipn/421/2024-11-14-NetAppls.pdf#page=7}{\textbf{Largest{-}Increase Rule}}: A rule for selecting the entering variable in the simplex method.%
\item \href{https://www.math.purdue.edu/~yipn/421/2024-11-14-NetAppls.pdf#page=7}{\textbf{Randomized Pivot Rule}}: A rule for selecting the entering and leaving variables in the simplex method randomly.%
\item \href{https://www.math.purdue.edu/~yipn/421/2024-11-14-NetAppls.pdf#page=7}{\textbf{Iterative Improvement Using the Simplex Method}}: The process of iteratively improving the solution in the simplex method.%
\item \href{https://www.math.purdue.edu/~yipn/421/2024-11-14-NetAppls.pdf#page=7}{\textbf{Updating the Simplex Tableau Using Matrix Operations}}: The process of updating the simplex tableau using matrix operations.%
\item \href{https://www.math.purdue.edu/~yipn/421/2024-11-14-NetAppls.pdf#page=7}{\textbf{Expressing the Objective Function in Terms of Basic and Non{-}Basic Variables}}: The process of rewriting the objective function in terms of basic and non-basic variables.%
\item \href{https://www.math.purdue.edu/~yipn/421/2024-11-14-NetAppls.pdf#page=7}{\textbf{Analyzing the Optimal Tableau and Its Properties}}: The process of analyzing the optimal tableau and its properties.%
\item \href{https://www.math.purdue.edu/~yipn/421/2024-11-14-NetAppls.pdf#page=7}{\textbf{Understanding Complementary Variables and Their Relationship in Simplex Iterations}}: The process of understanding complementary variables and their relationship in simplex iterations.%
\item \href{https://www.math.purdue.edu/~yipn/421/2024-11-14-NetAppls.pdf#page=7}{\textbf{Effect of Variable Scaling on Simplex Method Performance}}: The effect of scaling variables on the performance of the simplex method.%
\item \href{https://www.math.purdue.edu/~yipn/421/2024-11-14-NetAppls.pdf#page=7}{\textbf{Empirical vs. Theoretical Analysis of Simplex Method}}: Comparing empirical and theoretical analyses of the simplex method.%
\item \href{https://www.math.purdue.edu/~yipn/421/2024-11-14-NetAppls.pdf#page=7}{\textbf{Understanding Worst{-}Case Behavior of the Simplex Method}}: Analyzing the worst-case behavior of the simplex method.%
\item \href{https://www.math.purdue.edu/~yipn/421/2024-11-14-NetAppls.pdf#page=7}{\textbf{Smoothed Complexity}}: A measure of the complexity of an algorithm that takes into account the distribution of the input data.%
\item \href{https://www.math.purdue.edu/~yipn/421/2024-11-14-NetAppls.pdf#page=7}{\textbf{Impact of Pivot Rules on Simplex Method Efficiency}}: The impact of different pivot rules on the efficiency of the simplex method.%
\item \href{https://www.math.purdue.edu/~yipn/421/2024-11-14-NetAppls.pdf#page=7}{\textbf{Analyzing the Relationship Between Primal and Dual Problems}}: The process of analyzing the relationship between primal and dual problems.%
\item \href{https://www.math.purdue.edu/~yipn/421/2024-11-14-NetAppls.pdf#page=7}{\textbf{Analyzing the Relationship Between Primal and Dual Objective Function Values}}: The process of analyzing the relationship between primal and dual objective function values.%
\item \href{https://www.math.purdue.edu/~yipn/421/2024-11-14-NetAppls.pdf#page=7}{\textbf{Analyzing the Incidence Matrix of a Network}}: The process of analyzing the incidence matrix of a network.%
\item \href{https://www.math.purdue.edu/~yipn/421/2024-11-14-NetAppls.pdf#page=7}{\textbf{Determining the Four Possible Outcomes of Primal{-}Dual Problem Pairs}}: The process of determining the four possible outcomes of primal-dual problem pairs.%
\item \href{https://www.math.purdue.edu/~yipn/421/2024-11-14-NetAppls.pdf#page=7}{\textbf{Handling Infeasible Origins in Dual Based Phase I Algorithm}}: The process of handling infeasible origins in the dual based phase I algorithm.%
\item \href{https://www.math.purdue.edu/~yipn/421/2024-11-14-NetAppls.pdf#page=7}{\textbf{Handling Infeasible Origins}}: The process of handling infeasible origins in a linear program.%
\item \href{https://www.math.purdue.edu/~yipn/421/2024-11-14-NetAppls.pdf#page=7}{\textbf{Finding a Dual Feasible Solution}}: The process of finding a dual feasible solution.%
\item \href{https://www.math.purdue.edu/~yipn/421/2024-11-14-NetAppls.pdf#page=7}{\textbf{Finding a Primal Feasible Flow}}: The process of finding a primal feasible flow.%
\item \href{https://www.math.purdue.edu/~yipn/421/2024-11-14-NetAppls.pdf#page=7}{\textbf{Finding Primal and Dual Solutions Using Spanning Trees}}: The process of finding primal and dual solutions using spanning trees.%
\item \href{https://www.math.purdue.edu/~yipn/421/2024-11-14-NetAppls.pdf#page=7}{\textbf{Finding Primal Flows Using a Spanning Tree}}: The process of finding primal flows using a spanning tree.%
\item \href{https://www.math.purdue.edu/~yipn/421/2024-11-14-NetAppls.pdf#page=7}{\textbf{Finding an Initial Feasible Solution}}: The process of finding an initial feasible solution for a linear program.%
\item \href{https://www.math.purdue.edu/~yipn/421/2024-11-14-NetAppls.pdf#page=7}{\textbf{Formulating Network Flow Problems as Linear Programs}}: The process of formulating network flow problems as linear programs.%
\item \href{https://www.math.purdue.edu/~yipn/421/2024-11-14-NetAppls.pdf#page=7}{\textbf{Formulating Linear Programs in Matrix Notation}}: The process of formulating linear programs in matrix notation.%
\item \href{https://www.math.purdue.edu/~yipn/421/2024-11-14-NetAppls.pdf#page=7}{\textbf{Formulating Primal and Dual Linear Programs}}: The process of formulating primal and dual linear programs.%
\item \href{https://www.math.purdue.edu/~yipn/421/2024-11-14-NetAppls.pdf#page=7}{\textbf{Deriving the Dual Problem from the Lagrangian}}: The process of deriving the dual problem from the Lagrangian.%
\item \href{https://www.math.purdue.edu/~yipn/421/2024-11-14-NetAppls.pdf#page=7}{\textbf{Deriving the Dual of a General LP}}: The process of deriving the dual of a general linear program.%
\item \href{https://www.math.purdue.edu/~yipn/421/2024-11-14-NetAppls.pdf#page=7}{\textbf{Deriving the Dual of a General LP with Upper Bounds}}: The process of deriving the dual of a general linear program with upper bounds.%
\item \href{https://www.math.purdue.edu/~yipn/421/2024-11-14-NetAppls.pdf#page=7}{\textbf{Representing Linear Programs in Matrix Form}}: The process of representing linear programs in matrix form.%
\item \href{https://www.math.purdue.edu/~yipn/421/2024-11-14-NetAppls.pdf#page=7}{\textbf{Solving Linear Programs Using the Basis and Non{-}Basis Method}}: The process of solving linear programs using the basis and non-basis method.%
\item \href{https://www.math.purdue.edu/~yipn/421/2024-11-14-NetAppls.pdf#page=7}{\textbf{Solving Integer Programming Problems by Relaxing and Rounding}}: The process of solving integer programming problems by relaxing and rounding.%
\item \href{https://www.math.purdue.edu/~yipn/421/2024-11-14-NetAppls.pdf#page=7}{\textbf{Solving Linear Programs Graphically}}: The process of solving linear programs graphically.%
\item \href{https://www.math.purdue.edu/~yipn/421/2024-11-14-NetAppls.pdf#page=7}{\textbf{Applying the Simplex Method in Matrix Form}}: The process of applying the simplex method in matrix form.%
\item \href{https://www.math.purdue.edu/~yipn/421/2024-11-14-NetAppls.pdf#page=7}{\textbf{Applying the Primal Simplex Method to Network Flows}}: The process of applying the primal simplex method to network flows.%
\item \href{https://www.math.purdue.edu/~yipn/421/2024-11-14-NetAppls.pdf#page=7}{\textbf{Applying Complementary Slackness}}: The process of applying complementary slackness.%
\item \href{https://www.math.purdue.edu/~yipn/421/2024-11-14-NetAppls.pdf#page=7}{\textbf{Applying the Branch and Bound Method}}: The process of applying the branch and bound method.%
\item \href{https://www.math.purdue.edu/~yipn/421/2024-11-14-NetAppls.pdf#page=7}{\textbf{Applying Gomory Cuts to Solve Integer Programming Problems}}: The process of applying Gomory cuts to solve integer programming problems.%
\item \href{https://www.math.purdue.edu/~yipn/421/2024-11-14-NetAppls.pdf#page=7}{\textbf{Applying Weak Duality Theorem}}: The process of applying the weak duality theorem.%
\item \href{https://www.math.purdue.edu/~yipn/421/2024-11-14-NetAppls.pdf#page=7}{\textbf{Applying Caratheodory Theorem}}: The process of applying the Caratheodory theorem.%
\item \href{https://www.math.purdue.edu/~yipn/421/2024-11-14-NetAppls.pdf#page=7}{\textbf{Demonstrating the Preservation of the Negative Transpose Property}}: The process of demonstrating the preservation of the negative transpose property.%
\item \href{https://www.math.purdue.edu/~yipn/421/2024-11-14-NetAppls.pdf#page=7}{\textbf{Proving Farkas' Lemma}}: The process of proving Farkas' Lemma.%
\item \href{https://www.math.purdue.edu/~yipn/421/2024-11-14-NetAppls.pdf#page=7}{\textbf{Proving Weak Duality}}: The process of proving weak duality.%
\item \href{https://www.math.purdue.edu/~yipn/421/2024-11-14-NetAppls.pdf#page=7}{\textbf{Proving the Strong Duality Theorem}}: The process of proving the strong duality theorem.%
\item \href{https://www.math.purdue.edu/~yipn/421/2024-11-14-NetAppls.pdf#page=7}{\textbf{Proving Properties of Convex Sets}}: The process of proving properties of convex sets.%
\item \href{https://www.math.purdue.edu/~yipn/421/2024-11-14-NetAppls.pdf#page=7}{\textbf{Defining Convex Sets and Convex Combinations}}: The process of defining convex sets and convex combinations.%
\item \href{https://www.math.purdue.edu/~yipn/421/2024-11-14-NetAppls.pdf#page=7}{\textbf{Defining and Proving Properties of Convex Hulls}}: The process of defining and proving properties of convex hulls.%
\item \href{https://www.math.purdue.edu/~yipn/421/2024-11-14-NetAppls.pdf#page=7}{\textbf{Illustrating Simplex Iterations with a Numerical Example}}: The process of illustrating simplex iterations with a numerical example.%
\item \href{https://www.math.purdue.edu/~yipn/421/2024-11-14-NetAppls.pdf#page=7}{\textbf{Understanding Spanning Tree Properties}}: The process of understanding spanning tree properties.%
\item \href{https://www.math.purdue.edu/~yipn/421/2024-11-14-NetAppls.pdf#page=7}{\textbf{Understanding Complementary Variables in Simplex Iterations}}: The process of understanding complementary variables in simplex iterations.%
\item \href{https://www.math.purdue.edu/~yipn/421/2024-11-14-NetAppls.pdf#page=7}{\textbf{Verifying Optimality}}: The process of verifying optimality.%
\item \href{https://www.math.purdue.edu/~yipn/421/2024-11-14-NetAppls.pdf#page=7}{\textbf{Analyzing the Efficiency of the Simplex Method}}: The process of analyzing the efficiency of the simplex method.%
\item \href{https://www.math.purdue.edu/~yipn/421/2024-11-14-NetAppls.pdf#page=7}{\textbf{Analyzing the Sensitivity of the Objective Function Coefficients}}: The process of analyzing the sensitivity of the objective function coefficients.%
\item \href{https://www.math.purdue.edu/~yipn/421/2024-11-14-NetAppls.pdf#page=7}{\textbf{Analyzing the Sensitivity of the Constraint Values}}: The process of analyzing the sensitivity of the constraint values.%
\item \href{https://www.math.purdue.edu/~yipn/421/2024-11-14-NetAppls.pdf#page=7}{\textbf{Diet Problem Formulation and Interpretation}}: The process of formulating and interpreting the diet problem.%
\item \href{https://www.math.purdue.edu/~yipn/421/2024-11-14-NetAppls.pdf#page=7}{\textbf{Graphical Representation and Solution}}: The process of graphically representing and solving a linear program.%
\item \href{https://www.math.purdue.edu/~yipn/421/2024-11-14-NetAppls.pdf#page=7}{\textbf{Binary Classification Using a Hyperplane}}: The process of using a hyperplane for binary classification.%
\item \href{https://www.math.purdue.edu/~yipn/421/2024-11-14-NetAppls.pdf#page=7}{\textbf{Maximum Margin Classifier}}: A classifier that maximizes the margin between classes.%
\item \href{https://www.math.purdue.edu/~yipn/421/2024-11-14-NetAppls.pdf#page=7}{\textbf{SVM Optimization as a QP Problem (Version 1, 2, 3, 4)}}: Formulating SVM optimization as a quadratic programming problem.%
\item \href{https://www.math.purdue.edu/~yipn/421/2024-11-14-NetAppls.pdf#page=7}{\textbf{Lagrangian Formulation of SVM}}: Formulating SVM using the Lagrangian.%
\item \href{https://www.math.purdue.edu/~yipn/421/2024-11-14-NetAppls.pdf#page=7}{\textbf{Distance of a Point from a Hyperplane}}: The distance between a point and a hyperplane.%
\item \href{https://www.math.purdue.edu/~yipn/421/2024-11-14-NetAppls.pdf#page=7}{\textbf{Sparse Solutions of Linear Systems}}: Solutions to linear systems with many zero components.%
\item \href{https://www.math.purdue.edu/~yipn/421/2024-11-14-NetAppls.pdf#page=7}{\textbf{Sparse Solutions of Underdetermined System of Linear Equations}}: Solutions to underdetermined systems of linear equations with many zero components.%
\item \href{https://www.math.purdue.edu/~yipn/421/2024-11-14-NetAppls.pdf#page=7}{\textbf{General LP}}: A general linear programming problem.%
\item \href{https://www.math.purdue.edu/~yipn/421/2024-11-14-NetAppls.pdf#page=7}{\textbf{Primal Network Simplex}}: A primal simplex method applied to network flows.%
\item \href{https://www.math.purdue.edu/~yipn/421/2024-11-14-NetAppls.pdf#page=7}{\textbf{Dual Network Simplex}}: A dual simplex method applied to network flows.%
\item \href{https://www.math.purdue.edu/~yipn/421/2024-11-14-NetAppls.pdf#page=7}{\textbf{Developing Variants of the Network Simplex Method}}: The process of developing variants of the network simplex method.%
\item \href{https://www.math.purdue.edu/~yipn/421/2024-11-14-NetAppls.pdf#page=7}{\textbf{Dual Based Phase I Algorithm}}: A phase I algorithm for linear programming that uses the dual problem.%
\item \href{https://www.math.purdue.edu/~yipn/421/2024-11-14-NetAppls.pdf#page=7}{\textbf{Matrix Operations}}: Operations performed on matrices.%
\item \href{https://www.math.purdue.edu/~yipn/421/2024-11-14-NetAppls.pdf#page=7}{\textbf{Working with Dictionaries of Variables}}: The process of working with dictionaries of variables.%
\item \href{https://www.math.purdue.edu/~yipn/421/2024-11-14-NetAppls.pdf#page=7}{\textbf{Inequality Constraints}}: Constraints expressed as inequalities.%
\item \href{https://www.math.purdue.edu/~yipn/421/2024-11-14-NetAppls.pdf#page=7}{\textbf{Equality Constraints}}: Constraints expressed as equalities.%
\item \href{https://www.math.purdue.edu/~yipn/421/2024-11-14-NetAppls.pdf#page=7}{\textbf{Integer Programming}}: A type of optimization problem where the decision variables are restricted to integer values.%
\item \href{https://www.math.purdue.edu/~yipn/421/2024-11-14-NetAppls.pdf#page=7}{\textbf{Feasible Region}}: The set of all points that satisfy the constraints of a linear programming problem.%
\item \href{https://www.math.purdue.edu/~yipn/421/2024-11-14-NetAppls.pdf#page=7}{\textbf{Primal Flow}}: The flow in the primal network flow problem.%
\item \href{https://www.math.purdue.edu/~yipn/421/2024-11-14-NetAppls.pdf#page=7}{\textbf{Dual Flow}}: The flow in the dual network flow problem.%
\item \href{https://www.math.purdue.edu/~yipn/421/2024-11-14-NetAppls.pdf#page=7}{\textbf{Dual Slack}}: The slack variables in the dual problem.%
\item \href{https://www.math.purdue.edu/~yipn/421/2024-11-14-NetAppls.pdf#page=7}{\textbf{Theorem 10.1}}: A theorem related to linear programming.%
\item \href{https://www.math.purdue.edu/~yipn/421/2024-11-14-NetAppls.pdf#page=7}{\textbf{Theorem 10.2}}: A theorem related to linear programming.%
\item \href{https://www.math.purdue.edu/~yipn/421/2024-11-14-NetAppls.pdf#page=7}{\textbf{Theorem 14.1}}: A theorem related to linear programming.%
\item \href{https://www.math.purdue.edu/~yipn/421/2024-11-14-NetAppls.pdf#page=7}{\textbf{Mid{-}range}}: The average of the minimum and maximum values in a dataset.%
\item \href{https://www.math.purdue.edu/~yipn/421/2024-11-14-NetAppls.pdf#page=7}{\textbf{Klee{-}Minty Example}}: An example that demonstrates the worst-case behavior of the simplex method.%
\item \href{https://www.math.purdue.edu/~yipn/421/2024-11-14-NetAppls.pdf#page=7}{\textbf{Farkas' Lemma (Variant 1, 2, 3)}}: Variants of Farkas' Lemma.%
\item \href{https://www.math.purdue.edu/~yipn/421/2024-11-14-NetAppls.pdf#page=7}{\textbf{Smallest Ball Problem}}: The problem of finding the smallest ball that encloses a given set of points.%
\item \href{https://www.math.purdue.edu/~yipn/421/2024-11-14-NetAppls.pdf#page=7}{\textbf{Smallest Enclosing Ball Problem}}: The problem of finding the smallest sphere that encloses a given set of points.%
\item \href{https://www.math.purdue.edu/~yipn/421/2024-11-14-NetAppls.pdf#page=7}{\textbf{Linear Regression}}: A statistical method for modeling the relationship between a dependent variable and one or more independent variables.%
\item \href{https://www.math.purdue.edu/~yipn/421/2024-11-14-NetAppls.pdf#page=7}{\textbf{Linear Regression Using Different Error Functions}}: Using different error functions in linear regression.%
\item \href{https://www.math.purdue.edu/~yipn/421/2024-11-14-NetAppls.pdf#page=7}{\textbf{Graphical Representation}}: A visual representation of a linear program.%
\item \href{https://www.math.purdue.edu/~yipn/421/2024-11-14-NetAppls.pdf#page=7}{\textbf{Dictionary}}: A data structure used in the simplex method.%
\item \href{https://www.math.purdue.edu/~yipn/421/2024-11-14-NetAppls.pdf#page=7}{\textbf{Simplex Tableau}}: A table used in the simplex method to keep track of the variables and their values.%
\item \href{https://www.math.purdue.edu/~yipn/421/2024-11-14-NetAppls.pdf#page=7}{\textbf{Network (N, A)}}: A network with nodes N and arcs A.%
\item \href{https://www.math.purdue.edu/~yipn/421/2024-11-14-NetAppls.pdf#page=7}{\textbf{Primal Problem (P)}}: The original linear programming problem.%
\item \href{https://www.math.purdue.edu/~yipn/421/2024-11-14-NetAppls.pdf#page=7}{\textbf{Dual Problem (D)}}: The dual linear programming problem.%
\item \href{https://www.math.purdue.edu/~yipn/421/2024-11-14-NetAppls.pdf#page=7}{\textbf{Final Form of the Dual Problem with Upper Bounds}}: The final form of the dual problem when upper bounds are considered.%
\item \href{https://www.math.purdue.edu/~yipn/421/2024-11-14-NetAppls.pdf#page=7}{\textbf{Lagrangian Function (with Inequality and Equality Constraints)}}: The Lagrangian function with both inequality and equality constraints.%
\item \href{https://www.math.purdue.edu/~yipn/421/2024-11-14-NetAppls.pdf#page=7}{\textbf{Lagrangian Function for a General LP with Upper Bounds}}: The Lagrangian function for a general linear program with upper bounds.%
\item \href{https://www.math.purdue.edu/~yipn/421/2024-11-19-MaxFlowMinCut.pdf#page=1}{\textbf{Max Flow}}: The maximum amount of flow that can be sent from a source node to a sink node in a network, subject to capacity constraints on the edges. %
\item \href{https://www.math.purdue.edu/~yipn/421/2024-11-19-MaxFlowMinCut.pdf#page=2}{\textbf{Min Cut}}: A partition of the nodes in a network into two sets, one containing the source and the other containing the sink. The capacity of the cut is the sum of the capacities of the edges going from the source set to the sink set. %
\item \href{https://www.math.purdue.edu/~yipn/421/2024-11-19-MaxFlowMinCut.pdf#page=4}{\textbf{Max{-}Flow Min{-}Cut Theorem}}: The theorem stating that the maximum flow in a network is equal to the minimum capacity of any cut in the network. %
\item \href{https://www.math.purdue.edu/~yipn/421/2024-11-19-MaxFlowMinCut.pdf#page=10}{\textbf{Ford{-}Fulkerson Algorithm}}: An algorithm for finding the maximum flow in a network by iteratively finding augmenting paths and increasing the flow along these paths. %
\item \href{https://www.math.purdue.edu/~yipn/421/2024-11-19-MaxFlowMinCut.pdf#page=10}{\textbf{Augmenting Path}}: A path in a network from the source to the sink along which the flow can be increased. %
\item \href{https://www.math.purdue.edu/~yipn/421/2024-11-19-MaxFlowMinCut.pdf#page=14}{\textbf{Linear Programming Formulation (Max Flow)}}: A mathematical formulation of the maximum flow problem as a linear program, typically involving decision variables representing the flow on each edge, objective function to maximize total flow, and constraints representing capacity limits and flow conservation. %
\item \href{https://www.math.purdue.edu/~yipn/421/2024-11-21-EconNetSimplex.pdf#page=1}{\textbf{Network Simplex Method}}: A specialized simplex method for solving linear programs representing network flow problems. %
\item \href{https://www.math.purdue.edu/~yipn/421/2024-11-21-EconNetSimplex.pdf#page=2}{\textbf{Spanning Tree}}: A subgraph of a network that connects all nodes without forming any cycles. %
\item \href{https://www.math.purdue.edu/~yipn/421/2024-11-21-EconNetSimplex.pdf#page=1}{\textbf{Arcs}}: The directed edges in a network, representing connections between nodes. %
\item \href{https://www.math.purdue.edu/~yipn/421/2024-11-21-EconNetSimplex.pdf#page=1}{\textbf{Nodes}}: The points in a network, representing locations or entities. %
\item \href{https://www.math.purdue.edu/~yipn/421/2024-11-21-EconNetSimplex.pdf#page=1}{\textbf{Cost (\$c\_\{ij\}\$)}}: The cost associated with traversing arc $ij$ in a network. %
\item \href{https://www.math.purdue.edu/~yipn/421/2024-11-21-EconNetSimplex.pdf#page=1}{\textbf{Supply/Demand (\$b\_i\$)}}: The supply or demand at node $i$ in a network flow problem. %
\item \href{https://www.math.purdue.edu/~yipn/421/2024-11-21-EconNetSimplex.pdf#page=1}{\textbf{Objective Function}}: The function to be maximized or minimized in a linear program. %
\item \href{https://www.math.purdue.edu/~yipn/421/2024-11-21-EconNetSimplex.pdf#page=1}{\textbf{Constraints}}: The limitations or restrictions on the variables in a linear program. %
\item \href{https://www.math.purdue.edu/~yipn/421/2024-11-21-EconNetSimplex.pdf#page=2}{\textbf{Basic Variables}}: Variables that form a basis in the simplex method. %
\item \href{https://www.math.purdue.edu/~yipn/421/2024-11-21-EconNetSimplex.pdf#page=2}{\textbf{Non{-}basic Variables}}: Variables that are not in the basis in the simplex method. %
\item \href{https://www.math.purdue.edu/~yipn/421/2024-11-21-EconNetSimplex.pdf#page=7}{\textbf{Basis Matrix (B)}}: A matrix formed by the columns of the constraint matrix corresponding to the basic variables. %
\item \href{https://www.math.purdue.edu/~yipn/421/2024-11-21-EconNetSimplex.pdf#page=7}{\textbf{Inverse of the Basis Matrix (B⁻¹)}}: The inverse of the basis matrix, used in simplex method calculations. %
\item \href{https://www.math.purdue.edu/~yipn/421/2024-11-21-EconNetSimplex.pdf#page=7}{\textbf{Reduced Costs}}: The costs of bringing non-basic variables into the basis. %
\item \href{https://www.math.purdue.edu/~yipn/421/2024-11-21-EconNetSimplex.pdf#page=8}{\textbf{Primal Problem (P)}}: The original linear programming problem. %
\item \href{https://www.math.purdue.edu/~yipn/421/2024-11-21-EconNetSimplex.pdf#page=8}{\textbf{Dual Problem (D)}}: The problem derived from the primal problem, providing bounds on the optimal solution. %
\item \href{https://www.math.purdue.edu/~yipn/421/2024-11-21-EconNetSimplex.pdf#page=4}{\textbf{Complementary Slackness}}: A condition that must hold at optimality for primal and dual problems. %
\item \href{https://www.math.purdue.edu/~yipn/421/2024-11-21-EconNetSimplex.pdf#page=4}{\textbf{Strong Duality}}: The theorem stating that the optimal objective function values of the primal and dual problems are equal. %
\item \href{https://www.math.purdue.edu/~yipn/421/2024-11-21-EconNetSimplex.pdf#page=4}{\textbf{Weak Duality}}: The theorem stating that the objective function value of the dual problem provides an upper bound for the primal problem (for maximization problems). %
\item \href{https://www.math.purdue.edu/~yipn/421/2024-11-21-EconNetSimplex.pdf#page=1}{\textbf{Sensitivity Analysis}}: The study of how changes in the parameters of a linear program affect the optimal solution. %
\item \href{https://www.math.purdue.edu/~yipn/421/2024-11-21-EconNetSimplex.pdf#page=4}{\textbf{Path}}: A sequence of arcs connecting two nodes in a network. %
\item \href{https://www.math.purdue.edu/~yipn/421/2024-11-21-EconNetSimplex.pdf#page=4}{\textbf{Augmenting Path}}: A path in a network flow problem along which flow can be increased. %
\item \href{https://www.math.purdue.edu/~yipn/421/2024-11-21-EconNetSimplex.pdf#page=4}{\textbf{Flow (\$x\_\{ij\}\$)}}: The amount of flow along arc $ij$ in a network flow problem. %
\item \href{https://www.math.purdue.edu/~yipn/421/2024-11-21-EconNetSimplex.pdf#page=4}{\textbf{Max Flow}}: The maximum amount of flow that can be sent through a network. %
\item \href{https://www.math.purdue.edu/~yipn/421/2024-11-21-EconNetSimplex.pdf#page=4}{\textbf{Min Cut}}: The minimum capacity of a cut in a network, equal to the max flow. %
\item \href{https://www.math.purdue.edu/~yipn/421/2024-11-21-EconNetSimplex.pdf#page=4}{\textbf{Max{-}Flow Min{-}Cut Theorem}}: The theorem stating that the maximum flow in a network is equal to the minimum cut capacity. %
\item \href{https://www.math.purdue.edu/~yipn/421/2024-11-21-EconNetSimplex.pdf#page=4}{\textbf{Simplex Method}}: An iterative algorithm for solving linear programs. %
\item \href{https://www.math.purdue.edu/~yipn/421/2024-11-21-EconNetSimplex.pdf#page=7}{\textbf{Primal Simplex Method}}: The standard simplex method applied to the primal problem. %
\item \href{https://www.math.purdue.edu/~yipn/421/2024-11-21-EconNetSimplex.pdf#page=7}{\textbf{Dual Simplex Method}}: The simplex method applied to the dual problem. %
\item \href{https://www.math.purdue.edu/~yipn/421/2024-11-21-EconNetSimplex.pdf#page=7}{\textbf{Primal Network Simplex Method}}: The simplex method specialized for network flow problems, applied to the primal problem. %
\item \href{https://www.math.purdue.edu/~yipn/421/2024-11-21-EconNetSimplex.pdf#page=7}{\textbf{Dual Network Simplex Method}}: The simplex method specialized for network flow problems, applied to the dual problem. %
\item \href{https://www.math.purdue.edu/~yipn/421/2024-11-21-EconNetSimplex.pdf#page=2}{\textbf{Slack Variables}}: Variables added to inequality constraints to transform them into equalities. %
\item \href{https://www.math.purdue.edu/~yipn/421/2024-11-21-EconNetSimplex.pdf#page=2}{\textbf{Dictionary}}: A representation of the simplex tableau. %
\item \href{https://www.math.purdue.edu/~yipn/421/2024-11-21-EconNetSimplex.pdf#page=7}{\textbf{Pivot Operation}}: A step in the simplex method that changes the basis. %
\item \href{https://www.math.purdue.edu/~yipn/421/2024-11-21-EconNetSimplex.pdf#page=8}{\textbf{Parametric Analysis}}: The study of how changes in the parameters of a linear program affect the optimal solution. %
\item \href{https://www.math.purdue.edu/~yipn/421/2024-11-21-EconNetSimplex.pdf#page=4}{\textbf{Graphical Solution Method}}: A method for solving linear programs graphically. %
\item \href{https://www.math.purdue.edu/~yipn/421/2024-11-21-EconNetSimplex.pdf#page=4}{\textbf{Integer Programming}}: Linear programming where some or all variables are restricted to integer values. %
\item \href{https://www.math.purdue.edu/~yipn/421/2024-11-21-EconNetSimplex.pdf#page=4}{\textbf{Branch and Bound}}: A method for solving integer programming problems. %
\item \href{https://www.math.purdue.edu/~yipn/421/2024-11-21-EconNetSimplex.pdf#page=4}{\textbf{Gomory Cuts}}: Cuts added to the feasible region of an integer program to improve the solution. %
\item \href{https://www.math.purdue.edu/~yipn/421/2024-11-21-EconNetSimplex.pdf#page=4}{\textbf{Farkas' Lemma}}: A theorem that provides conditions for the existence of a solution to a system of linear inequalities. %
\item \href{https://www.math.purdue.edu/~yipn/421/2024-11-21-EconNetSimplex.pdf#page=4}{\textbf{Separation Theorem}}: A theorem that states that two disjoint convex sets can be separated by a hyperplane. %
\item \href{https://www.math.purdue.edu/~yipn/421/2024-11-21-EconNetSimplex.pdf#page=4}{\textbf{Convex Set}}: A set where any line segment connecting two points in the set is also contained in the set. %
\item \href{https://www.math.purdue.edu/~yipn/421/2024-11-21-EconNetSimplex.pdf#page=4}{\textbf{Convex Hull}}: The smallest convex set containing a given set of points. %
\item \href{https://www.math.purdue.edu/~yipn/421/2024-11-21-EconNetSimplex.pdf#page=4}{\textbf{Caratheodory's Theorem}}: A theorem stating that any point in the convex hull of a set of points can be expressed as a convex combination of at most d+1 points, where d is the dimension of the space. %
\item \href{https://www.math.purdue.edu/~yipn/421/2024-11-21-EconNetSimplex.pdf#page=4}{\textbf{Hyperplane}}: A flat subspace of dimension one less than the dimension of the ambient space. %
\item \href{https://www.math.purdue.edu/~yipn/421/2024-11-21-EconNetSimplex.pdf#page=4}{\textbf{Support Vector Machine (SVM)}}: A machine learning algorithm that uses hyperplanes to perform binary classification. %
\item \href{https://www.math.purdue.edu/~yipn/421/2024-11-21-EconNetSimplex.pdf#page=4}{\textbf{Linear Regression}}: A statistical method for modeling the relationship between a dependent variable and one or more independent variables. %
\item \href{https://www.math.purdue.edu/~yipn/421/2024-11-21-EconNetSimplex.pdf#page=4}{\textbf{Smallest Enclosing Ball}}: The smallest ball that contains a given set of points. %
\item \href{https://www.math.purdue.edu/~yipn/421/2024-11-21-EconNetSimplex.pdf#page=4}{\textbf{Shortest Path Problem}}: The problem of finding the shortest path between two nodes in a network. %
\item \href{https://www.math.purdue.edu/~yipn/421/2024-11-21-EconNetSimplex.pdf#page=4}{\textbf{Lagrangian Function}}: A function used in optimization problems with constraints. %
\item \href{https://www.math.purdue.edu/~yipn/421/2024-11-21-EconNetSimplex.pdf#page=4}{\textbf{Dijkstra's Algorithm}}: An algorithm for finding the shortest paths from a single source node to all other nodes in a graph. %
\item \href{https://www.math.purdue.edu/~yipn/421/2024-11-21-EconNetSimplex.pdf#page=4}{\textbf{Bellman's Equation}}: A dynamic programming equation used to solve the shortest path problem. %
\item \href{https://www.math.purdue.edu/~yipn/421/2024-11-21-EconNetSimplex.pdf#page=4}{\textbf{Parameter λ}}: A parameter used in various optimization problems. %
\item \href{https://www.math.purdue.edu/~yipn/421/2024-11-21-EconNetSimplex.pdf#page=4}{\textbf{Sparse Solutions of Linear Systems}}: Solutions to linear systems where many of the variables are zero. %
\item \href{https://www.math.purdue.edu/~yipn/421/2024-12-03-Duality.pdf#page=1}{\textbf{Duality}}: In linear programming, duality refers to the relationship between a primal problem and its corresponding dual problem. The dual problem provides a lower bound (for maximization problems) or an upper bound (for minimization problems) on the optimal value of the primal problem. %
\item \href{https://www.math.purdue.edu/~yipn/421/2024-12-03-Duality.pdf#page=2}{\textbf{Complementary Slackness}}: At the optimal solution of a linear program and its dual, complementary slackness states that for each constraint, either the constraint is binding (the slack variable is zero) or the corresponding dual variable is zero. %
\item \href{https://www.math.purdue.edu/~yipn/421/2024-12-03-Duality.pdf#page=3}{\textbf{Graphical Method}}: A method for solving linear programming problems with two variables by plotting the constraints as lines on a graph and identifying the feasible region. The optimal solution is found at a vertex of the feasible region. %
\item \href{https://www.math.purdue.edu/~yipn/421/2024-12-03-Duality.pdf#page=11}{\textbf{Sensitivity Analysis}}: The study of how changes in the objective function coefficients or the constraint values of a linear program affect the optimal solution. %
\item \href{https://www.math.purdue.edu/~yipn/421/2024-12-03-Duality.pdf#page=15}{\textbf{Shadow Price}}: In sensitivity analysis, the shadow price of a constraint represents the rate of change of the optimal objective function value with respect to a change in the right-hand side value of that constraint. It represents the marginal value of a resource. %
\item \href{https://www.math.purdue.edu/~yipn/421/2024-12-03-Duality.pdf#page=2}{\textbf{Optimal Solution}}: The solution to a linear programming problem that maximizes (or minimizes) the objective function while satisfying all constraints. %
\item \href{https://www.math.purdue.edu/~yipn/421/2024-12-03-Duality.pdf#page=1}{\textbf{Primal Problem (P)}}: The original linear programming problem being considered. %
\item \href{https://www.math.purdue.edu/~yipn/421/2024-12-03-Duality.pdf#page=1}{\textbf{Dual Problem (D)}}: A linear programming problem derived from the primal problem, providing a lower bound (for maximization problems) or an upper bound (for minimization problems) on the optimal value of the primal problem. %
\item \href{https://www.math.purdue.edu/~yipn/421/2024-12-03-Duality.pdf#page=7}{\textbf{Basis Matrix (B)}}: A square matrix formed by selecting a subset of columns from the constraint matrix of a linear program. Used in the simplex method to represent the basic variables. %
\item \href{https://www.math.purdue.edu/~yipn/421/2024-12-03-Duality.pdf#page=1}{\textbf{Objective Function}}: The function being maximized or minimized in a linear programming problem. %
\item \href{https://www.math.purdue.edu/~yipn/421/2024-12-03-Duality.pdf#page=2}{\textbf{Slack Variables}}: Variables added to inequality constraints to transform them into equality constraints. %
\item \href{https://www.math.purdue.edu/~yipn/421/2024-12-03-Duality.pdf#page=2}{\textbf{Non{-}basic Variables}}: Variables in a linear program that are not in the basis. Their values are zero in a basic feasible solution. %
\item \href{https://www.math.purdue.edu/~yipn/421/2024-12-03-Duality.pdf#page=2}{\textbf{Basic Variables}}: Variables in a linear program that are in the basis. Their values are determined by the basis matrix and the right-hand side values. %
\item \href{https://www.math.purdue.edu/~yipn/421/2024-12-03-Duality.pdf#page=3}{\textbf{Feasible Region}}: The set of all points that satisfy the constraints of a linear programming problem. %
\item \href{https://www.math.purdue.edu/~yipn/421/2024-12-03-Duality.pdf#page=1}{\textbf{Constraints}}: The limitations or restrictions on the values of the decision variables in a linear programming problem. %
\item \href{https://www.math.purdue.edu/~yipn/421/2024-12-03-Duality.pdf#page=1}{\textbf{Decision Variables}}: The variables whose values are being optimized in a linear programming problem. %
\item \href{https://www.math.purdue.edu/~yipn/421/2024-12-03-Duality.pdf#page=4}{\textbf{Parametric Analysis}}: A technique used in sensitivity analysis to study the effect of changes in a parameter (e.g., objective function coefficient or constraint value) on the optimal solution. %
\item \href{https://www.math.purdue.edu/~yipn/421/2024-12-03-Duality.pdf#page=5}{\textbf{Negative Transpose Property}}: The relationship between the coefficient matrix and the right-hand side vector of the primal and dual problems. The coefficient matrix of the dual is the transpose of the coefficient matrix of the primal, and the right-hand side vector of the dual is the objective function vector of the primal. %
\item \href{https://www.math.purdue.edu/~yipn/421/2024-12-03-Duality.pdf#page=9}{\textbf{Sparse Solutions}}: Solutions to linear systems where many of the variables have a value of zero. %
\item \href{https://www.math.purdue.edu/~yipn/421/2024-12-03-Duality.pdf#page=2}{\textbf{Complementary Variables}}: In the simplex method, complementary variables are pairs of variables such that at most one variable in each pair can be non-zero at optimality. This is a consequence of complementary slackness. %
\item \href{https://www.math.purdue.edu/~yipn/421/2024-12-03-Duality.pdf#page=1}{\textbf{Strong Duality Theorem}}: If both the primal and dual problems have feasible solutions, then they both have optimal solutions, and the optimal objective function values are equal. %
\item \href{https://www.math.purdue.edu/~yipn/421/2024-12-03-Duality.pdf#page=1}{\textbf{Weak Duality Theorem}}: The objective function value of the dual problem provides a lower bound for the objective function value of the primal problem (for maximization problems). %
\item \href{https://www.math.purdue.edu/~yipn/421/2024-12-03-Duality.pdf#page=1}{\textbf{Farkas' Lemma}}: A fundamental theorem in linear programming that provides conditions for the existence of a non-negative solution to a system of linear inequalities. %
\item \href{https://www.math.purdue.edu/~yipn/421/2024-12-03-Duality.pdf#page=1}{\textbf{Caratheodory's Theorem}}: A theorem stating that any point in the convex hull of a set of points in R<sup>n</sup> can be expressed as a convex combination of at most n+1 points from the set. %
\item \href{https://www.math.purdue.edu/~yipn/421/2024-12-03-Duality.pdf#page=1}{\textbf{Separation Theorem}}: A theorem stating that two disjoint convex sets can be separated by a hyperplane. %
\item \href{https://www.math.purdue.edu/~yipn/421/2024-12-03-Duality.pdf#page=1}{\textbf{Convex Set}}: A set such that for any two points in the set, the line segment connecting them is also contained in the set. %
\item \href{https://www.math.purdue.edu/~yipn/421/2024-12-03-Duality.pdf#page=1}{\textbf{Convex Combination}}: A linear combination of points where the coefficients are non-negative and sum to one. %
\item \href{https://www.math.purdue.edu/~yipn/421/2024-12-03-Duality.pdf#page=1}{\textbf{Convex Hull}}: The smallest convex set containing a given set of points. %
\item \href{https://www.math.purdue.edu/~yipn/421/2024-12-03-Duality.pdf#page=2}{\textbf{Dictionary}}: A representation of a linear program in a form suitable for the simplex method. %
\item \href{https://www.math.purdue.edu/~yipn/421/2024-12-03-Duality.pdf#page=1}{\textbf{Pivot Operation}}: A step in the simplex method that involves selecting an entering and leaving variable and updating the dictionary. %
\item \href{https://www.math.purdue.edu/~yipn/421/2024-12-03-Duality.pdf#page=1}{\textbf{Simplex Method}}: An iterative algorithm for solving linear programming problems. %
\item \href{https://www.math.purdue.edu/~yipn/421/2024-12-03-Duality.pdf#page=1}{\textbf{Simplex Tableau}}: A tabular representation of the simplex method's iterations. %
\item \href{https://www.math.purdue.edu/~yipn/421/2024-12-03-Duality.pdf#page=1}{\textbf{Network Simplex Method}}: A specialized version of the simplex method for solving network flow problems. %
\item \href{https://www.math.purdue.edu/~yipn/421/2024-12-03-Duality.pdf#page=1}{\textbf{Dual Network Simplex Method}}: A specialized version of the dual simplex method for solving network flow problems. %
\item \href{https://www.math.purdue.edu/~yipn/421/2024-12-03-Duality.pdf#page=1}{\textbf{Ford{-}Fulkerson Algorithm}}: An algorithm for finding the maximum flow in a network. %
\item \href{https://www.math.purdue.edu/~yipn/421/2024-12-03-Duality.pdf#page=1}{\textbf{Augmenting Path}}: A path in a network that can be used to increase the flow. %
\item \href{https://www.math.purdue.edu/~yipn/421/2024-12-03-Duality.pdf#page=1}{\textbf{Max{-}Flow Min{-}Cut Theorem}}: A theorem stating that the maximum flow in a network is equal to the minimum capacity of a cut. %
\item \href{https://www.math.purdue.edu/~yipn/421/2024-12-03-Duality.pdf#page=1}{\textbf{Min Cut}}: A partition of the nodes in a network into two sets, such that the capacity of the edges crossing the partition is minimized. %
\item \href{https://www.math.purdue.edu/~yipn/421/2024-12-03-Duality.pdf#page=1}{\textbf{Shortest Path Problem}}: The problem of finding the shortest path between two nodes in a network. %
\item \href{https://www.math.purdue.edu/~yipn/421/2024-12-03-Duality.pdf#page=1}{\textbf{Dijkstra's Algorithm}}: An algorithm for finding the shortest path between two nodes in a network. %
\item \href{https://www.math.purdue.edu/~yipn/421/2024-12-03-Duality.pdf#page=1}{\textbf{Bellman's Equation}}: A dynamic programming equation for solving the shortest path problem. %
\item \href{https://www.math.purdue.edu/~yipn/421/2024-12-03-Duality.pdf#page=1}{\textbf{Transportation Problem}}: A linear programming problem involving the transportation of goods from sources to destinations. %
\item \href{https://www.math.purdue.edu/~yipn/421/2024-12-03-Duality.pdf#page=1}{\textbf{Assignment Problem}}: A special case of the transportation problem where each source has a supply of one unit and each destination has a demand of one unit. %
\item \href{https://www.math.purdue.edu/~yipn/421/2024-12-03-Duality.pdf#page=1}{\textbf{Cutting Stock Problem}}: A linear programming problem involving the cutting of stock material into smaller pieces to meet demands. %
\item \href{https://www.math.purdue.edu/~yipn/421/2024-12-03-Duality.pdf#page=1}{\textbf{Knapsack Problem}}: A combinatorial optimization problem where a set of items with weights and values must be selected to maximize the total value while not exceeding a weight limit. %
\item \href{https://www.math.purdue.edu/~yipn/421/2024-12-03-Duality.pdf#page=1}{\textbf{Machine Scheduling}}: A combinatorial optimization problem involving the scheduling of jobs on machines to minimize the makespan or other objective. %
\item \href{https://www.math.purdue.edu/~yipn/421/2024-12-03-Duality.pdf#page=1}{\textbf{Workforce Planning}}: A linear programming problem involving the determination of the optimal number of workers to hire to meet demands. %
\item \href{https://www.math.purdue.edu/~yipn/421/2024-12-03-Duality.pdf#page=1}{\textbf{Portfolio Selection (Markowitz Model)}}: A linear programming problem involving the selection of investments to maximize return while minimizing risk. %
\item \href{https://www.math.purdue.edu/~yipn/421/2024-12-03-Duality.pdf#page=1}{\textbf{Ice Cream Production Planning}}: A linear programming problem involving the determination of the optimal production plan for ice cream to meet demands. %
\item \href{https://www.math.purdue.edu/~yipn/421/2024-12-03-Duality.pdf#page=1}{\textbf{Diet Problem}}: A linear programming problem involving the determination of the optimal diet to meet nutritional requirements while minimizing cost. %
\item \href{https://www.math.purdue.edu/~yipn/421/2024-12-03-Duality.pdf#page=1}{\textbf{Support Vector Machine (SVM)}}: A machine learning algorithm for binary classification that finds the optimal hyperplane to separate data points. %
\item \href{https://www.math.purdue.edu/~yipn/421/2024-12-03-Duality.pdf#page=1}{\textbf{Hyperplane}}: A linear subspace of dimension one less than the dimension of the ambient space. %
\item \href{https://www.math.purdue.edu/~yipn/421/2024-12-03-Duality.pdf#page=1}{\textbf{Linear Separability}}: The ability to separate data points into two classes using a hyperplane. %
\item \href{https://www.math.purdue.edu/~yipn/421/2024-12-03-Duality.pdf#page=1}{\textbf{Support Vectors}}: The data points closest to the hyperplane in an SVM. %
\item \href{https://www.math.purdue.edu/~yipn/421/2024-12-03-Duality.pdf#page=1}{\textbf{Maximum Margin Classifier}}: An SVM that maximizes the margin between the hyperplane and the data points. %
\item \href{https://www.math.purdue.edu/~yipn/421/2024-12-03-Duality.pdf#page=1}{\textbf{Lagrangian Function}}: A function used in optimization problems with constraints. %
\item \href{https://www.math.purdue.edu/~yipn/421/2024-12-03-Duality.pdf#page=1}{\textbf{Lagrange Multiplier}}: A parameter in the Lagrangian function that represents the penalty for violating a constraint. %
\item \href{https://www.math.purdue.edu/~yipn/421/2024-12-03-Duality.pdf#page=1}{\textbf{Quadratic Programming (QP)}}: A type of optimization problem where the objective function is quadratic and the constraints are linear. %
\item \href{https://www.math.purdue.edu/~yipn/421/2024-12-03-Duality.pdf#page=1}{\textbf{Gomory Cuts}}: Cutting planes used in integer programming to improve the linear programming relaxation. %
\item \href{https://www.math.purdue.edu/~yipn/421/2024-12-03-Duality.pdf#page=1}{\textbf{Integer Programming}}: A type of optimization problem where the decision variables are restricted to integer values. %
\item \href{https://www.math.purdue.edu/~yipn/421/2024-12-03-Duality.pdf#page=1}{\textbf{Relaxed Problem}}: A linear programming problem obtained by relaxing the integer constraints of an integer programming problem. %
\item \href{https://www.math.purdue.edu/~yipn/421/2024-12-03-Duality.pdf#page=1}{\textbf{Branch and Bound}}: An algorithm for solving integer programming problems. %
\item \href{https://www.math.purdue.edu/~yipn/421/2024-12-03-Duality.pdf#page=1}{\textbf{Spanning Tree}}: A tree that connects all nodes in a graph. %
\item \href{https://www.math.purdue.edu/~yipn/421/2024-12-03-Duality.pdf#page=1}{\textbf{Incidence Matrix}}: A matrix representing the connections between nodes and arcs in a network. %
\item \href{https://www.math.purdue.edu/~yipn/421/2024-12-03-Duality.pdf#page=1}{\textbf{Network}}: A graph representing connections between nodes. %
\item \href{https://www.math.purdue.edu/~yipn/421/2024-12-03-Duality.pdf#page=1}{\textbf{Nodes}}: Points in a network. %
\item \href{https://www.math.purdue.edu/~yipn/421/2024-12-03-Duality.pdf#page=1}{\textbf{Arcs}}: Connections between nodes in a network. %
\item \href{https://www.math.purdue.edu/~yipn/421/2024-12-03-Duality.pdf#page=1}{\textbf{Flow}}: The amount of material flowing through an arc in a network. %
\item \href{https://www.math.purdue.edu/~yipn/421/2024-12-03-Duality.pdf#page=1}{\textbf{Path}}: A sequence of arcs connecting two nodes in a network. %
\item \href{https://www.math.purdue.edu/~yipn/421/2024-12-03-Duality.pdf#page=1}{\textbf{Supply/Demand (bᵢ)}}: The supply or demand at a node in a network. %
\item \href{https://www.math.purdue.edu/~yipn/421/2024-12-03-Duality.pdf#page=1}{\textbf{Auxiliary Problem}}: A problem used to find an initial feasible solution for a linear program. %
\item \href{https://www.math.purdue.edu/~yipn/421/2024-12-03-Duality.pdf#page=1}{\textbf{Entering Arc}}: An arc selected to enter the basis in the network simplex method. %
\item \href{https://www.math.purdue.edu/~yipn/421/2024-12-03-Duality.pdf#page=1}{\textbf{Leaving Arc}}: An arc selected to leave the basis in the network simplex method. %
\item \href{https://www.math.purdue.edu/~yipn/421/2024-12-03-Duality.pdf#page=1}{\textbf{Cycle}}: A closed path in a network. %
\item \href{https://www.math.purdue.edu/~yipn/421/2024-12-03-Duality.pdf#page=1}{\textbf{Tree}}: A connected graph with no cycles. %
\item \href{https://www.math.purdue.edu/~yipn/421/2024-12-03-Duality.pdf#page=1}{\textbf{Tree Solution}}: A solution to a network flow problem that corresponds to a spanning tree. %
\item \href{https://www.math.purdue.edu/~yipn/421/2024-12-03-Duality.pdf#page=1}{\textbf{Rank(A)}}: The number of linearly independent rows or columns in a matrix A. %
\item \href{https://www.math.purdue.edu/~yipn/421/2024-12-03-Duality.pdf#page=1}{\textbf{Reduced Costs}}: The costs of non-basic variables in the simplex method. %
\item \href{https://www.math.purdue.edu/~yipn/421/2024-12-03-Duality.pdf#page=1}{\textbf{Parameter λ}}: A parameter used in various optimization problems. %
\item \href{https://www.math.purdue.edu/~yipn/421/2024-12-03-Duality.pdf#page=1}{\textbf{Parameter δbᵢ}}: A parameter representing a change in the right-hand side value of a constraint. %
\item \href{https://www.math.purdue.edu/~yipn/421/2024-12-03-Duality.pdf#page=1}{\textbf{Fluctuation (r̃ⱼ)}}: A measure of the variability of a variable. %
\item \href{https://www.math.purdue.edu/~yipn/421/2024-12-03-Duality.pdf#page=1}{\textbf{Monthly Demands}}: The demands for a product in each month. %
\item \href{https://www.math.purdue.edu/~yipn/421/2024-12-03-Duality.pdf#page=1}{\textbf{Production at Month i (xᵢ)}}: The production level in month i. %
\item \href{https://www.math.purdue.edu/~yipn/421/2024-12-03-Duality.pdf#page=1}{\textbf{Cost of Storage}}: The cost of storing a product. %
\item \href{https://www.math.purdue.edu/~yipn/421/2024-12-03-Duality.pdf#page=1}{\textbf{Cost of Changing Production}}: The cost of changing the production level. %
\item \href{https://www.math.purdue.edu/~yipn/421/2024-12-03-Duality.pdf#page=1}{\textbf{Average Demand}}: The average demand for a product over a period of time. %
\item \href{https://www.math.purdue.edu/~yipn/421/2024-12-03-Duality.pdf#page=1}{\textbf{Proportion of Investment in i (xᵢ)}}: The proportion of investment in investment i. %
\item \href{https://www.math.purdue.edu/~yipn/421/2024-12-03-Duality.pdf#page=1}{\textbf{Annual Return of ith Investment (rᵢ)}}: The annual return of investment i. %
\item \href{https://www.math.purdue.edu/~yipn/421/2024-12-03-Duality.pdf#page=1}{\textbf{Surplus at End of Month i (sᵢ)}}: The surplus of a product at the end of month i. %
\item \href{https://www.math.purdue.edu/~yipn/421/2024-12-03-Duality.pdf#page=1}{\textbf{L¹ Regression}}: A type of regression that minimizes the sum of absolute errors. %
\item \href{https://www.math.purdue.edu/~yipn/421/2024-12-03-Duality.pdf#page=1}{\textbf{L² Regression (Least Square)}}: A type of regression that minimizes the sum of squared errors. %
\item \href{https://www.math.purdue.edu/~yipn/421/2024-12-03-Duality.pdf#page=1}{\textbf{L∞ Regression}}: A type of regression that minimizes the maximum absolute error. %
\item \href{https://www.math.purdue.edu/~yipn/421/2024-12-03-Duality.pdf#page=1}{\textbf{Smallest Enclosing Ball}}: The smallest ball that contains a given set of points. %
\item \href{https://www.math.purdue.edu/~yipn/421/2024-12-03-Duality.pdf#page=1}{\textbf{Smallest Ball Problem}}: The problem of finding the smallest ball that contains a given set of points. %
\item \href{https://www.math.purdue.edu/~yipn/421/2024-12-03-Duality.pdf#page=1}{\textbf{Maximum Inscribed Circle in a Convex Polyhedron}}: The largest circle that can be inscribed in a convex polyhedron. %
\item \href{https://www.math.purdue.edu/~yipn/421/2024-12-03-Duality.pdf#page=1}{\textbf{Median}}: The middle value in a sorted list of numbers. %
\item \href{https://www.math.purdue.edu/~yipn/421/2024-12-03-Duality.pdf#page=1}{\textbf{Mean}}: The average of a set of numbers. %
\item \href{https://www.math.purdue.edu/~yipn/421/2024-12-03-Duality.pdf#page=1}{\textbf{Mid{-}Range}}: The average of the maximum and minimum values in a set of numbers. %
\item \href{https://www.math.purdue.edu/~yipn/421/2024-12-03-Duality.pdf#page=1}{\textbf{Klee{-}Minty Example}}: An example of a linear program that demonstrates the worst-case behavior of the simplex method. %
\item \href{https://www.math.purdue.edu/~yipn/421/2024-12-03-Duality.pdf#page=1}{\textbf{Bland's Rule}}: A pivot rule for the simplex method that prevents cycling. %
\item \href{https://www.math.purdue.edu/~yipn/421/2024-12-03-Duality.pdf#page=1}{\textbf{Largest{-}Coefficient Rule}}: A pivot rule for the simplex method that selects the variable with the largest coefficient in the objective function. %
\item \href{https://www.math.purdue.edu/~yipn/421/2024-12-03-Duality.pdf#page=1}{\textbf{Largest{-}Increase Rule}}: A pivot rule for the simplex method that selects the variable that leads to the largest increase in the objective function value. %
\item \href{https://www.math.purdue.edu/~yipn/421/2024-12-03-Duality.pdf#page=1}{\textbf{Lexicographic Method}}: A method for preventing cycling in the simplex method. %
\item \href{https://www.math.purdue.edu/~yipn/421/2024-12-03-Duality.pdf#page=1}{\textbf{Randomized Pivot Rule}}: A pivot rule for the simplex method that randomly selects the entering variable. %
\item \href{https://www.math.purdue.edu/~yipn/421/2024-12-03-Duality.pdf#page=1}{\textbf{Empirical vs. Theoretical Analysis of Simplex Method}}: A comparison of the observed performance of the simplex method with its theoretical worst-case behavior. %
\item \href{https://www.math.purdue.edu/~yipn/421/2024-12-03-Duality.pdf#page=1}{\textbf{Smoothed Complexity}}: A measure of the complexity of an algorithm that takes into account the distribution of the input data. %
\item \href{https://www.math.purdue.edu/~yipn/421/2024-12-03-Duality.pdf#page=1}{\textbf{Impact of Pivot Rules on Simplex Method Efficiency}}: The effect of different pivot rules on the efficiency of the simplex method. %
\item \href{https://www.math.purdue.edu/~yipn/421/2024-12-03-Duality.pdf#page=1}{\textbf{Understanding Worst{-}Case Behavior of the Simplex Method}}: An analysis of the worst-case performance of the simplex method. %
\item \href{https://www.math.purdue.edu/~yipn/421/2024-12-03-Duality.pdf#page=1}{\textbf{Analyzing the Efficiency of the Simplex Method}}: An analysis of the efficiency of the simplex method. %
\item \href{https://www.math.purdue.edu/~yipn/421/2024-12-03-Duality.pdf#page=1}{\textbf{Illustrating Simplex Iterations with a Numerical Example}}: A demonstration of the simplex method using a numerical example. %
\item \href{https://www.math.purdue.edu/~yipn/421/2024-12-03-Duality.pdf#page=1}{\textbf{Iterative Improvement Using the Simplex Method}}: A description of how the simplex method iteratively improves the solution. %
\item \href{https://www.math.purdue.edu/~yipn/421/2024-12-03-Duality.pdf#page=1}{\textbf{Updating the Simplex Tableau Using Matrix Operations}}: A description of how matrix operations are used to update the simplex tableau. %
\item \href{https://www.math.purdue.edu/~yipn/421/2024-12-03-Duality.pdf#page=1}{\textbf{Expressing the Objective Function in Terms of Basic and Non{-}Basic Variables}}: A description of how the objective function can be expressed in terms of basic and non-basic variables. %
\item \href{https://www.math.purdue.edu/~yipn/421/2024-12-03-Duality.pdf#page=1}{\textbf{Analyzing the Optimal Tableau and Its Properties}}: An analysis of the properties of the optimal simplex tableau. %
\item \href{https://www.math.purdue.edu/~yipn/421/2024-12-03-Duality.pdf#page=1}{\textbf{Understanding Complementary Variables and Their Relationship in Simplex Iterations}}: An explanation of the relationship between complementary variables in simplex iterations. %
\item \href{https://www.math.purdue.edu/~yipn/421/2024-12-03-Duality.pdf#page=1}{\textbf{Understanding Spanning Tree Properties}}: An explanation of the properties of spanning trees. %
\item \href{https://www.math.purdue.edu/~yipn/421/2024-12-03-Duality.pdf#page=1}{\textbf{Adjusting Shipments to Maintain Feasibility}}: A description of how shipments can be adjusted to maintain feasibility in a transportation problem. %
\item \href{https://www.math.purdue.edu/~yipn/421/2024-12-03-Duality.pdf#page=1}{\textbf{Handling Infeasible Origins in Dual Based Phase I Algorithm}}: A description of how infeasible origins are handled in the dual based phase I algorithm. %
\item \href{https://www.math.purdue.edu/~yipn/421/2024-12-03-Duality.pdf#page=1}{\textbf{Handling Infeasible Origins}}: A description of how infeasible origins are handled in a network flow problem. %
\item \href{https://www.math.purdue.edu/~yipn/421/2024-12-03-Duality.pdf#page=1}{\textbf{Finding a Feasible Solution on a Spanning Tree}}: A description of how a feasible solution can be found on a spanning tree. %
\item \href{https://www.math.purdue.edu/~yipn/421/2024-12-03-Duality.pdf#page=1}{\textbf{Finding Primal Flows Using a Spanning Tree}}: A description of how primal flows can be found using a spanning tree. %
\item \href{https://www.math.purdue.edu/~yipn/421/2024-12-03-Duality.pdf#page=1}{\textbf{Finding Dual Variables Using a Spanning Tree}}: A description of how dual variables can be found using a spanning tree. %
\item \href{https://www.math.purdue.edu/~yipn/421/2024-12-03-Duality.pdf#page=1}{\textbf{Finding Primal and Dual Solutions Using Spanning Trees}}: A description of how primal and dual solutions can be found using spanning trees. %
\item \href{https://www.math.purdue.edu/~yipn/421/2024-12-03-Duality.pdf#page=1}{\textbf{Finding a Dual Feasible Solution}}: A description of how a dual feasible solution can be found. %
\item \href{https://www.math.purdue.edu/~yipn/421/2024-12-03-Duality.pdf#page=1}{\textbf{Finding an Initial Feasible Solution}}: A description of how an initial feasible solution can be found for a linear program. %
\item \href{https://www.math.purdue.edu/~yipn/421/2024-12-03-Duality.pdf#page=1}{\textbf{Identifying Profitable Opportunities for Competitors}}: A description of how profitable opportunities for competitors can be identified. %
\item \href{https://www.math.purdue.edu/~yipn/421/2024-12-03-Duality.pdf#page=1}{\textbf{Solving the Upper Bounded Transshipment Problem}}: A description of how the upper bounded transshipment problem can be solved. %
\item \href{https://www.math.purdue.edu/~yipn/421/2024-12-03-Duality.pdf#page=1}{\textbf{Solving Integer Programming Problems by Relaxing and Rounding}}: A description of how integer programming problems can be solved by relaxing and rounding. %
\item \href{https://www.math.purdue.edu/~yipn/421/2024-12-03-Duality.pdf#page=1}{\textbf{Solving Linear Programs Using the Basis and Non{-}Basis Method}}: A description of how linear programs can be solved using the basis and non-basis method. %
\item \href{https://www.math.purdue.edu/~yipn/421/2024-12-03-Duality.pdf#page=1}{\textbf{Solving a Specific Integer Programming Problem}}: A description of how a specific integer programming problem can be solved. %
\item \href{https://www.math.purdue.edu/~yipn/421/2024-12-03-Duality.pdf#page=1}{\textbf{Solving Linear Programs Graphically}}: A description of how linear programs can be solved graphically. %
\item \href{https://www.math.purdue.edu/~yipn/421/2024-12-03-Duality.pdf#page=1}{\textbf{Demonstrating the Preservation of the Negative Transpose Property}}: A demonstration of how the negative transpose property is preserved in the primal and dual problems. %
\item \href{https://www.math.purdue.edu/~yipn/421/2024-12-03-Duality.pdf#page=1}{\textbf{Deriving the Dual Problem from the Lagrangian}}: A description of how the dual problem can be derived from the Lagrangian. %
\item \href{https://www.math.purdue.edu/~yipn/421/2024-12-03-Duality.pdf#page=1}{\textbf{Deriving the Dual of a General LP}}: A description of how the dual of a general linear program can be derived. %
\item \href{https://www.math.purdue.edu/~yipn/421/2024-12-03-Duality.pdf#page=1}{\textbf{Deriving the Dual of a General LP with Upper Bounds}}: A description of how the dual of a general linear program with upper bounds can be derived. %
\item \href{https://www.math.purdue.edu/~yipn/421/2024-12-03-Duality.pdf#page=1}{\textbf{Formulating Linear Programs in Matrix Notation}}: A description of how linear programs can be formulated in matrix notation. %
\item \href{https://www.math.purdue.edu/~yipn/421/2024-12-03-Duality.pdf#page=1}{\textbf{Formulating Primal and Dual Linear Programs}}: A description of how primal and dual linear programs can be formulated. %
\item \href{https://www.math.purdue.edu/~yipn/421/2024-12-03-Duality.pdf#page=1}{\textbf{Formulating Transportation Problems with Inequality Constraints (Algebraic Representation)}}: A description of how transportation problems with inequality constraints can be formulated algebraically. %
\item \href{https://www.math.purdue.edu/~yipn/421/2024-12-03-Duality.pdf#page=1}{\textbf{Formulating Transportation Problems with Inequality Constraints (Graphical Representation)}}: A description of how transportation problems with inequality constraints can be formulated graphically. %
\item \href{https://www.math.purdue.edu/~yipn/421/2024-12-03-Duality.pdf#page=1}{\textbf{Formulating the Maximum Flow Problem as a Linear Program}}: A description of how the maximum flow problem can be formulated as a linear program. %
\item \href{https://www.math.purdue.edu/~yipn/421/2024-12-03-Duality.pdf#page=1}{\textbf{Formulating Network Flow Problems as Linear Programs}}: A description of how network flow problems can be formulated as linear programs. %
\item \href{https://www.math.purdue.edu/~yipn/421/2024-12-03-Duality.pdf#page=1}{\textbf{Reformulating a General LP Without Non{-}Negativity Constraints}}: A description of how a general linear program without non-negativity constraints can be reformulated. %
\item \href{https://www.math.purdue.edu/~yipn/421/2024-12-03-Duality.pdf#page=1}{\textbf{Representing Linear Programs in Matrix Form}}: A description of how linear programs can be represented in matrix form. %
\item \href{https://www.math.purdue.edu/~yipn/421/2024-12-03-Duality.pdf#page=1}{\textbf{Applying the Branch and Bound Method}}: A description of how the branch and bound method can be applied to solve integer programming problems. %
\item \href{https://www.math.purdue.edu/~yipn/421/2024-12-03-Duality.pdf#page=1}{\textbf{Applying the Caratheodory Theorem}}: A description of how the Caratheodory theorem can be applied. %
\item \href{https://www.math.purdue.edu/~yipn/421/2024-12-03-Duality.pdf#page=1}{\textbf{Applying Complementary Slackness}}: A description of how complementary slackness can be applied. %
\item \href{https://www.math.purdue.edu/~yipn/421/2024-12-03-Duality.pdf#page=1}{\textbf{Applying Gomory Cuts to Solve Integer Programming Problems}}: A description of how Gomory cuts can be applied to solve integer programming problems. %
\item \href{https://www.math.purdue.edu/~yipn/421/2024-12-03-Duality.pdf#page=1}{\textbf{Applying the Primal Simplex Method to Network Flows}}: A description of how the primal simplex method can be applied to network flows. %
\item \href{https://www.math.purdue.edu/~yipn/421/2024-12-03-Duality.pdf#page=1}{\textbf{Applying Weak Duality Theorem}}: A description of how the weak duality theorem can be applied. %
\item \href{https://www.math.purdue.edu/~yipn/421/2024-12-03-Duality.pdf#page=1}{\textbf{Applying the Simplex Method in Matrix Form}}: A description of how the simplex method can be applied in matrix form. %
\item \href{https://www.math.purdue.edu/~yipn/421/2024-12-03-Duality.pdf#page=1}{\textbf{Analyzing the Incidence Matrix of a Network}}: An analysis of the incidence matrix of a network. %
\item \href{https://www.math.purdue.edu/~yipn/421/2024-12-03-Duality.pdf#page=1}{\textbf{Analyzing the Relationship Between Primal and Dual Problems}}: An analysis of the relationship between primal and dual problems. %
\item \href{https://www.math.purdue.edu/~yipn/421/2024-12-03-Duality.pdf#page=1}{\textbf{Analyzing the Relationship Between Primal and Dual Objective Function Values}}: An analysis of the relationship between primal and dual objective function values. %
\item \href{https://www.math.purdue.edu/~yipn/421/2024-12-03-Duality.pdf#page=1}{\textbf{Analyzing the Sensitivity of the Objective Function Coefficients}}: An analysis of the sensitivity of the objective function coefficients. %
\item \href{https://www.math.purdue.edu/~yipn/421/2024-12-03-Duality.pdf#page=1}{\textbf{Analyzing the Sensitivity of the Constraint Values}}: An analysis of the sensitivity of the constraint values. %
\item \href{https://www.math.purdue.edu/~yipn/421/2024-12-03-Duality.pdf#page=1}{\textbf{Proving Farkas' Lemma}}: A proof of Farkas' lemma. %
\item \href{https://www.math.purdue.edu/~yipn/421/2024-12-03-Duality.pdf#page=1}{\textbf{Proving Properties of Convex Sets}}: A proof of the properties of convex sets. %
\item \href{https://www.math.purdue.edu/~yipn/421/2024-12-03-Duality.pdf#page=1}{\textbf{Proving the Separation Theorem}}: A proof of the separation theorem. %
\item \href{https://www.math.purdue.edu/~yipn/421/2024-12-03-Duality.pdf#page=1}{\textbf{Proving the Strong Duality Theorem}}: A proof of the strong duality theorem. %
\item \href{https://www.math.purdue.edu/~yipn/421/2024-12-03-Duality.pdf#page=1}{\textbf{Proving Weak Duality}}: A proof of weak duality. %
\item \href{https://www.math.purdue.edu/~yipn/421/2024-12-03-Duality.pdf#page=1}{\textbf{Defining Convex Sets and Convex Combinations}}: A definition of convex sets and convex combinations. %
\item \href{https://www.math.purdue.edu/~yipn/421/2024-12-03-Duality.pdf#page=1}{\textbf{Defining and Proving Properties of Convex Hulls}}: A definition and proof of the properties of convex hulls. %
\item \href{https://www.math.purdue.edu/~yipn/421/2024-12-03-Duality.pdf#page=1}{\textbf{Determining the Existence of a Non{-}Negative Solution to Ax ≤ b}}: A determination of the existence of a non-negative solution to Ax ≤ b. %
\item \href{https://www.math.purdue.edu/~yipn/421/2024-12-03-Duality.pdf#page=1}{\textbf{Determining the Existence of a Non{-}Negative Solution to Ax = b}}: A determination of the existence of a non-negative solution to Ax = b. %
\item \href{https://www.math.purdue.edu/~yipn/421/2024-12-03-Duality.pdf#page=1}{\textbf{Determining the Four Possible Outcomes of Primal{-}Dual Problem Pairs}}: A determination of the four possible outcomes of primal-dual problem pairs. %
\item \href{https://www.math.purdue.edu/~yipn/421/2024-12-03-Duality.pdf#page=1}{\textbf{Determining the Existence of a Solution to Ax ≤ b}}: A determination of the existence of a solution to Ax ≤ b. %
\item \href{https://www.math.purdue.edu/~yipn/421/2024-12-03-Duality.pdf#page=1}{\textbf{Justifying Depth{-}First Search in Branch and Bound}}: A justification of the use of depth-first search in branch and bound. %
\item \href{https://www.math.purdue.edu/~yipn/421/2024-12-03-Duality.pdf#page=1}{\textbf{Illustrating Simplex Iterations with a Numerical Example}}: An illustration of simplex iterations with a numerical example. %
\item \href{https://www.math.purdue.edu/~yipn/421/2024-12-03-Duality.pdf#page=1}{\textbf{Graphical Representation}}: A graphical representation of a linear program. %
\item \href{https://www.math.purdue.edu/~yipn/421/2024-12-03-Duality.pdf#page=1}{\textbf{Graphical Representation and Solution}}: A graphical representation and solution of a linear program. %
\item \href{https://www.math.purdue.edu/~yipn/421/2024-12-03-Duality.pdf#page=1}{\textbf{Geometric Interpretation of Dual Constraints}}: A geometric interpretation of the dual constraints. %
\item \href{https://www.math.purdue.edu/~yipn/421/2024-12-03-Duality.pdf#page=1}{\textbf{Interpreting the Relationship Between Primal and Dual Variables}}: An interpretation of the relationship between primal and dual variables. %
\item \href{https://www.math.purdue.edu/~yipn/421/2024-12-03-Duality.pdf#page=1}{\textbf{Formulating the Dual of a General LP Problem}}: A formulation of the dual of a general linear programming problem. %
\item \href{https://www.math.purdue.edu/~yipn/421/2024-12-03-Duality.pdf#page=1}{\textbf{Formulating the Final Form of the Dual Problem with Upper Bounds}}: A formulation of the final form of the dual problem with upper bounds. %
\item \href{https://www.math.purdue.edu/~yipn/421/2024-12-03-Duality.pdf#page=1}{\textbf{Implementing the Ford{-}Fulkerson Algorithm}}: An implementation of the Ford-Fulkerson algorithm. %
\item \href{https://www.math.purdue.edu/~yipn/421/2024-12-03-Duality.pdf#page=1}{\textbf{Verifying Optimality}}: A verification of the optimality of a solution. %
\item \href{https://www.math.purdue.edu/~yipn/421/2024-12-03-Duality.pdf#page=1}{\textbf{Understanding Complementary Variables in Simplex Iterations}}: An understanding of complementary variables in simplex iterations. %
\item \href{https://www.math.purdue.edu/~yipn/421/2024-12-03-Duality.pdf#page=1}{\textbf{Understanding Spanning Tree Properties}}: An understanding of spanning tree properties. %
\end{itemize}

%
\section{Problem Types}%
\label{sec:ProblemTypes}%
\begin{itemize}[leftmargin=*]%
\item \href{https://www.math.purdue.edu/~yipn/421/2024-08-27-ExSimplex.pdf#page=32}{\textbf{Finding an Initial Feasible Solution}}: The problem of finding a starting feasible solution for the Simplex method, particularly when the origin is not feasible. An auxiliary problem is introduced to find a feasible solution, which is then used to initialize the Simplex method. %
\item \href{https://www.math.purdue.edu/~yipn/421/2024-08-27-ExSimplex.pdf#page=30}{\textbf{Handling Infeasible Origins}}: Addressing situations where the origin (0,0) is not within the feasible region defined by the constraints of a linear programming problem. This requires a different approach to find an initial feasible solution before applying the Simplex method. %
\item \href{https://www.math.purdue.edu/~yipn/421/2024-08-27-ExSimplex.pdf#page=2}{\textbf{Graphical Representation and Solution}}: Solving a linear programming problem graphically by identifying the feasible region and determining the optimal solution at a vertex of the feasible region. This method is particularly useful for visualizing the problem and understanding the Simplex method's iterative nature. %
\item \href{https://www.math.purdue.edu/~yipn/421/2024-08-27-ExSimplex.pdf#page=6}{\textbf{Iterative Improvement using the Simplex Method}}: The core process of the Simplex method, which involves iteratively moving from one vertex of the feasible region to another, improving the objective function value at each step until the optimal solution is found. This involves pivot operations and updating the dictionary of variables. %
\item \href{https://www.math.purdue.edu/~yipn/421/2024-08-27-ExSimplex.pdf#page=26}{\textbf{Working with Dictionaries of Variables}}: Utilizing a structured representation of the linear programming problem, called a dictionary, to organize basic and non-basic variables and facilitate the pivot operations within the Simplex method. Each dictionary represents a vertex of the feasible region. %
\item \href{https://www.math.purdue.edu/~yipn/421/2024-09-05-SimplexEfficiency.pdf#page=1}{\textbf{Analyzing the Efficiency of the Simplex Method}}: This problem involves investigating the average-case and worst-case performance of the Simplex method, considering factors like the number of iterations and the influence of different pivot rules. The analysis uses empirical data and theoretical bounds to assess efficiency. %
\item \href{https://www.math.purdue.edu/~yipn/421/2024-09-05-SimplexEfficiency.pdf#page=2}{\textbf{Understanding Worst{-}Case Behavior of the Simplex Method}}: This problem focuses on identifying and analyzing instances where the Simplex method exhibits poor performance, requiring an exponential number of iterations. The Klee-Minty example is used to illustrate this worst-case behavior. %
\item \href{https://www.math.purdue.edu/~yipn/421/2024-09-05-SimplexEfficiency.pdf#page=3}{\textbf{Impact of Pivot Rules on Simplex Method Efficiency}}: This problem explores how the choice of pivot rule (e.g., largest-coefficient, largest-increase, Bland's rule, randomized rule) affects the number of iterations required by the Simplex method. The analysis considers both empirical data and theoretical implications. %
\item \href{https://www.math.purdue.edu/~yipn/421/2024-09-05-SimplexEfficiency.pdf#page=5}{\textbf{Effect of Variable Scaling on Simplex Method Performance}}: This problem investigates how scaling the variables in a linear program can significantly impact the performance of the Simplex method, particularly when using rules like the largest-coefficient rule. %
\item \href{https://www.math.purdue.edu/~yipn/421/2024-09-05-SimplexEfficiency.pdf#page=9}{\textbf{Empirical vs. Theoretical Analysis of Simplex Method}}: This problem compares empirical observations of Simplex method performance with theoretical bounds and analyses, highlighting the gap between average-case behavior and worst-case scenarios. The discussion includes both experimental data and theoretical results on smoothed complexity. %
\item \href{https://www.math.purdue.edu/~yipn/421/2024-09-12-Duality.pdf#page=1}{\textbf{Formulating Primal and Dual Linear Programs}}: This problem involves expressing a given optimization problem in both primal and dual forms, identifying the objective functions, constraints, and variables in each. %
\item \href{https://www.math.purdue.edu/~yipn/421/2024-09-12-Duality.pdf#page=2}{\textbf{Applying Weak Duality Theorem}}: This problem involves verifying the inequality relationship between the objective function values of feasible primal and dual solutions. %
\item \href{https://www.math.purdue.edu/~yipn/421/2024-09-12-Duality.pdf#page=3}{\textbf{Representing Linear Programs in Matrix Form}}: This problem focuses on expressing primal and dual linear programs using matrices and vectors, highlighting the relationship between the primal and dual coefficient matrices. %
\item \href{https://www.math.purdue.edu/~yipn/421/2024-09-12-Duality.pdf#page=5}{\textbf{Understanding Complementary Variables and their Relationship in Simplex Iterations}}: This problem involves identifying and tracking complementary variables (primal and dual variables and their slack variables) during simplex iterations, observing their behavior as variables enter and leave the basis. %
\item \href{https://www.math.purdue.edu/~yipn/421/2024-09-12-Duality.pdf#page=6}{\textbf{Demonstrating the Preservation of the Negative Transpose Property}}: This problem involves showing that the negative transpose relationship between the primal and dual tableaux is maintained throughout the simplex method iterations. %
\item \href{https://www.math.purdue.edu/~yipn/421/2024-09-12-Duality.pdf#page=9}{\textbf{Analyzing the Optimal Tableau and its Properties}}: This problem involves examining the characteristics of the optimal tableau for both primal and dual problems, specifically focusing on the non-negativity of coefficients in the objective row and constraint rows. %
\item \href{https://www.math.purdue.edu/~yipn/421/2024-09-12-Duality.pdf#page=11}{\textbf{Proving the Strong Duality Theorem}}: This problem involves demonstrating the equality of optimal objective function values for the primal and dual problems, given that an optimal solution exists for the primal problem. %
\item \href{https://www.math.purdue.edu/~yipn/421/2024-09-12-Duality.pdf#page=21}{\textbf{Applying the Complementary Slackness Theorem}}: This problem involves verifying the optimality of a given primal-dual solution pair by checking the complementary slackness conditions. %
\item \href{https://www.math.purdue.edu/~yipn/421/2024-09-17-DualityGeneral.pdf#page=1}{\textbf{Formulating primal and dual linear programs}}: This problem involves expressing a given optimization problem in both primal and dual forms, identifying the objective function and constraints for each. %
\item \href{https://www.math.purdue.edu/~yipn/421/2024-09-17-DualityGeneral.pdf#page=11}{\textbf{Determining the existence of a non{-}negative solution to Ax ≤ b}}: This problem focuses on finding a vector x with non-negative components that satisfies the inequality Ax ≤ b. %
\item \href{https://www.math.purdue.edu/~yipn/421/2024-09-17-DualityGeneral.pdf#page=5}{\textbf{Analyzing the relationship between primal and dual objective function values}}: This problem involves comparing the optimal values of the primal and dual objective functions, establishing the weak and strong duality theorems. %
\item \href{https://www.math.purdue.edu/~yipn/421/2024-09-17-DualityGeneral.pdf#page=22}{\textbf{Applying the complementary slackness theorem}}: This problem involves using the complementary slackness theorem to determine relationships between optimal primal and dual solutions. %
\item \href{https://www.math.purdue.edu/~yipn/421/2024-09-17-DualityGeneral.pdf#page=3}{\textbf{Maximizing the Lagrangian function}}: This problem involves finding the maximum value of the Lagrangian function with respect to the primal variables, under different constraint conditions. %
\item \href{https://www.math.purdue.edu/~yipn/421/2024-09-19-DualityExamples.pdf#page=1}{\textbf{Formulating Primal and Dual Linear Programs}}: This problem involves constructing both the primal and dual linear programming formulations for a given problem, such as the diet problem, where the primal minimizes cost and the dual maximizes profit. %
\item \href{https://www.math.purdue.edu/~yipn/421/2024-09-19-DualityExamples.pdf#page=1}{\textbf{Analyzing the Relationship Between Primal and Dual Objective Function Values}}: This problem focuses on understanding the relationship between the optimal objective function values of the primal and dual problems, particularly in the context of strong duality. %
\item \href{https://www.math.purdue.edu/~yipn/421/2024-09-19-DualityExamples.pdf#page=1}{\textbf{Applying the Complementary Slackness Theorem}}: This problem involves using the complementary slackness theorem to analyze the relationship between primal and dual optimal solutions. %
\item \href{https://www.math.purdue.edu/~yipn/421/2024-09-19-DualityExamples.pdf#page=2}{\textbf{Diet Problem Formulation and Interpretation}}: This problem involves formulating and interpreting the diet problem as a linear program, including its constraints and objective function, and understanding its dual interpretation as a profit maximization problem. %
\item \href{https://www.math.purdue.edu/~yipn/421/2024-09-24-SimplexMatrix.pdf#page=1}{\textbf{Formulating Linear Programs in Matrix Notation}}: Representing linear programming problems using matrices and vectors for efficient computation. %
\item \href{https://www.math.purdue.edu/~yipn/421/2024-09-24-SimplexMatrix.pdf#page=4}{\textbf{Partitioning Variables into Basic and Non{-}Basic Sets}}: Dividing variables into basic and non-basic sets to facilitate iterative solutions using the simplex method. %
\item \href{https://www.math.purdue.edu/~yipn/421/2024-09-24-SimplexMatrix.pdf#page=8}{\textbf{Expressing the Objective Function in Terms of Basic and Non{-}Basic Variables}}: Rewriting the objective function to depend only on non-basic variables for efficient updates during simplex iterations. %
\item \href{https://www.math.purdue.edu/~yipn/421/2024-09-24-SimplexMatrix.pdf#page=10}{\textbf{Updating the Simplex Tableau Using Matrix Operations}}: Performing matrix operations to update the simplex tableau during each iteration of the simplex method. %
\item \href{https://www.math.purdue.edu/~yipn/421/2024-09-24-SimplexMatrix.pdf#page=14}{\textbf{Understanding Complementary Variables in Simplex Iterations}}: Analyzing the relationship between primal and dual variables during simplex iterations. %
\item \href{https://www.math.purdue.edu/~yipn/421/2024-09-24-SimplexMatrix.pdf#page=16}{\textbf{Solving a Linear Program Graphically}}: Finding the optimal solution of a linear program by visualizing the feasible region and objective function. %
\item \href{https://www.math.purdue.edu/~yipn/421/2024-09-24-SimplexMatrix.pdf#page=15}{\textbf{Illustrating Simplex Iterations with a Numerical Example}}: Step-by-step demonstration of the simplex method using a specific linear programming problem. %
\item \href{https://www.math.purdue.edu/~yipn/421/2024-09-26-Sensitivity.pdf#page=4}{\textbf{Analyzing the Sensitivity of the Objective Function Coefficients}}: This problem investigates how changes in the coefficients of the objective function affect the optimal solution. A parameter λ is introduced to represent the change in the coefficient of x₁, and the analysis determines the range of λ for which the optimal solution remains unchanged. %
\item \href{https://www.math.purdue.edu/~yipn/421/2024-09-26-Sensitivity.pdf#page=8}{\textbf{Analyzing the Sensitivity of the Constraint Values}}: This problem explores how changes in the right-hand side (RHS) values of the constraints affect the optimal solution. Parameters Δb₁, Δb₂, and Δb₃ are introduced to represent these changes, and the analysis determines the range of these parameters for which the optimal solution remains unchanged. %
\item \href{https://www.math.purdue.edu/~yipn/421/2024-10-01-DualGeneralLP.pdf#page=1}{\textbf{Deriving the Dual of a General Linear Program}}: This problem focuses on deriving the dual linear program from a given primal linear program, particularly when the primal problem does not have non-negativity constraints on all variables. The derivation involves introducing slack variables and using the Lagrangian function. %
\item \href{https://www.math.purdue.edu/~yipn/421/2024-10-01-DualGeneralLP.pdf#page=5}{\textbf{Proving Weak Duality}}: This problem involves demonstrating that the objective function value of any feasible solution to the primal problem is less than or equal to the objective function value of any feasible solution to the dual problem. %
\item \href{https://www.math.purdue.edu/~yipn/421/2024-10-01-SensitivityDuality.pdf#page=13}{\textbf{Verifying Optimality}}: A method for verifying the optimality of a solution is presented, involving checking the objective function value and the feasibility of the dual solution. %
\item \href{https://www.math.purdue.edu/~yipn/421/2024-10-01-SensitivityDuality.pdf#page=1}{\textbf{Sensitivity Analysis of Constraint Values}}: Analyzing how changes in the right-hand side of constraints affect the optimal solution and objective function value. %
\item \href{https://www.math.purdue.edu/~yipn/421/2024-10-01-SensitivityDuality.pdf#page=3}{\textbf{Applying Complementary Slackness}}: Using complementary slackness to verify the optimal solution by relating non-zero primal variables to zero dual slack variables. %
\item \href{https://www.math.purdue.edu/~yipn/421/2024-10-01-SensitivityDuality.pdf#page=4}{\textbf{Solving Linear Programs Using the Basis and Non{-}Basis Method}}: Utilizing the basis and non-basis method (likely the simplex method) to solve linear programming problems, involving matrix operations and iterative improvements. %
\item \href{https://www.math.purdue.edu/~yipn/421/2024-10-01-SensitivityDuality.pdf#page=2}{\textbf{Analyzing the Relationship Between Primal and Dual Problems}}: Examining the relationship between the primal and dual problems, including strong duality and complementary slackness. %
\item \href{https://www.math.purdue.edu/~yipn/421/2024-10-03-DualSimplexMatrix.pdf#page=2}{\textbf{Formulating Primal and Dual Linear Programs}}: This problem involves constructing both the primal and dual linear programming formulations from a given problem description. %
\item \href{https://www.math.purdue.edu/~yipn/421/2024-10-03-DualSimplexMatrix.pdf#page=1}{\textbf{Solving Linear Programs Graphically}}: This problem focuses on finding the optimal solution of a linear program by visualizing the feasible region and objective function in a graphical representation. %
\item \href{https://www.math.purdue.edu/~yipn/421/2024-10-03-DualSimplexMatrix.pdf#page=7}{\textbf{Applying the Simplex Method in Matrix Form}}: This problem involves using matrix operations to perform simplex iterations, including selecting entering and leaving variables and updating the basis matrix. %
\item \href{https://www.math.purdue.edu/~yipn/421/2024-10-03-DualSimplexMatrix.pdf#page=1}{\textbf{Analyzing the Relationship Between Primal and Dual Problems}}: This problem explores the connection between primal and dual solutions, including their optimal values and the implications of feasibility and optimality for both problems. %
\item \href{https://www.math.purdue.edu/~yipn/421/2024-10-03-DualSimplexMatrix.pdf#page=5}{\textbf{Handling Infeasible Origins in Dual Based Phase I Algorithm}}: This problem addresses situations where the initial point (origin) is not feasible for either the primal or dual problem, requiring a Phase I procedure to find a feasible starting point. %
\item \href{https://www.math.purdue.edu/~yipn/421/2024-10-22-Convexity.pdf#page=1}{\textbf{Defining Convex Sets and Convex Combinations}}: This problem involves understanding the definition of a convex set and how to form convex combinations of points within a set. The geometric interpretation of convexity is also explored. %
\item \href{https://www.math.purdue.edu/~yipn/421/2024-10-22-Convexity.pdf#page=2}{\textbf{Proving Properties of Convex Sets}}: This problem focuses on proving theorems related to convex sets, such as the equivalence between a set being convex and containing all convex combinations of its points. %
\item \href{https://www.math.purdue.edu/~yipn/421/2024-10-22-Convexity.pdf#page=5}{\textbf{Defining and Proving Properties of Convex Hulls}}: This problem involves understanding the definition of the convex hull of a set and proving theorems related to its properties, such as its characterization as the set of all convex combinations of finitely many points from the original set. %
\item \href{https://www.math.purdue.edu/~yipn/421/2024-10-22-Convexity.pdf#page=9}{\textbf{Applying the Caratheodory Theorem}}: This problem involves understanding and applying the Caratheodory Theorem, which states that any point in the convex hull of a set in $\mathbb{R}^m$ can be expressed as a convex combination of at most $m+1$ points from the original set. This is demonstrated through the use of linear programming. %
\item \href{https://www.math.purdue.edu/~yipn/421/2024-10-22-Convexity.pdf#page=10}{\textbf{Proving the Separation Theorem}}: This problem involves understanding and proving the Separation Theorem, which states that two disjoint nonempty polyhedra can be separated by two disjoint half-spaces. The proof utilizes Farkas' Lemma. %
\item \href{https://www.math.purdue.edu/~yipn/421/2024-10-22-Regression.pdf#page=1}{\textbf{Linear Regression using different error functions}}: This problem focuses on finding the best-fit line (or hyperplane) for a given dataset by minimizing different error functions, such as the sum of squared errors (L²), sum of absolute errors (L¹), and maximum absolute error (L∞). %
\item \href{https://www.math.purdue.edu/~yipn/421/2024-10-22-Regression.pdf#page=7}{\textbf{Binary Classification using a hyperplane}}: This problem involves finding a hyperplane that optimally separates data points into two classes. The goal is to find a hyperplane that maximizes the margin between the classes. %
\item \href{https://www.math.purdue.edu/~yipn/421/2024-10-22-Regression.pdf#page=15}{\textbf{Finding the Maximum Inscribed Circle in a Convex Polyhedron}}: This problem aims to find the largest circle that can fit entirely within a given convex polyhedron. %
\item \href{https://www.math.purdue.edu/~yipn/421/2024-10-22-Regression.pdf#page=22}{\textbf{Finding the Smallest Enclosing Ball}}: This problem involves finding the smallest-radius sphere that contains a given set of points in a d-dimensional space. %
\item \href{https://www.math.purdue.edu/~yipn/421/2024-10-28-ApplsLinProg.pdf#page=1}{\textbf{Production Planning}}: Developing a yearly production plan for an ice cream manufacturer to meet fluctuating monthly demands while minimizing the costs of changing production levels and storing surplus inventory. %
\item \href{https://www.math.purdue.edu/~yipn/421/2024-10-28-ApplsLinProg.pdf#page=7}{\textbf{Workforce Planning}}: Determining the optimal mix of regular and temporary employees to meet fluctuating monthly labor demands while minimizing total labor costs. %
\item \href{https://www.math.purdue.edu/~yipn/421/2024-10-28-ApplsLinProg.pdf#page=9}{\textbf{Portfolio Selection (Markowitz Model)}}: Optimizing an investment portfolio to maximize expected return while considering and managing risk. %
\item \href{https://www.math.purdue.edu/~yipn/421/2024-10-28-ApplsLinProg.pdf#page=17}{\textbf{Smallest Enclosing Ball Problem}}: Finding the smallest ball (circle in 2D, hypersphere in higher dimensions) that contains a given set of points. %
\item \href{https://www.math.purdue.edu/~yipn/421/2024-10-28-ApplsLinProg.pdf#page=19}{\textbf{Maximum Weight Matching}}: Finding a matching in a weighted graph that maximizes the total weight of the matched edges. %
\item \href{https://www.math.purdue.edu/~yipn/421/2024-10-28-ApplsLinProg.pdf#page=22}{\textbf{Machine Scheduling}}: Determining the optimal order to process a set of printing jobs on different types of machines to minimize the total processing time. %
\item \href{https://www.math.purdue.edu/~yipn/421/2024-10-28-ApplsLinProg.pdf#page=24}{\textbf{Knapsack Problem}}: Selecting a subset of items with weights and values to maximize the total value while not exceeding a weight capacity. %
\item \href{https://www.math.purdue.edu/~yipn/421/2024-10-28-ApplsLinProg.pdf#page=25}{\textbf{Cutting Stock Problem}}: Determining how to cut large rolls of paper into smaller rolls of various widths to minimize waste while meeting customer demands. %
\item \href{https://www.math.purdue.edu/~yipn/421/2024-10-28-ApplsLinProg.pdf#page=30}{\textbf{Sparse Solutions of Linear Systems}}: Finding a solution to an underdetermined system of linear equations that has a limited number of non-zero elements. %
\item \href{https://www.math.purdue.edu/~yipn/421/2024-10-28-Farkas.pdf#page=2}{\textbf{Determining the existence of a non{-}negative solution to `Ax = b`}}: This problem investigates whether there exists a vector `x` with non-negative components that satisfies the equation `Ax = b`. Several equivalent conditions are presented using a vector `y`. %
\item \href{https://www.math.purdue.edu/~yipn/421/2024-10-28-Farkas.pdf#page=4}{\textbf{Determining the existence of a solution to `Ax ≤ b`}}: This problem explores whether a solution `x` exists that satisfies the system of inequalities `Ax ≤ b`. Conditions for the existence (and non-existence) of such a solution are examined using a vector `y`. %
\item \href{https://www.math.purdue.edu/~yipn/421/2024-10-28-Farkas.pdf#page=1}{\textbf{Proving Farkas' Lemma}}: This problem focuses on proving different variants of Farkas' Lemma, which provides conditions for the existence or non-existence of solutions to systems of linear inequalities and equations of the form `Ax = b` and `Ax ≤ b`, with or without non-negativity constraints on `x`. The proof utilizes geometric interpretations and properties of convex sets. %
\item \href{https://www.math.purdue.edu/~yipn/421/2024-10-31-IntLinProg.pdf#page=2}{\textbf{Solving Integer Programming Problems by Relaxing and Rounding}}: This problem explores the limitations of solving an integer program by first solving its linear relaxation and then rounding the solution. The method is shown to be unreliable because the rounded solution may be infeasible or far from the true optimal integer solution. %
\item \href{https://www.math.purdue.edu/~yipn/421/2024-10-31-IntLinProg.pdf#page=3}{\textbf{Solving a Specific Integer Programming Problem}}: This problem involves maximizing a linear objective function subject to linear inequality constraints, with the added constraint that the decision variables must be integers. A graphical method is used to illustrate the feasible region and the difference between the optimal solution of the relaxed problem and the optimal integer solution. %
\item \href{https://www.math.purdue.edu/~yipn/421/2024-10-31-IntLinProg.pdf#page=5}{\textbf{Applying the Branch and Bound Method}}: This problem demonstrates the application of the branch and bound algorithm to solve an integer programming problem. The algorithm is illustrated using a graphical representation of the feasible region and an enumeration tree to track the branching and bounding steps. %
\item \href{https://www.math.purdue.edu/~yipn/421/2024-10-31-IntLinProg.pdf#page=11}{\textbf{Justifying Depth{-}First Search in Branch and Bound}}: This problem analyzes the efficiency of depth-first search versus breadth-first search in the branch and bound algorithm. It argues that depth-first search is generally preferred due to the tendency of integer solutions to lie deep in the search tree, enabling early feasible solution identification and pruning. %
\item \href{https://www.math.purdue.edu/~yipn/421/2024-10-31-IntLinProg.pdf#page=15}{\textbf{Applying Gomory Cuts to Solve Integer Programming Problems}}: This problem demonstrates the use of Gomory cuts to find integer solutions to linear programs. The method involves deriving cuts from the optimal dictionary of the relaxed problem and iteratively adding them to refine the feasible region until an integer solution is found. %
\item \href{https://www.math.purdue.edu/~yipn/421/2024-11-05-NetFlows.pdf#page=1}{\textbf{Formulating Network Flow Problems as Linear Programs}}: This problem involves representing a network flow problem, with its nodes, arcs, supplies, demands, and costs, as a linear program with an objective function and constraints. %
\item \href{https://www.math.purdue.edu/~yipn/421/2024-11-05-NetFlows.pdf#page=5}{\textbf{Analyzing the Incidence Matrix of a Network}}: This problem focuses on understanding the properties of the incidence matrix, including its rank and the relationship between its submatrices and spanning trees of the network. %
\item \href{https://www.math.purdue.edu/~yipn/421/2024-11-05-NetFlows.pdf#page=10}{\textbf{Finding Primal and Dual Solutions using Spanning Trees}}: This problem involves finding primal and dual feasible solutions to a network flow problem by utilizing spanning trees and applying complementary slackness. %
\item \href{https://www.math.purdue.edu/~yipn/421/2024-11-05-NetFlows.pdf#page=22}{\textbf{Applying the Primal Simplex Method to Network Flows}}: This problem involves using the primal simplex method to iteratively improve a feasible tree solution until an optimal solution is reached. %
\item \href{https://www.math.purdue.edu/~yipn/421/2024-11-07-NetSimplex.pdf#page=2}{\textbf{Finding a Primal Feasible Flow}}: This problem involves finding a flow vector x that satisfies the constraints Bx = b, where x is non-negative and some components of x are zero based on the spanning tree. %
\item \href{https://www.math.purdue.edu/~yipn/421/2024-11-07-NetSimplex.pdf#page=3}{\textbf{Finding a Dual Feasible Solution}}: This problem involves finding a dual solution (y, z) that satisfies the dual constraints, which relate the dual variables to the costs associated with the edges in the network. %
\item \href{https://www.math.purdue.edu/~yipn/421/2024-11-07-NetSimplex.pdf#page=4}{\textbf{Understanding Spanning Tree Properties}}: This problem focuses on understanding how adding or removing an arc from a spanning tree affects the tree's connectivity and the presence of cycles. %
\item \href{https://www.math.purdue.edu/~yipn/421/2024-11-07-NetSimplex.pdf#page=5}{\textbf{Primal Network Simplex Method}}: This problem involves iteratively improving a primal feasible solution by selecting entering and leaving arcs to modify the spanning tree and update primal and dual flows. %
\item \href{https://www.math.purdue.edu/~yipn/421/2024-11-07-NetSimplex.pdf#page=32}{\textbf{Dual Network Simplex Method}}: This problem involves iteratively improving a dual feasible solution by selecting entering and leaving arcs to modify the spanning tree and update primal and dual flows. %
\item \href{https://www.math.purdue.edu/~yipn/421/2024-11-07-NetSimplex.pdf#page=45}{\textbf{Developing Variants of the Network Simplex Method}}: This problem explores different approaches to solving network flow problems using the network simplex method, including two-phased methods and a parametric self-dual method. %
\item \href{https://www.math.purdue.edu/~yipn/421/2024-11-14-NetAppls.pdf#page=1}{\textbf{Formulating Transportation Problems with Inequality Constraints}}: This problem involves modeling transportation problems where supply may exceed demand or vice versa, requiring the use of dummy nodes and slack variables to ensure balance. %
\item \href{https://www.math.purdue.edu/~yipn/421/2024-11-14-NetAppls.pdf#page=5}{\textbf{Formulating Transportation Problems with Inequality Constraints (Graphical Representation)}}: This problem focuses on visually representing transportation problems with inequality constraints using network flow diagrams, illustrating the flow of goods between sources and destinations, including the use of dummy nodes to handle excess supply or unmet demand. %
\item \href{https://www.math.purdue.edu/~yipn/421/2024-11-14-NetAppls.pdf#page=6}{\textbf{Formulating Transportation Problems with Inequality Constraints (Algebraic Representation)}}: This problem involves expressing the constraints of a transportation problem algebraically, using matrices and vectors to represent the supply, demand, and flow variables, including the use of slack variables to convert inequalities into equalities. %
\item \href{https://www.math.purdue.edu/~yipn/421/2024-11-14-NetAppls.pdf#page=9}{\textbf{Solving the Upper Bounded Transshipment Problem}}: This problem involves finding optimal flows in a network with upper bounds on the flow through each arc, requiring the introduction of slack variables to handle the upper bound constraints. %
\item \href{https://www.math.purdue.edu/~yipn/421/2024-11-14-NetAppls.pdf#page=16}{\textbf{Solving the Shortest Path Problem using Bellman's Equation}}: This problem focuses on finding the shortest path between two nodes in a network using Bellman's equation, an iterative dynamic programming approach that approximates the shortest distance from each node to the destination. %
\item \href{https://www.math.purdue.edu/~yipn/421/2024-11-14-NetAppls.pdf#page=20}{\textbf{Solving the Shortest Path Problem using Dijkstra's Algorithm}}: This problem involves finding the shortest path between two nodes in a network using Dijkstra's algorithm, an iterative algorithm that maintains a set of processed nodes and updates shortest distance estimates for unprocessed nodes. %
\item \href{https://www.math.purdue.edu/~yipn/421/2024-11-19-MaxFlowMinCut.pdf#page=1}{\textbf{Finding the Maximum Flow in a Network}}: This problem involves determining the maximum amount of flow that can be sent from a source node to a sink node in a network, subject to capacity constraints on the edges. %
\item \href{https://www.math.purdue.edu/~yipn/421/2024-11-19-MaxFlowMinCut.pdf#page=8}{\textbf{Proving the Max{-}Flow Min{-}Cut Theorem}}: This problem focuses on proving the equivalence between the maximum flow in a network and the minimum capacity of a cut separating the source and sink nodes. %
\item \href{https://www.math.purdue.edu/~yipn/421/2024-11-19-MaxFlowMinCut.pdf#page=10}{\textbf{Implementing the Ford{-}Fulkerson Algorithm}}: This problem involves understanding and implementing the Ford-Fulkerson algorithm, a method for finding the maximum flow in a network by iteratively finding augmenting paths. %
\item \href{https://www.math.purdue.edu/~yipn/421/2024-11-19-MaxFlowMinCut.pdf#page=14}{\textbf{Formulating the Maximum Flow Problem as a Linear Program}}: This problem involves expressing the maximum flow problem in the standard form of a linear program, defining the objective function, decision variables, and constraints. %
\item \href{https://www.math.purdue.edu/~yipn/421/2024-11-21-EconNetSimplex.pdf#page=4}{\textbf{Finding Fair Prices in Network Simplex}}: Determining a set of market prices ($y_i$) at each node such that the price at node $j$ equals the price at node $i$ plus the cost of shipping along arc $ij$ ($c_{ij}$) for each arc $ij$ in the spanning tree. %
\item \href{https://www.math.purdue.edu/~yipn/421/2024-11-21-EconNetSimplex.pdf#page=5}{\textbf{Identifying Profitable Opportunities for Competitors}}: Finding an arc $ij$ where the cost of buying at node $i$ and shipping to node $j$ ($y_i + c_{ij}$) is less than the selling price at node $j$ ($y_j$). %
\item \href{https://www.math.purdue.edu/~yipn/421/2024-11-21-EconNetSimplex.pdf#page=6}{\textbf{Adjusting Shipments to Maintain Feasibility}}: Adjusting shipments along the arcs of the spanning tree to maintain feasibility after introducing a new arc identified in the previous step. %
\item \href{https://www.math.purdue.edu/~yipn/421/2024-11-21-EconNetSimplex.pdf#page=3}{\textbf{Finding a Feasible Solution on a Spanning Tree}}: Finding an initial feasible solution to a network flow problem that is represented by a spanning tree. %
\item \href{https://www.math.purdue.edu/~yipn/421/2024-12-03-Duality.pdf#page=1}{\textbf{Formulating the Dual Problem}}: This problem involves transforming a given primal linear programming problem into its corresponding dual problem. The transformation involves transposing the constraint matrix and swapping the objective function coefficients with the right-hand side constants of the primal constraints. %
\item \href{https://www.math.purdue.edu/~yipn/421/2024-12-03-Duality.pdf#page=3}{\textbf{Solving Linear Programs Graphically}}: This problem focuses on solving linear programming problems with two variables using a graphical method. The method involves plotting the constraints to define a feasible region and then finding the point within this region that maximizes or minimizes the objective function. %
\item \href{https://www.math.purdue.edu/~yipn/421/2024-12-03-Duality.pdf#page=2}{\textbf{Applying Complementary Slackness}}: This problem involves using the complementary slackness theorem to determine the optimal solutions to both the primal and dual linear programming problems. The theorem states that at optimality, the product of a primal variable and its corresponding dual slack variable must be zero, and similarly for the dual variable and its corresponding primal slack variable. %
\item \href{https://www.math.purdue.edu/~yipn/421/2024-12-03-Duality.pdf#page=11}{\textbf{Performing Sensitivity Analysis}}: This problem involves analyzing how changes in the objective function coefficients or constraint values affect the optimal solution of a linear programming problem. This includes determining the range of parameter changes that will not affect the optimality of the solution. %
\item \href{https://www.math.purdue.edu/~yipn/421/2024-12-03-Duality.pdf#page=9}{\textbf{Geometric Interpretation of Dual Constraints}}: This problem involves providing a graphical interpretation of the dual constraints in the context of linear programming. This includes visualizing the half-spaces defined by the dual constraints and determining the optimal solution to the dual problem as the intersection of these half-spaces. %
\item \href{https://www.math.purdue.edu/~yipn/421/2024-12-03-Duality.pdf#page=22}{\textbf{Solving Linear Programs Using Matrix Methods}}: This problem involves solving linear programming problems using matrix methods, such as the simplex method in matrix form. This includes working with the basis matrix, non-basic variables, and performing matrix operations to find the optimal solution. %
\end{itemize}

%
\section{Algorithm Solutions}%
\label{sec:AlgorithmSolutions}%
\begin{itemize}[leftmargin=*]%
\item \href{https://www.math.purdue.edu/~yipn/421/2024-08-27-ExSimplex.pdf#page=6}{\textbf{Simplex Method}}: Iteratively move from one vertex of the feasible region to another by setting non-basic variables to zero and choosing a new set of (non)basic variables to improve the objective function, until the optimal solution is found.%
\item \href{https://www.math.purdue.edu/~yipn/421/2024-08-27-ExSimplex.pdf#page=30}{\textbf{Handling Infeasible Origins}}: Introduce an auxiliary problem with an additional variable  $x_0$ to shift the feasible region, making the origin feasible. Solve the auxiliary problem using the Simplex method. If the optimal solution has $x_0 = 0$, then this solution is feasible for the original problem.%
\item \href{https://www.math.purdue.edu/~yipn/421/2024-08-27-ExSimplex.pdf#page=9}{\textbf{Pivot Operation}}: Interchange one basic and one non-basic variable in the dictionary by rewriting the equations to express the basic variables in terms of the non-basic variables. This corresponds to moving from one vertex to an adjacent vertex in the feasible region.%
\item \href{https://www.math.purdue.edu/~yipn/421/2024-09-05-SimplexEfficiency.pdf#page=3}{\textbf{Largest{-}Coefficient Rule}}: Choose the variable with the largest coefficient in the objective function row to enter the basis. %
\item \href{https://www.math.purdue.edu/~yipn/421/2024-09-05-SimplexEfficiency.pdf#page=6}{\textbf{Largest{-}Increase Rule}}: Choose the variable whose entrance into the basis brings about the largest increase in the objective function. %
\item \href{https://www.math.purdue.edu/~yipn/421/2024-09-05-SimplexEfficiency.pdf#page=10}{\textbf{Bland's Rule}}: A rule for choosing entering and leaving variables in the simplex method to avoid cycling. %
\item \href{https://www.math.purdue.edu/~yipn/421/2024-09-05-SimplexEfficiency.pdf#page=10}{\textbf{Randomized Pivot Rule}}: Randomly permute the indices of the variables in the input linear program; then use Bland's rule for choosing the entering variable, and the lexicographic method for choosing the leaving variable. %
\item \href{https://www.math.purdue.edu/~yipn/421/2024-09-12-Duality.pdf#page=6}{\textbf{Simplex Method}}: An iterative algorithm that moves from one feasible solution to another, improving the objective function value at each step until an optimal solution is found. The algorithm involves selecting an entering and leaving variable based on rules like the largest-increase rule, and updating the tableau using pivot operations. %
\item \href{https://www.math.purdue.edu/~yipn/421/2024-09-12-Duality.pdf#page=21}{\textbf{Complementary Slackness}}: Primal and dual feasible solutions x and y are optimal if and only if $x_j z_j = 0$ for all j and $w_i y_i = 0$ for all i, where $z_j$ and $w_i$ are the dual and primal slack variables, respectively. %
\item \href{https://www.math.purdue.edu/~yipn/421/2024-09-17-DualityGeneral.pdf#page=2}{\textbf{Lagrange Multiplier}}: $L(x,\mu) = f_0(x) - \sum_{j=1}^n \mu_j g_j(x)$%
\item \href{https://www.math.purdue.edu/~yipn/421/2024-09-17-DualityGeneral.pdf#page=7}{\textbf{Lagrangian Function (with inequality and equality constraints)}}: $L(x,\lambda,\mu) = f_0(x) - \sum_{i=1}^m \lambda_i f_i(x) - \sum_{j=1}^n \mu_j g_j(x)$%
\item \href{https://www.math.purdue.edu/~yipn/421/2024-09-17-DualityGeneral.pdf#page=5}{\textbf{Weak Duality}}: $\max_{x \in \mathcal{G}} \xi(x) \le \min_{\lambda \ge 0, \mu} \xi(\lambda,\mu)$%
\item \href{https://www.math.purdue.edu/~yipn/421/2024-09-17-DualityGeneral.pdf#page=17}{\textbf{Strong Duality}}: If $x^$ is optimal for (P) and $\lambda^$ is optimal for (D), then $\xi(x^) = \xi(\lambda^)$%
\item \href{https://www.math.purdue.edu/~yipn/421/2024-09-17-DualityGeneral.pdf#page=22}{\textbf{Complementary Slackness}}: If $x^$ is optimal for (P) and $\lambda^$ is optimal for (D), then $x^_i \lambda^_i = 0$ for all $i$%
\item \href{https://www.math.purdue.edu/~yipn/421/2024-09-19-DualityExamples.pdf#page=5}{\textbf{Simplex Method}}: An iterative algorithm that moves from one feasible solution to another, improving the objective function value at each step until an optimal solution is found. It involves pivot operations on a tableau representing the linear program. %
\item \href{https://www.math.purdue.edu/~yipn/421/2024-09-19-DualityExamples.pdf#page=8}{\textbf{Dual Simplex Method}}: Similar to the simplex method, but starts with an infeasible solution and iteratively improves feasibility while maintaining optimality. %
\item \href{https://www.math.purdue.edu/~yipn/421/2024-09-19-DualityExamples.pdf#page=1}{\textbf{Strong Duality}}: The optimal objective function values of the primal and dual problems are equal. %
\item \href{https://www.math.purdue.edu/~yipn/421/2024-09-19-DualityExamples.pdf#page=1}{\textbf{Complementary Slackness}}: A condition that relates the optimal solutions of the primal and dual problems. Specifically, if a primal constraint is not binding (strict inequality), then the corresponding dual variable is zero, and vice versa. %
\item \href{https://www.math.purdue.edu/~yipn/421/2024-09-24-SimplexMatrix.pdf#page=8}{\textbf{Simplex Method}}: Iteratively improve a feasible solution by exchanging a basic variable with a non-basic variable to increase the objective function value until optimality is reached. This involves updating the simplex tableau using matrix operations, partitioning variables into basic and non-basic sets, and utilizing pivot operations. %
\item \href{https://www.math.purdue.edu/~yipn/421/2024-09-24-SimplexMatrix.pdf#page=4}{\textbf{Expressing the Objective Function in Terms of Basic and Non{-}Basic Variables}}: Partition the variable vector $\mathbf{x}$ into basic ($\mathbf{x}_B$) and non-basic ($\mathbf{x}_N$) variables. Express the objective function as $\mathcal{Z} = \mathbf{c}_B^T \mathbf{x}_B + \mathbf{c}_N^T \mathbf{x}_N$. Solve for $\mathbf{x}_B$ in terms of $\mathbf{x}_N$ using the constraint equation $B\mathbf{x}_B + N\mathbf{x}_N = \mathbf{b}$ to obtain $\mathbf{x}_B = B^{-1}\mathbf{b} - B^{-1}N\mathbf{x}_N$. Substitute this into the objective function to express it solely in terms of non-basic variables. %
\item \href{https://www.math.purdue.edu/~yipn/421/2024-09-24-SimplexMatrix.pdf#page=16}{\textbf{Graphical Solution Method}}: Graphically represent the constraints as lines in the decision variable space. The feasible region is the intersection of all constraint regions. The optimal solution is the point in the feasible region that maximizes the objective function, often found at a vertex of the feasible region. %
\item \href{https://www.math.purdue.edu/~yipn/421/2024-09-26-Sensitivity.pdf#page=7}{\textbf{Sensitivity Analysis of Objective Function Coefficients}}: Determine $\Delta \tilde{Z}_N = (B^{-1}N)^T \Delta C_B - \Delta C_N$. The current optimal solution remains optimal if $\tilde{Z}_N^ + \Delta \tilde{Z}_N \ge 0$.%
\item \href{https://www.math.purdue.edu/~yipn/421/2024-09-26-Sensitivity.pdf#page=11}{\textbf{Sensitivity Analysis of Constraint Values}}: Determine $\Delta X_B^ = B^{-1}\Delta b$. The current optimal solution remains optimal if $X_B^ + \Delta X_B^ \ge 0$.%
\item \href{https://www.math.purdue.edu/~yipn/421/2024-10-01-DualGeneralLP.pdf#page=1}{\textbf{Dual of a General LP Problem}}: Given a primal problem $\max \vec{c}^T\vec{x}$ subject to $A\vec{x} \le \vec{b}$ and $\vec{x} \ge \vec{0}$, its dual is $\min \vec{b}^T\vec{y}$ subject to $A^T\vec{y} \ge \vec{c}$ and $\vec{y} \ge \vec{0}$.%
\item \href{https://www.math.purdue.edu/~yipn/421/2024-10-01-DualGeneralLP.pdf#page=2}{\textbf{Reformulating a General LP without Non{-}negativity Constraints}}: To solve a problem of the form $\max \vec{c}^T\vec{x}$ subject to $A\vec{x} \le \vec{b}$ and $\vec{x} \le \vec{u}$ (where $\vec{x}$ is not necessarily non-negative), rewrite it as $\max \vec{c}^T(\vec{x}^+ - \vec{x}^-)$ subject to $A(\vec{x}^+ - \vec{x}^-) \le \vec{b}$, $\vec{x}^+ \ge \vec{0}$, $\vec{x}^- \ge \vec{0}$, where $\vec{x} = \vec{x}^+ - \vec{x}^-$.%
\item \href{https://www.math.purdue.edu/~yipn/421/2024-10-01-DualGeneralLP.pdf#page=3}{\textbf{Deriving the Dual of a General LP with Upper Bounds}}: The dual of $\max \vec{c}^T\vec{x}$ subject to $A\vec{x} \le \vec{b}$ and $\vec{x} \le \vec{u}$ is $\min \vec{b}^T\vec{y} + \vec{u}^T\vec{z}$ subject to $A^T\vec{y} + \vec{z} \ge \vec{c}$ and $\vec{y}, \vec{z} \ge \vec{0}$.%
\item \href{https://www.math.purdue.edu/~yipn/421/2024-10-01-DualGeneralLP.pdf#page=4}{\textbf{Lagrangian Function for a General LP with Upper Bounds}}: For the primal problem $\max \vec{c}^T\vec{x}$ subject to $A\vec{x} \le \vec{b}$ and $\vec{x} \le \vec{u}$, the Lagrangian is $\mathcal{L}(\vec{x}, \vec{y}, \vec{z}) = \vec{c}^T\vec{x} - \vec{y}^T(A\vec{x} - \vec{b}) - \vec{z}^T(\vec{x} - \vec{u})$, where $\vec{y} \ge \vec{0}$ and $\vec{z} \ge \vec{0}$.%
\item \href{https://www.math.purdue.edu/~yipn/421/2024-10-01-DualGeneralLP.pdf#page=7}{\textbf{Final Form of the Dual Problem with Upper Bounds}}: The dual problem is $\min_{\vec{y}, \vec{z} \ge \vec{0}} \vec{y}^T\vec{b} + \vec{z}^T\vec{u}$ subject to $A^T\vec{y} + \vec{z} = \vec{c}$.%
\item \href{https://www.math.purdue.edu/~yipn/421/2024-10-01-SensitivityDuality.pdf#page=4}{\textbf{Simplex Method}}: This algorithm iteratively improves a feasible solution to a linear program by moving from one vertex of the feasible region to an adjacent vertex with a better objective function value, until an optimal solution is found.%
\item \href{https://www.math.purdue.edu/~yipn/421/2024-10-01-SensitivityDuality.pdf#page=3}{\textbf{Complementary Slackness}}: This theorem states that for an optimal solution to a linear program and its dual, if a primal variable is positive, then the corresponding dual constraint is binding (and vice versa).%
\item \href{https://www.math.purdue.edu/~yipn/421/2024-10-01-SensitivityDuality.pdf#page=1}{\textbf{Sensitivity Analysis}}: This technique analyzes how changes in the objective function coefficients or constraint values affect the optimal solution and objective function value of a linear program.%
\item \href{https://www.math.purdue.edu/~yipn/421/2024-10-01-SensitivityDuality.pdf#page=10}{\textbf{Shadow Price}}: This is the rate of change in the optimal objective function value for a unit change in the right-hand side of a constraint, given by the optimal dual variable corresponding to that constraint.%
\item \href{https://www.math.purdue.edu/~yipn/421/2024-10-03-DualSimplexMatrix.pdf#page=2}{\textbf{Dual Simplex Method}}: This algorithm solves linear programs by iteratively improving a dual-feasible solution until primal feasibility and optimality are achieved. It starts with a dual feasible solution and moves towards primal feasibility while maintaining dual feasibility. %
\item \href{https://www.math.purdue.edu/~yipn/421/2024-10-03-DualSimplexMatrix.pdf#page=5}{\textbf{Dual Based Phase I Algorithm}}: This algorithm finds an initial feasible solution for a linear program by solving its dual problem. If the origin is not feasible for both primal and dual problems, it modifies the dual problem to make the origin feasible, then solves the modified dual problem to obtain a feasible solution for the primal problem. %
\item \href{https://www.math.purdue.edu/~yipn/421/2024-10-03-DualSimplexMatrix.pdf#page=7}{\textbf{Simplex Method in Matrix Form}}: This algorithm solves linear programs using matrix operations. It iteratively improves a basic feasible solution by selecting an entering variable based on reduced costs, a leaving variable based on the ratio test, and updating the basis matrix using elementary matrices. %
\item \href{https://www.math.purdue.edu/~yipn/421/2024-10-22-Convexity.pdf#page=1}{\textbf{Convex Set}}: $S \subseteq \mathbb{R}^m$ is convex if for any $z_1, z_2 \in S \implies \lambda z_1 + (1-\lambda) z_2 \in S$, $0 \le \lambda \le 1$.%
\item \href{https://www.math.purdue.edu/~yipn/421/2024-10-22-Convexity.pdf#page=1}{\textbf{Convex Combination}}: $\sum_{i=1}^n \lambda_i z_i$, $\lambda_i \ge 0$, $\sum_{i=1}^n \lambda_i = 1$.%
\item \href{https://www.math.purdue.edu/~yipn/421/2024-10-22-Convexity.pdf#page=2}{\textbf{Theorem 10.1}}: $S$ is convex $\iff$ it contains all convex combinations of points in $S$.%
\item \href{https://www.math.purdue.edu/~yipn/421/2024-10-22-Convexity.pdf#page=5}{\textbf{Convex Hull}}: $\text{Conv}(S) = \bigcap \{C \text{ convex } | S \subseteq C\}$%
\item \href{https://www.math.purdue.edu/~yipn/421/2024-10-22-Convexity.pdf#page=5}{\textbf{Theorem 10.2}}: $S \subseteq \mathbb{R}^m$, $\text{Conv}(S) = \{\text{collection of all convex combinations of finitely many points from } S\}$.%
\item \href{https://www.math.purdue.edu/~yipn/421/2024-10-22-Convexity.pdf#page=8}{\textbf{Caratheodory's Theorem}}: If $z \in \text{Conv}(S)$, then $z$ can be written as a convex combination of at most $m+1$ points from $S$.%
\item \href{https://www.math.purdue.edu/~yipn/421/2024-10-22-Convexity.pdf#page=10}{\textbf{Separation Theorem}}: Let $P$ and $\bar{P}$ be two disjoint nonempty polyhedra in $\mathbb{R}^n$. Then there exist disjoint half-spaces $H$ and $\bar{H}$ such that $P \subset H$ and $\bar{P} \subset \bar{H}$.%
\item \href{https://www.math.purdue.edu/~yipn/421/2024-10-22-Convexity.pdf#page=12}{\textbf{Farkas' Lemma}}: The system $Ax \le b$ has no solution if and only if there exists a vector $y \ge 0$ such that $y^T A = 0$ and $y^T b < 0$.%
\item \href{https://www.math.purdue.edu/~yipn/421/2024-10-22-Regression.pdf#page=1}{\textbf{Mean}}: $x = \frac{x_1 + ... + x_m}{m}$%
\item \href{https://www.math.purdue.edu/~yipn/421/2024-10-22-Regression.pdf#page=1}{\textbf{Median}}: $x = \text{median}$%
\item \href{https://www.math.purdue.edu/~yipn/421/2024-10-22-Regression.pdf#page=1}{\textbf{Mid{-}range}}: $x = \frac{\min x_i + \max x_i}{2}$%
\item \href{https://www.math.purdue.edu/~yipn/421/2024-10-22-Regression.pdf#page=9}{\textbf{Maximum Margin Classifier}}: Maximize the minimum distance from the hyperplane to the closest data points.%
\item \href{https://www.math.purdue.edu/~yipn/421/2024-10-22-Regression.pdf#page=10}{\textbf{Support Vector Machine (SVM)}}: Find a hyperplane that maximizes the margin. %
\item \href{https://www.math.purdue.edu/~yipn/421/2024-10-22-Regression.pdf#page=11}{\textbf{Lagrangian Formulation of SVM}}: $\min_{\alpha} \frac{1}{2} \sum_{i=1}^N \sum_{j=1}^N \alpha_i \alpha_j y_i y_j x_i^T x_j - \sum_{i=1}^N \alpha_i$ s.t. $\sum_{i=1}^N \alpha_i y_i = 0$ and $0 \le \alpha_i \le C, \quad \forall i$%
\item \href{https://www.math.purdue.edu/~yipn/421/2024-10-22-Regression.pdf#page=13}{\textbf{SVM Optimization as a QP Problem (version 1)}}: Maximize $\delta$ s.t. $y_i (a_1 x_i^1 + a_2 x_i^2 - b) \ge \delta, \quad \forall i$%
\item \href{https://www.math.purdue.edu/~yipn/421/2024-10-22-Regression.pdf#page=14}{\textbf{SVM Optimization as a QP Problem (version 2)}}: Maximize $\delta$ s.t. $y_i (a_1 x_i^1 + a_2 x_i^2 - b) \ge \delta, \quad \forall i$ and $-1 \le a_1, a_2 \le 1$%
\item \href{https://www.math.purdue.edu/~yipn/421/2024-10-22-Regression.pdf#page=18}{\textbf{Distance of a Point from a Hyperplane}}: $d_i = \frac{|a_1 s_i + a_2 s_i^2 - b'|}{\sqrt{a_1^2 + a_2^2}}$%
\item \href{https://www.math.purdue.edu/~yipn/421/2024-10-22-Regression.pdf#page=19}{\textbf{SVM Optimization as a QP Problem (version 3)}}: Maximize $d$ s.t. $|a_1 s_i + a_2 s_i^2 - b'| \le d \sqrt{a_1^2 + a_2^2}, \quad \forall i$ and $a_1 s_i + a_2 s_i^2 - b' \le 0, \quad \forall i$%
\item \href{https://www.math.purdue.edu/~yipn/421/2024-10-22-Regression.pdf#page=20}{\textbf{SVM Optimization as a QP Problem (version 4)}}: Maximize $d$ s.t. $d \le \frac{|a_1 s_i + a_2 s_i^2 - b'|}{\sqrt{a_1^2 + a_2^2}}, \quad \forall i$ and $a_1 s_i + a_2 s_i^2 - b' \le 0, \quad \forall i$%
\item \href{https://www.math.purdue.edu/~yipn/421/2024-10-22-Regression.pdf#page=22}{\textbf{Smallest Ball Problem}}: Find a ball of the smallest radius that contains all the points $p_1, \dots, p_n \in \mathbb{R}^d$.%
\item \href{https://www.math.purdue.edu/~yipn/421/2024-10-28-ApplsLinProg.pdf#page=5}{\textbf{Ice Cream Production Planning}}: $\min \sum_{i=1}^{12} 50(Y_i + Z_i) + 20 \sum_{i=1}^{12} S_i$  s.t. $X_i + S_{i-1} - S_i = d_i$, $X_i - X_{i-1} = Y_i - Z_i$, $X_0 = 0$, $S_0 = 0$, $S_{12} = 0$, $X_i, Y_i, Z_i, S_i \ge 0$ for $i=1, \dots, 12$.%
\item \href{https://www.math.purdue.edu/~yipn/421/2024-10-28-ApplsLinProg.pdf#page=8}{\textbf{Workforce Planning}}: $\min 25 \sum_{i=1}^{12} Y_i + (17.5 \times 12) X$ s.t. $a_i - X \le Y_i$, $0 \le Y_i$, $X \ge 0$.%
\item \href{https://www.math.purdue.edu/~yipn/421/2024-10-28-ApplsLinProg.pdf#page=16}{\textbf{Portfolio Selection (Markowitz Model)}}: $\max \mu \sum_{j} x_j \overline{r_j} - \frac{\epsilon}{2} \sum_{i} \sum_{j} x_i x_j \sigma_{ij}$ s.t. $\sum_{j} x_j = 1$, $x_j \ge 0$.%
\item \href{https://www.math.purdue.edu/~yipn/421/2024-10-28-ApplsLinProg.pdf#page=18}{\textbf{Smallest Enclosing Ball}}: minimize $x^T Q^T Q x - \sum_{j=1}^n x_j p_j^T p_j$ subject to $\sum_{j=1}^n x_j = 1$, $x_j \ge 0$.%
\item \href{https://www.math.purdue.edu/~yipn/421/2024-10-28-ApplsLinProg.pdf#page=21}{\textbf{Maximum Weight Matching}}: $\max \sum_{e \in E} w_e x_e$ s.t. $\sum_{e \in \delta(v)} x_e = 1$ for all vertices $v$, $x_e \in \{0, 1\}$ for all edges $e$.%
\item \href{https://www.math.purdue.edu/~yipn/421/2024-10-28-ApplsLinProg.pdf#page=23}{\textbf{Machine Scheduling}}: $\min \pi$ s.t. $\sum_{i \in M} x_{ij} = 1$ for $j \in J$, $\sum_{j \in J} d_{ij} x_{ij} \le \pi$ for $i \in M$, $x_{ij} \in \{0, 1\}$ for all $i, j$.%
\item \href{https://www.math.purdue.edu/~yipn/421/2024-10-28-ApplsLinProg.pdf#page=24}{\textbf{Knapsack Problem}}: maximize $\sum_{j=1}^n c_j x_j$ subject to $\sum_{j=1}^n a_j x_j \le b$, $x_j \in \{0, 1\}$, $j = 1, 2, \dots, n$.%
\item \href{https://www.math.purdue.edu/~yipn/421/2024-10-28-ApplsLinProg.pdf#page=29}{\textbf{Cutting Stock Problem}}: $\min \sum_{i=1}^{12} x_i$ subject to constraints based on cutting patterns and demand.%
\item \href{https://www.math.purdue.edu/~yipn/421/2024-10-28-ApplsLinProg.pdf#page=34}{\textbf{Sparse Solutions of Linear System}}: $\min \|x\|_1$ s.t. $Ax = b$.%
\item \href{https://www.math.purdue.edu/~yipn/421/2024-10-28-Farkas.pdf#page=4}{\textbf{Farkas' Lemma (Variant 1)}}: The system $Ax = b$ has a nonnegative solution $x \ge 0$ if and only if every $y \in \mathbb{R}^m$ with $y^T A \ge 0^T$ also satisfies $y^T b \ge 0$.%
\item \href{https://www.math.purdue.edu/~yipn/421/2024-10-28-Farkas.pdf#page=4}{\textbf{Farkas' Lemma (Variant 2)}}: The system $Ax \le b$ has a nonnegative solution $x \ge 0$ if and only if every $y \in \mathbb{R}^m$ with $y^T A \ge 0^T$ also satisfies $y^T b \ge 0$.%
\item \href{https://www.math.purdue.edu/~yipn/421/2024-10-28-Farkas.pdf#page=4}{\textbf{Farkas' Lemma (Variant 3)}}: The system $Ax \le b$ has a solution $x \in \mathbb{R}^n$ if and only if every $y \in \mathbb{R}^m$ with $y^T A = 0^T$ also satisfies $y^T b \ge 0$.%
\item \href{https://www.math.purdue.edu/~yipn/421/2024-10-28-Farkas.pdf#page=6}{\textbf{Farkas' Lemma (Infeasibility Condition)}}: $Ax \le b$ has no solution if and only if there is $y \ge 0$ such that $y^T A \ge 0^T$ and $y^T b < 0$.%
\item \href{https://www.math.purdue.edu/~yipn/421/2024-10-31-IntLinProg.pdf#page=5}{\textbf{Branch and Bound}}: This algorithm solves integer programming problems by recursively partitioning the feasible region into subproblems, solving the relaxed linear program for each subproblem, and pruning branches that cannot yield a better solution than the best integer solution found so far. A depth-first search strategy is often preferred due to the tendency for integer solutions to lie deep in the search tree, enabling early pruning. %
\item \href{https://www.math.purdue.edu/~yipn/421/2024-10-31-IntLinProg.pdf#page=15}{\textbf{Gomory Cuts}}: This algorithm solves integer programming problems by iteratively adding new constraints (Gomory cuts) to the relaxed linear program. These cuts are derived from the fractional parts of the basic variables in the optimal dictionary of the relaxed problem, cutting off non-integer solutions while preserving integer feasible solutions. The process is repeated until an integer optimal solution is found. %
\item \href{https://www.math.purdue.edu/~yipn/421/2024-11-05-NetFlows.pdf#page=22}{\textbf{Primal Simplex Method for Network Flows}}: Iteratively improve a spanning tree solution by identifying an entering variable (arc with most negative reduced cost), a leaving variable (arc in the cycle formed by the entering variable and the spanning tree), updating the spanning tree, and updating primal and dual variables until all dual slacks are non-negative.%
\item \href{https://www.math.purdue.edu/~yipn/421/2024-11-05-NetFlows.pdf#page=14}{\textbf{Finding Primal Flows using a Spanning Tree}}: Start at leaf nodes of the spanning tree and iteratively use flow balance equations to calculate flow values along arcs, working inward until all flows are determined.%
\item \href{https://www.math.purdue.edu/~yipn/421/2024-11-05-NetFlows.pdf#page=19}{\textbf{Finding Dual Variables using a Spanning Tree}}: Starting at the root node, iteratively use the dual flow equation ($y_j - y_i + z_{ij} = c_{ij}$) to compute dual variables ($y_i$) for each node in the spanning tree. Then, compute dual slacks ($z_{ij}$) for arcs not in the spanning tree using $z_{ij} = y_j - y_i + c_{ij}$.%
\item \href{https://www.math.purdue.edu/~yipn/421/2024-11-05-NetFlows.pdf#page=15}{\textbf{Theorem 14.1}}: A square submatrix of the modified incidence matrix ($\tilde{A}$) is a basis if and only if the arcs corresponding to its columns form a spanning tree.%
\item \href{https://www.math.purdue.edu/~yipn/421/2024-11-07-NetSimplex.pdf#page=5}{\textbf{Primal Network Simplex}}: Iteratively choose an entering arc with $z_{ij} < 0$, find a cycle including the entering arc and a leaving arc (the arc with the smallest flow in the reverse direction of the entering arc in the cycle), and update primal and dual flows along the cycle. %
\item \href{https://www.math.purdue.edu/~yipn/421/2024-11-07-NetSimplex.pdf#page=32}{\textbf{Dual Network Simplex}}: Choose a leaving arc with $x_{ij} < 0$, find two subtrees resulting from removing the leaving arc, select an entering arc bridging the subtrees in the opposite direction of the leaving arc with the smallest dual slack, and update primal and dual flows. %
\item \href{https://www.math.purdue.edu/~yipn/421/2024-11-14-NetAppls.pdf#page=3}{\textbf{Transportation Problem with Inequality Constraints}}: Introduce a dummy node to convert inequality constraints into equality constraints. The total supply equals the total demand plus any unmet demand or excess supply flowing to/from the dummy node. %
\item \href{https://www.math.purdue.edu/~yipn/421/2024-11-14-NetAppls.pdf#page=9}{\textbf{Upper Bounded Transshipment Problem}}: Introduce slack variables ($h_{ij}$) to transform upper bound constraints ($0 \le x_{ij} \le u_{ij}$) into equality constraints ($x_{ij} + h_{ij} = u_{ij}$). %
\item \href{https://www.math.purdue.edu/~yipn/421/2024-11-14-NetAppls.pdf#page=17}{\textbf{Shortest Path Problem (Bellman's Equation)}}: Iteratively approximate the shortest distance ($v_i$) from node $i$ to the destination node $r$ using the recursive formula $v_i = \min_{j} \{c_{ij} + v_j\}$, where $c_{ij}$ is the cost of the edge from $i$ to $j$. The algorithm requires at most $m$ iterations, where $m$ is the number of nodes. %
\item \href{https://www.math.purdue.edu/~yipn/421/2024-11-14-NetAppls.pdf#page=21}{\textbf{Shortest Path Problem (Dijkstra's Algorithm)}}: Iteratively build a set $F$ of finished nodes. In each iteration, select the unprocessed node $j$ with the minimum distance estimate $v_j$, add it to $F$, and update the distance estimates of its neighbors. %
\item \href{https://www.math.purdue.edu/~yipn/421/2024-11-19-MaxFlowMinCut.pdf#page=12}{\textbf{Ford{-}Fulkerson Algorithm (Augmented Path)}}: This algorithm iteratively finds augmenting paths in a network and updates the flow along these paths to increase the total flow from the source to the sink until no more augmenting paths can be found.%
\item \href{https://www.math.purdue.edu/~yipn/421/2024-11-19-MaxFlowMinCut.pdf#page=4}{\textbf{Max{-}Flow Min{-}Cut Theorem}}: The maximum flow from the source to the sink in a network is equal to the minimum capacity of any cut in the network.%
\item \href{https://www.math.purdue.edu/~yipn/421/2024-12-03-Duality.pdf#page=3}{\textbf{Graphical Method}}: Solve a linear program with two variables by plotting the constraints to define the feasible region and then finding the vertex that maximizes (or minimizes) the objective function. %
\item \href{https://www.math.purdue.edu/~yipn/421/2024-12-03-Duality.pdf#page=7}{\textbf{Simplex Method in Matrix Form}}: Solve a linear program by iteratively improving a basic feasible solution using matrix operations on the constraint matrix and objective function coefficients. The algorithm involves partitioning variables into basic and non-basic sets, updating the basis matrix (B⁻¹), and checking optimality conditions (reduced costs). %
\item \href{https://www.math.purdue.edu/~yipn/421/2024-12-03-Duality.pdf#page=11}{\textbf{Sensitivity Analysis}}: Determine how changes in the objective function coefficients (c) or constraint values (b) affect the optimal solution. For changes in c, ΔS = (Δc)ᵀx. For changes in b, ΔS = (Δb)ᵀy, where y are the shadow prices. %
\item \href{https://www.math.purdue.edu/~yipn/421/2024-12-03-Duality.pdf#page=22}{\textbf{Network Simplex Method}}: A specialized simplex method for network flow problems. It maintains a spanning tree solution and iteratively improves the flow by finding augmenting paths. %
\end{itemize}

%
\end{document}